\documentclass[11pt,fleqn]{report}

\usepackage{cml}
\usepackage{amssymb}
\usepackage{lmodern}
\usepackage{graphicx}
\usepackage{makecell}
\usepackage{float}

\newcommand{\ModelName}{Math Utilities}
\newcommand{\ModelAuthor}{Shubhodeep Mukherji \\ Gary Turner}
\newcommand{\ModelAffiliation}{Odyssey Space Research}
\newcommand{\ModelDate}{June 2016}

\setcounter{tocdepth}{3}
\setcounter{secnumdepth}{3}

\begin{document}
\pagestyle{empty}
\titlePage
\pagenumbering{roman}
\tableofcontents % Print table of contents
\vfill
{\Large \textbf {Revision History}}
{\large
\begin{center}
\begin{tabular}{|l||l|l|l|}
\hline
 \textbf{Version} & \textbf{Date} & \textbf{Author} & \textbf{Purpose} \\
\hline \hline
1 & June 2016 & Shubhodeep Mukherji & Initial Version \\ \hline
2 & June 2021 & Brent Caughron & Partial revision to cover new capabilities \\
\hline
3 & July 2021 & Gary Turner & Additional capabilities \\ \hline
4 & October 2022 & Gary Turner & Additional capabilities \\ \hline
5 & May 2024 & Gary Turner & Additional capabilities \\ \hline
6 & October 2025 & Hirad Mirhashemi & Additional capabilities \\ \hline
\end{tabular}
\end{center}
}


\chapter{Introduction} % Make a "chapter" heading
\pagestyle{plain}
\pagenumbering{arabic} % Start text with arabic 1
The Math Utilities model provides several mathematical tools intended for use
across multiple models.These include:
\begin{itemize}
\item Generation of transformation matrices to perform the following
coordinate transformations:
\begin{itemize}
\item Between the Inertial and Local-Vertical-Local-Horizontal (LVLH)
frames, given position and velocity vectors in the inertial
frame.
\item Between the Inertial and UVW frames, given position and velocity
vectors in the inertial frame.
\item Between the Inertial frame and a user-specified reference-frame,
with the user-specified frame defined by the location of its
origin in the inertial frame and the direction of its x-axis,
also in the inertial frame.
\item Between the Inertial and Vertical-Normal-Conormal (VNC)
frames, given position and velocity vectors in the inertial
frame.
\item Between the Cartesian planet-fixed and ENU (East-North-Up)
frames, given a position in the planet-fixed frame.
\end{itemize}

\item Generation of left-transformation quaternions to facilitate the
following coordinate transformations:
\begin{itemize}
\item Between the ENU and Cartesian planet-fixed frames.
\end{itemize}

\item Out-of-domain protection in mathematical calculation with limited
domains, including:
\begin{itemize}
\item arc-cosine
\item arc-sine
\item log (base 10 and natural)
\item square-root
\end{itemize}

\item Protected division, protecting against overflow when dividing by
values close to 0.
\item Evaluation of a polynomial function of arbitrary order using a vector
of predefined coefficients.
\item Computation of the roots for a quadratic equation.
\item Cholesky decomposition of a square, symmetric, positive-semi-definite
matrix (i.e. determinant $|A| \geqslant 0, A_{ij} = A_{ji}$)
to generate the "square-root" of a matrix.
\item Numerical differencing to obtain differential values using a
backward-difference operator utilizing a flexible history of up to 5
data points.
\item Computation of the derivative of a unit-vector from the current
values of the vector and the vector derivative.
\item Evaluation of whether two floating-point numbers are numerically equal.
\item Evaluation of whether a floating-point number is positive, zero, or negative.
\item A type-flexible templated method to evaluate whether two values are
within a specified absolute or relative tolerance of each other.
\item Transforming a \textit{6x6} position-velocity covariance matrix from one frame
to another by independently transforming each \textit{3x3} submatrix using a 
\textit{3x3} transformation matrix.
\item Extracting the correlation coefficient matrix from a flexible template-based
size covariance matrix, assuming the covariance is at least populated in the 
lower diagonal.

\item A set of vector and matrix utilities using a flexible template-based
size to support broader use-cases than the conventional 3-vector and
3x3 matrix manipulations provided elsewhere. These include:
\begin{itemize}
\item Populating a one-dimensional array with elements of zero.
\item Normalizing a one-dimensional array to produce an array of unit length.
\item Computing the scalar product of two one-dimensional arrays
of the same size.
\item Computing the vector magnitude of a one-dimensional array.
\item Computing the vector magnitude squared of a one-dimensional array.
\item Applying the square-root to each element of a one-dimensional array independently.
\item Performing element-wise addition between two one-dimensional arrays
of the same size.
\item Performing element-wise subtraction between two one-dimensional arrays
of the same size.
\item Populating a \textit{mxn} matrix with elements of a one-dimensional
array of size \textit{m*n} arranged in column-major format.
\item Populating a \textit{mxn} matrix with elements of a one-dimensional
array of size \textit{m*n} arranged in row-major format.
\item Copying a \textit{nxm} matrix.
\item Copying a \textit{pxq} submatrix from a larger \textit{mxn} matrix into
a smaller \textit{pxq} target matrix, starting at a specified upper-left
corner in the larger matrix.
\item Copying a \textit{pxq} submatrix from a smaller \textit{pxq} matrix into
a larger \textit{mxn} target matrix, starting at a specified upper-left
corner in the larger matrix.
\item Incrementing a \textit{mxn} matrix by another of identical
size.
\item Decrementing a \textit{mxn} matrix by another of identical
size.
\item Scaling of the elements of a \textit{mxn} matrix by a scalar
multiple.
\item Transposing a \textit{mxn} matrix into a \textit{nxm}
matrix.
\item Multiplying a \textit{mxn} matrix with a \textit{nxp} matrix
to produce a \textit{mxp} matrix.
\item Multiplying the transpose of a \textit{mxn} matrix with
a \textit{mxp} matrix to produce a \textit{nxp} matrix.
\item Multiplying a \textit{mxn} matrix with the transpose of a
\textit{pxn} matrix to produce a \textit{mxp} matrix.
\item Multiplying the transpose of a \textit{mxn} matrix with
the transpose of a \textit{pxm} matrix to produce a
\textit{nxp} matrix.
\item Transforming a \textit{mxm} matrix, $A$, by application of a
\textit{nxm} transformation matrix, $T$ to produce a \textit{nxn} 
matrix, B with $B = T A T'$
\item Transforming a \textit{mxm} matrix, $A$, by application of a
\textit{mxn} inverse transformation matrix, $T$ to produce a \textit{nxn} 
matrix, B with $B = T' A T$
\item Inverting a \textit{mxm} matrix. This process uses the
Cholesky decomposition so the matrix must be
symmetric and positive-definite.
\end{itemize}

\item A set of utilities for the conventional 3-vector represented as a one-dimensional 
arrays of size three. These include:
\begin{itemize}
\item Computing the cross product of two 3D vectors.
\item Computing the magnitude of a 3D vector's projection onto the $xy$-plane.
\item Computing the magnitude of a 3D vector's projection onto the $xz$-plane.
\item Computing the magnitude of a 3D vector's projection onto the $yz$-plane.
\end{itemize}
\end{itemize}


\chapter{Requirements}
\begin{enumerate}
\item The model shall provide a convenient interface to support general
mathematical techniques
\begin {itemize}
\item Methods will typically be declared static and thus made
accessible without requiring instantiation of the model's
primary class.
\end {itemize}

\item Specifically, for frame transformations, the model shall provide the
tools to generate:
\begin{enumerate}
\item a transformation matrix to transform between coordinates in the
inertial frame and those in the LVLH frame.
\item a transformation matrix to transform between coordinates in the
inertial frame and those in the UVW frame.
\item a transformation matrix to transform between coordinates in the
Cartesian planet-fixed frame and those in the ENU frame.
\item a transformation matrix to transform between coordinates in the
inertial frame and those in the VNC frame.
\item a transformation quaternion to transform between coordinates
in the Cartesian planet-fixed frame and those in the ENU frame.
\item a transformation matrix to transform coordinates in the
inertial frame to those in a user-specified reference-frame.
\end{enumerate}
\item The model shall provide convenient utilties to protect against
floating-point exceptions when using functions and operations
operating outside their legal domain, including:
\begin{itemize}
\item division, with protection against divide-by-zero
\item inverse sine and cosine, with protection against arguments
with magnitude greater than 1.
\item logarithms, with protection against negative arguments.
\item square-root, with protection against negative arguments.
\end{itemize}
\item The model shall provide an efficient polynomial evaluation function
of arbitrary order.
\item The model shall provide the capability to compute the roots of a 
quadratic equation, even for numerically unstable cases.
\item The model shall provide finite-difference and derivative
computations including:
\begin{itemize}
\item Backward differencing finite difference algorithm of
flexible order.
\item Derivative of unit vectors.
\end{itemize}
\item The model shall provide the ability to determine whether two values
are arbitrarily close to each other, including specifically the
ability to safely evaluate whether two floating-point values are
equivalent to the limit of numerical sensitivity.
\item The model shall provide the ability to determine whether a value is
positive, zero or negative.
\item The model shall provide matrix and vector utilities for
one- and two-dimensional arrays with arbitrary size, as weel as for 
commonly used fixed-size vectors of length 3.
\item The model shall provide the capability to decompose and invert
square matrices with arbitrary size.
\item The model shall support the transformation of a \textit{6x6} 
position-velocity covariance matrix between references frames.
\item The model shall provide functionality to extract a correlation
coefficient matrix from a covariance matrix of arbitrary size
with at least the lower diagonal populated.
\end{enumerate}

%******************************************************************************
%******************************************************************************
%******************************************************************************

\chapter{Model Specification}
\section{Code Structure}
The model is centered around the \texttt{MathUtils} class, which offers a variety of methods
described below. Unless otherwise specified, all methods are static and belong to 
\texttt{MathUtils} class. This class privately inherits from \texttt{MathUtilsPrivate},
allowing it to expose private functionality through \texttt{MathUtils}'s public interfaces,
while keeping the private implementations in a separate header. Although the methods are 
all at the same level in a single class, they are organized into the following groups for
documentation purposes:
\begin{enumerate}
\item Frame transformations provide transformations between the inertial frame
and other frames as specified.
\item ``Protection'' methods provide out-of-domain protection for various
mathematical functions by interfacing with the private implementations
defined in \texttt{MathUtilsPrivate}.
\item Numerical methods provide simple operations on single-value variables.
\item Matrix and Vector methods provide operations on matrices and vectors.
\item Calculus methods provide methods relating to differencing and
differentiation.
\item Geometry methods provide utilities for performing geometric computations
such as ellipsoid intersections.
\end{enumerate}

This arrangement will be followed throughout the rest of the document.

In addition to the \texttt{MathUtils} class, the model includes a set of global
operator overloads for \texttt{std::array<double, N>}. Some of these operators 
can be useful for quick vector math.

%******************************************************************************
%******************************************************************************
\subsection{Frame Transformations}
\subsubsection{Inertial to LVLH}
\begin{code}
static void generate_inertial_to_lvlh(const double position[3],
                                      const double velocity[3],
                                      double T_inrtl_lvlh[3][3])
\end{code}

A utility to generate the transformation matrix for transforming vectors in the
inertial frame to the LVLH frame. This utility requires the position and
velocity of the LVLH frame origin expressed in the inertial frame.

\subsubsection{Inertia to UVW}
\begin{code}
static void generate_inertial_to_uvw(const double position[3],
                                     const double velocity[3],
                                     double T_inrtl_uvw[3][3])
\end{code}

A utility to generate the transformation matrix for transforming vectors in the
inertial frame to the UVW frame. This utility requires the position and
velocity of the UVW frame origin expressed in the inertial frame.

\subsubsection{Inertial to User-specified Reference}
\begin{code}
static void generate_inrtl_to_reference(const double x_axis_inrtl[3],
                                        const double position[3],
                                        double T_inrtl_reference[3][3])
\end{code}

A utility to generate the transformation matrix for transforming vectors in the
inertial frame to a user-specified reference-frame. This utility requires the
direction of the user-specified reference-frame's x-axis and the location of
its origin, both expressed in the inertial frame.

\subsubsection{Inertia to VNC}
\begin{code}
static void generate_inrtl_to_vnc(const double (&position)[3],
                                  const double (&velocity)[3],
                                  double (&T_inrtl_vnc)[3][3])
\end{code}

A utility to generate the transformation matrix for transforming vectors in the
inertial frame to the VNC frame. This utility requires the position and
velocity of the VNC frame origin expressed in the inertial frame.

\subsubsection{Cartesian planet-fixed to East-North-Up (ENU)}
\begin{code}
static void generate_T_pfix_to_enu(const double position_ecef[3],
                                   double T_pfix_to_enu[3][3])

static void generate_Q_enu_to_pfix(double longitude,
                                   double latitude,
                                   jeod::Quaternion & Q_enu_to_pfix)
\end{code}

A utility to generate the transformation matrix for transforming vectors
expressed in the Cartesian planet-fixed frame to the ENU frame, or generate
the left-transformation quaternion for transforming vectors
expressed in the ENU frame to the Cartesian planet-fixed frame.
The transformation-matrix generation requires the vehicle position expressed
in the Cartesian planet-fixed frame. The quaternion generation requires the
vehicle position expressed as a latitude-longitude pair.

Other utilities are available to generate a quaternion from a transformation
matrix (and vice-versa) so either could be generated from either set of
inputs.

Note: The ENU frame becomes ill-defined near the poles due to the
singularity in the local east and north directions. Motion over or near the poles,
such as gyroscope-based navigation crossing the poles can result in large errors
or undefined states in the transformation. It is recommended to not use the ENU
frame near these regions.

%******************************************************************************
%******************************************************************************
\subsection{Protection Methods}

\subsubsection{divide\_protected}
\begin{code}
static float divide_protected(const float dividend,
                              const float divisor,
                              const float failed_val= 0.0f,
                              const boolfailed_flag = false)

static double divide_protected(const double dividend,
                               const double divisor,
                               const double failed_val = 0.0,
                               const bool failed_flag = false)
\end{code}

Two versions of \texttt{divide\_protected(divident, divisor,
failed\_val, failed\_flag)} access an internal templated method
(\texttt{div\_fp}) which centralizes the capability
of handling failure cases for values of type \texttt{float} or
\texttt{double} when dividing the divident by the divisor.

The third and fourth arguments (\texttt{failed\_val, failed\_flag}) are optional
and default to \texttt{0} and \texttt{false} respectively.
These two arguments are only considered in the event that the division results
in a floating-point exception (such as not-a-number or infinity).
In this case, the \texttt{failed\_flag} setting determines the outcome.

\begin{itemize}
\item \textbf{false (default)} the result of the division will be assigned
the specified value \texttt{failed\_val}, an error message will be
broadcast and execution will continue.
\item \textbf{true} execution will be terminated at this point.
\end{itemize}

\subsubsection{acos\_protected}
\begin{code}
static float acos_protected(const float val)

static double acos_protected(const double val)
\end{code}

As with \texttt{divide\_protected}, two versions of
\texttt{acos\_protected(value)} are provided for types
\texttt{float} and \texttt{double}; these access a common
internal templated method (\texttt{acos\_fp}) which
centralizes the capability of handling failure cases when calculating
the arc-cosine of the input value.

\begin{equation*}
y = acos (x) = \left\{
 \begin{array}{ll}
\pi& \quad x \leqslant -1 \\
cos^{-1} x & \quad -1 < x < 1 \\
0& \quad x \geqslant1 \\
 \end{array}
\right.
\end{equation*}

\subsubsection{asin\_protected}
\begin{code}
static float asin_protected(const float val)

static double asin_protected(const double val)
\end{code}

Follows the same design as \texttt{acos\_protected} (above)
by accessing a common internal templated method (\texttt{asin\_fp})
for handling the arc-sine function.
\begin{equation*}
y = asin (x) = \left\{
 \begin{array}{ll}
-\frac{\pi}{2}& \quad x \leqslant -1 \\
sin^{-1} x& \quad -1 < x < 1 \\
\frac{\pi}{2} & \quad x \geqslant1 \\
 \end{array}
\right.
\end{equation*}

\subsubsection{log\_protected}
\begin{code}
static float log_protected(const float val,
                           const float failed_val = 0.0,
                           const bool failed_flag = false)

static double log_protected(const double val,
                            const double failed_val = 0.0,
                            const bool failed_flag = false)
\end{code}

This method follows the same design concept, with two user-interface methods
-- one for \texttt{float} and one for \texttt{double} -- accessing a common
internal templated method (\texttt{log\_fp}) that centralizes the capability
of handling failure cases when calculating the logarithm of the input value.
The signature \texttt{log\_protected(value, failed\_val, failed\_flag)}
forwards to (\texttt{log\_fp(value, failed\_val, failed\_flag, false)} with
the 4th argument indicating that this is to take the natural logarithm
(base e) of the input value.

The second and third arguments (\texttt{failed\_val, failed\_flag}) are
optional and default to \texttt{0.0} and \texttt{false} respectively.
These two arguments are only considered in the event that the input value is
outside the domain of the natural function (i.e. is less than 0.0).
In this case, the \texttt{failed\_flag} setting determines the outcome.
\begin{itemize}
\item \textbf{false (default)} the result will be assigned the specified
value \texttt{failed\_val}, an error message will be broadcast and
execution will continue.
\item \textbf{true} execution will be terminated at this point.
\end{itemize}


\subsubsection{log10\_protected}
\begin{code}
static float log10_protected(const float val,
                             const float failed_val = 0.0,
                             const bool failed_flag = false)

static double log10_protected(const double val,
                              const double failed_val = 0.0,
                              const bool failed_flag = false)
\end{code}

See \textit{log\_protected} above; this method takes the logarithm to base
10 by passing \texttt{true} as the 4th argument to \texttt{log\_fp}.
It otherwise functions identically to \texttt{log\_protected}.

\subsubsection{sqrt\_protected}
\begin{code}
static float sqrt_protected(float val)

static double sqrt_protected(double val)
\end{code}

This method follows the same design concept as the other methods, 
with two user-interface methods -- one for \texttt{float} and one 
for \texttt{double} -- accessing a common internal templated method 
(\texttt{sqrt\_fp}) that centralizes the capability of handling
failure cases when calculating the square-root of the input value. The 
internal method will protect in the event that the input value is outside
the domain of the square-root function (i.e. is less than 0.0). In this case,
a value of 0.0 will be returned, an error message will be broadcast, and 
execution will continue.

%******************************************************************************
%******************************************************************************
\subsection{Numerical Methods}

\subsubsection{is\_near\_equal}
\begin{code}
static bool is_near_equal(const float val1,
                          const float ulp = 0.5f)
                          const float val2,

static bool is_near_equal(const double val1,
                          const double val2,
                          const double ulp = 0.5)

template <size_t N>
static bool is_near_equal(const std::array<double, N>& val1,
                          const std::array<double, N>& val2,
                          const double ulp = 0.5)

template <size_t N>
static bool is_near_equal(const std::array<float, N>& val1,
                          const std::array<float, N>& val2,
                          const float ulp = 0.5)
\end{code}

A version of \texttt{is\_near\_equal(value1, value2, ulp)}
is provided for \texttt{double} and \texttt{float} variables which centralizes
the capability of determining whether two values are effectively equal (i.e.
so close to equal to be indistinguishable at some specified level of
precision) encapsulated within the internal templated method, \texttt{check\_equal}.

A "ulp" (unit in last place) is used
to control the threshold on determining the difference needed to show that the
two compared value are near equal to each other. The \texttt{ulp} defaults to
0.5 (so the values must agree to within 0.5 of the least-significant digit
available in order to be considered ``nearly equal''); this setting can be
adjusted by the user.

\subsubsection{has\_changed\_from}
\begin{code}
template< typename T>
static bool has_changed_from(T val1, T val2)
\end{code}

This is a data-type-templated method that determines whether a variable has changed
from a previously stored value. It compares two values of the same type and returns
\texttt{true} if they differ. Instead of performing a strict inequality comparison,
it checks whether either value is less than the other (i.e. \texttt{(val1<val2) || (val2<val1)}). 
While functionally similar to using \texttt{!=}, this method is semantically different 
as it is intended to detect changes in value over time, and supports more type-specific
comparisons.

\subsubsection{is\_equal}
\begin{code}
template <typename T>
static bool is_equal(T val1, T val2)
\end{code}

This is a data-type-templated method that tests equality between two values of 
the same type. By default, it uses the \texttt{==} operator, but for 
\texttt{float} and \texttt{double} types, it specializes to use 
\texttt{!has\_changed\_from}. This method is intended for use in type-generic code 
to test equality of values where it could be testing floats, doubles or other common types. 
It should only be used to compare floating-point values that may have been previously 
assigned to that variable. For comparisons against independently computed values,
\texttt{is\_within\_abs\_tolerance} is preferred.

\subsubsection{is\_within\_abs\_tolerance}
\begin{code}
template <typename T>
static bool is_within_abs_tolerance(T val1,
                                    T val2,
                                    T tol)
\end{code}

This is a data-type-templated method that determines whether two values are 
within a specified absolute tolerance of each other. It avoids direct 
differencing between the two values to prevent issues with unsigned types
and negative differences. Instead, it compares the values using conditional 
logic to ensure that the larger value does not exceed the smaller value plus
the tolerance.

If the provided tolerance is negative, the method returns \texttt{false}, 
broadcasts an error, and continues execution, as negative tolerances can be 
ambiguous in this context.

\subsubsection{is\_within\_rel\_tolerance}
\begin{code}
template <typename T>
static bool is_within_rel_tolerance(T value,
                                    T expected,
                                    double tol)
\end{code}

This is a data-type-templated method that determines whether a value is within
a specified relative tolerance of an expected value. It converts the relative
tolerance into an absolute tolerance by scaling it as a fraction of the 
expected value, and then delegates the comparison to \texttt{is\_within\_abs\_tolerance}
to avoid direction differencing and ensure compatibility with unsigned types.

If the provided tolerance is negative, the method returns \texttt{false}, 
broadcasts an error, and continues execution, as negative tolerances can be 
ambiguous in this context.

\subsubsection{sign}
\begin{code}
template<typename T>
static int sign (T value)
\end{code}

This is a data-type-templated method that returns the sign of a given value. It 
evaluates to \texttt{1} if the value is positive, \texttt{-1} if the value is negative, 
and \texttt{0} if the value is zero. The implementation compares the value to zero
using conditional logic to work with any type that supports comparison with zero.

Mathematically, this would be:
\[
\text{sign}(x) = 
\begin{cases}
1 & \text{if } x > 0 \\
0 & \text{if } x = 0 \\
 -1 & \text{if } x < 0
\end{cases}
\]

\subsubsection{Polynomial}
\begin{code}
static double polynomial(double x,
                         std::vector<double> & coeffs,
                         const double failed_val= 0.0,
                         const bool failed_flag = false)
\end{code}

This function calculates the result of a polynomial given an input value $x$
and a \texttt{std::vector} of input coefficients \texttt{coeffs} is provided.
\begin{equation*}
f(x) = \sum_{i=0}^{n} c_i x^i
\end{equation*}

The third and fourth arguments (\texttt{failed\_val, failed\_flag}) are
optional and default to \textit{0.0} and \textit{false} respectively.
These two arguments are only considered in the event that the computed result is
either NaN or infinity.
In this case, the \texttt{failed\_flag} setting determines the outcome.
\begin{itemize}
\item \textbf{false (default)} the result will be assigned the specified
value \texttt{failed\_val}, a warning message will be broadcast and
execution will continue.
\item \textbf{true} execution will be terminated at this point.
\end{itemize}

\subsubsection{Quadratic Solver}
\begin{code}
bool QuadraticSolver::solve(bool compute_roots = true)
\end{code}

This capability is encapsulated in the \texttt{QuadraticSolver} class,
and is used to find the roots of the quadratic equation:
\begin{equation*}
ax^2 + bx + c = 0
\end{equation*}

The argument \texttt{compute\_roots} is optional and defaults to \texttt{true}.
It controls whether the method should compute and store the actual roots of the
equation. When set to \texttt{false}, the method will only check and return whether real
roots exist (i.e. whether the discriminant is non-negative, $b^2 - 4ac \geq 0$), 
without computing or modifying the roots. This is intended to be used
when a model only needs to verify the existence of real solutions without
requiring their exact values. 

If \texttt{compute\_roots} is set to \texttt{false} or if the roots don't exist,
the method will reset both roots (\texttt{root1} and \texttt{root2}) to 0.
Therefore, before accessing the values of the roots, it is important to consider
whether the input \texttt{compute\_roots} has been set to \texttt{false} or the
\texttt{solve(...)} method returned \texttt{false}; the latter case can also be
identified by querying the flag \texttt{roots\_exist}.

%******************************************************************************
%******************************************************************************
\subsection {Matrix and Vector Methods}
Matrix and Vector operations are grouped together because they share a
common architecture -- with a ``vector'' being a one-dimensional array, and a
``matrix'' being a two-dimensional array.

\subsubsection{Cholesky Decomposition}
A function that applies a Cholesky decomposition to a matrix (or a sub-matrix
of a matrix), effectively computing the ``square-root'' of the matrix.
\begin{code}
static bool cholesky_decomposition(std::string caller_id,
                                   const double * in_array,
                                   double * out_array,
                                   size_t size_in,
                                   size_t size_out = 0)
\end{code}

\subsubsection{Position-Velocity Covariance Matrix Transformation}
Transforms a 6x6 position-velocity covariance (PV) matrix from one frame to another
by independently applying the provided 3x3 transformation matrix to each 3x3 
submatrix.
\begin{code}
static void transform_pv_matrix(const double (&T_mx)[3][3],
                                const double (&pv_in)[6][6],
                                double (&pv_out)[6][6])
\end{code}

\subsubsection{Correlation Matrix}
Extracts correlation coefficients from a square covariance matrix of arbitrary size
and collects the coefficients into a correlation matrix.
\begin{code}
template <size_t N>
static bool extract_correlation_coefficients(const double cov[N][N],
                                             double corr[N][N])
\end{code}

\subsubsection {Matrix and Vector Algebra}\label{sec:spec:matrix_methods}
The methods implemented for handling basic vector and matrix algebra use
method-templates, with the size of each dimension used as the template
parameters.

\textbf{Zero Vector}\\
Populates a vector of arbitrary size with elements of zero.
\begin{equation*}
\textbf{A} = \textbf{0}
\end{equation*}
\begin{code}
template <size_t N>
static void zero_vector(std::array<double, N>& vec)

template <size_t N>
static void zero_vector(double (&vec)[N])
\end{code}

\textbf{Unit Vector}\\
Normalizes a vector of arbitrary size to produce a unit vector of arbitrary size.
\begin{equation*}
\hat{A} = \textbf{A}
\end{equation*}
\begin{code}
template <size_t N>
static std::array<double, N> unit_vector(const std::array<double, N>& vec)

template <size_t N>
static std::array<double, N> unit_vector( const double (&vec)[N])
\end{code}

\textbf{Scalar Product of Two Vectors}\\
Computes the scalar product of two vectors of identical but otherwise
arbitrary size.
\begin{equation*}
y = \textbf{A} \cdot \textbf{B}
\end{equation*}
\begin{code}
template< size_t N>
static void vector_scalar_product(const double (&left)[N],
                                  const double (&right)[N],
                                  double & result)

template< size_t N>
static double vector_scalar_product(const double (&left)[N],
                                    const double (&right)[N])

template <size_t N>
static double vector_scalar_product(const std::array<double, N>& lhs,
                                    const std::array<double, N>& rhs) 

template <size_t N>
static double vector_scalar_product(const double (&lhs)[N],
                                    const std::array<double, N>& rhs)

template <size_t N>
static double vector_scalar_product(const std::array<double, N>& lhs,
                                    const double (&rhs)[N]) 
\end{code}

\textbf{Cross Product of Two Vectors}\\
Computes the cross product of two vectors of identical size three.
\begin{equation*}
\textbf{Y} = \textbf{A} \times \textbf{B}
\end{equation*}
\begin{code}
static std::array<double, 3> vector_cross_product(
                                        const std::array<double, 3>& lhs,
                                        const std::array<double, 3>& rhs)

static std::array<double, 3> vector_cross_product(
                                        const double (&lhs)[3],
                                        const std::array<double, 3>& rhs)

static std::array<double, 3> vector_cross_product(
                                        const std::array<double, 3>& lhs,
                                        const double (&rhs)[3])

static std::array<double, 3> vector_cross_product(
                                        const double (&lhs)[3],
                                        const double (&rhs)[3])
\end{code}

\textbf{Vector Magnitude}\\
Computes the magnitude of a vector of arbitrary size.
\begin{equation*}
y = ||A||
\end{equation*}
\begin{code}
template <size_t N>
static double vec_mag(const std::array<double, N>& vec)

template <size_t N>
static double vec_mag(const double (&vec)[N])
\end{code}

\textbf{Vector Magnitude Squared}\\
Computes the squared magnitude of a vector of arbitrary size.
\begin{equation*}
y = ||A||^2
\end{equation*}
\begin{code}
template <size_t N>
static double vec_mag_sq(const std::array<double, N>& vec)

template <size_t N>
static double vec_mag_sq(const double (&vec)[N])
\end{code}

\textbf{Vector X-Y Projection Magnitude}\\
Computes the magnitude of a vector of size three projected onto the X-Y plane.
\begin{code}
static double vec_mag_xy( const double (&vec)[3])

static double vec_mag_xy( const std::array<double, 3>& vec)
\end{code}

\textbf{Vector X-Z Projection Magnitude}\\
Computes the magnitude of a vector of size three projected onto the X-Z plane.
\begin{code}
static double vec_mag_xz( const double (&vec)[3])

static double vec_mag_xz( const std::array<double, 3>& vec)
\end{code}

\textbf{Vector Y-Z Projection Magnitude}\\
Computes the magnitude of a vector of size three projected onto the Y-Z plane.
\begin{code}
static double vec_mag_yz( const double (&vec)[3])

static double vec_mag_yz( const std::array<double, 3>& vec)
\end{code}

\textbf{Vector Element-wise Square-Root}\\
Computes the element-wise square-root of a vector of arbitrary size.
\begin{code}
template <size_t N>
static std::array<double, N> vector_elementwise_sqrt(
                                        const std::array<double, N>& vec)

template <size_t N>
static std::array<double, N> vector_elementwise_sqrt(
                                        const std::array<double, N>&& vec)

template <size_t N>
static std::array<double, N> vector_elementwise_sqrt(const double (&vec)[N])

template <size_t N>
static void vector_elementwise_sqrt(const double (&vec)[N],
                                        double (&result)[N])
\end{code}

\textbf{Vector Difference}\\
Computes the element-wise difference of two vectors of identical but otherwise
arbitrary size.
\begin{equation*}
\textbf{Y} = \textbf{A} - \textbf{B}
\end{equation*}
\begin{code}
template <size_t N>
static std::array<double, N> diff(const double (&lhs)[N],
                                  const double (&rhs)[N])
\end{code}

\textbf{Vector Addition}\\
Computes the element-wise addition of two vectors of identical but otherwise
arbitrary size.
\begin{equation*}
\textbf{Y} = \textbf{A} + \textbf{B}
\end{equation*}
\begin{code}
template <size_t N>
static std::array<double, N> sum(const double (&lhs)[N],
                                 const double (&rhs)[N])
\end{code}

\textbf{Column-Major Vector to Matrix}\\
Builds a matrix of arbitrary size from a vector with an identical 
number of elements arranged in column-major format.
\begin{code}
template<size_t N, size_t M>
static void col_maj_vec_to_matrix(const double (&vec)[N*M],
                                  double (&mx)[N][M])
\end{code}

\textbf{Row-Major Vector to Matrix}\\
Builds a matrix of arbitrary size from a vector with an identical 
number of elements arranged in row-major format.
\begin{code}
template<size_t N, size_t M>
static void row_maj_vec_to_matrix(const double (&vec)[N*M],
                                  double (&mx)[N][M])
\end{code}

\textbf{Zero Matrix}\\
Populates a matrix of arbitrary size with elements of zero.
\begin{equation*}
[\textbf{Y}] = [\textbf{0}]
\end{equation*}
\begin{code}
template< size_t dim1, size_t dim2>
static void zero_matrix(double (&out)[dim1][dim2])
\end{code}

\textbf{Matrix Scale}\\
Scales all the elements of a matrix of arbitrary size by a scale
factor.
\begin{equation*}
[\textbf{Y}] = C*[\textbf{A}]
\end{equation*}
\begin{code}
template< size_t dim1, size_t dim2>
static void matrix_scale(const double sf,
                         double (&object)[dim1][dim2])
\end{code}

\textbf{Matrix Increment}\\
Computes the element-wise addition of two matrices of identical 
but otherwise arbitrary size.
\begin{equation*}
[\textbf{Y}] = [\textbf{A}] + [\textbf{B}]
\end{equation*}
\begin{code}
template< size_t dim1, size_t dim2>
static void matrix_incr(const double (&increment)[dim1][dim2],
                        double (&object)[dim1][dim2])
\end{code}

\textbf{Matrix Decrement}\\
Computes the element-wise difference of two matrices of identical 
but otherwise arbitrary size.
\begin{equation*}
[\textbf{Y}] = [\textbf{A}] - [\textbf{B}]
\end{equation*}
\begin{code}
template< size_t dim1, size_t dim2>
static void matrix_decr(const double (&increment)[dim1][dim2],
                        double (&object)[dim1][dim2])
\end{code}

\textbf{Matrix Copy}\\
Copies the elements of a matrix of arbitrary size into one with identical
size.
\begin{equation*}
[\textbf{Y}] = [\textbf{A}]
\end{equation*}
\begin{code}
template< size_t dim1, size_t dim2>
static void matrix_copy(const double (&original)[dim1][dim2],
                        double (&copy)[dim1][dim2])
\end{code}

\textbf{Extract Submatrix}\\
Copies the elements of a submatrix in a larger source matrix of arbitrary size into a
smaller target matrix of arbitrary size. The submatrix is extracted starting at 
the specified top-left corner indices of the source matrix and extends according to the 
dimensions of the target matrix. If the submatrix exceeds the bounds of the
source matrix, only the overlapping portion is copied, and the remaining elements of the 
target matrix are populated as zeros.
\begin{code}
template< size_t row1, size_t col1, size_t row2, size_t col2>
static void matrix_copy_submatrix_out(const double (&original)[row1][col1],
                                      double (&copy)[row2][col2],
                                      size_t top_row_ix,
                                      size_t left_col_ix)
\end{code}

\textbf{Insert Submatrix}\\
Copies the elements of a smaller source matrix of arbitrary size into a
larger target matrix of arbitrary size as a submatrix. The submatrix is inserted starting at 
the specified top-left corner indices of the target matrix and extends according to the 
dimensions of the source matrix. If the inserted submatrix extends beyond the bounds of the
target matrix, only the overlapping portion is copied, and the remaining elements of the 
source matrix are ignored.
\begin{code}
template< size_t row1, size_t col1, size_t row2, size_t col2>
static void matrix_copy_submatrix_in(const double (&original)[row1][col1],
                                     double (&copy)[row2][col2],
                                     size_t top_row_ix,
                                     size_t left_col_ix)
\end{code}

\textbf{Matrix Transpose}\\
Transposes the matrix of arbitrary size.
\begin{equation*}
[\textbf{Y}] = [\textbf{A}]^T
\end{equation*}
\begin{code}
template< size_t dim1, size_t dim2>
static void matrix_trans(const double (&in)[dim1][dim2],
                         double (&result)[dim2][dim1])
\end{code}

\textbf{Basic Matrix Product:}
Performs matrix multiplication on two matrices with arbitrary but 
compatible sizes. The second dimension in the left matrix must equal
the first dimension in the right matrix.
\begin{equation*}
[\textbf{Y}] = [\textbf{A}] [\textbf{B}]
\end{equation*}
\begin{code}
template< size_t dim1, size_t dim2, size_t dim3>
static void matrix_mult(const double (&left)[dim1][dim2],
                        const double (&right)[dim2][dim3],
                        double (&result)[dim1][dim3])
\end{code}

\textbf{Left-Transpose Product:}
Performs matrix multiplication between the transpose of the left matrix 
and the right matrix. The matrices may be of arbitrary size, provided that
the first dimensions are equal size.
\begin{equation*}
[\textbf{Y}] = [\textbf{A}]^T [\textbf{B}]
\end{equation*}
\begin{code}
template< size_t dim1, size_t dim2, size_t dim3>
static void matrix_mult_left_trans(const double (&left)[dim1][dim2],
                                   const double (&right)[dim1][dim3],
                                   double (&result)[dim2][dim3])
\end{code}

\textbf{Right-Transpose Product:}
Performs matrix multiplication between the left matrix and the transpose of 
the right matrix. The matrices may be of arbitrary size, provided that
the second dimensions are equal size.
\begin{equation*}
[\textbf{Y}] = [\textbf{A}] [\textbf{B}]^T
\end{equation*}
\begin{code}
template< size_t dim1, size_t dim2, size_t dim3>
static void matrix_mult_right_trans(const double (&left)[dim1][dim2],
                                    const double (&right)[dim3][dim2],
                                    double (&result)[dim1][dim3])
\end{code}

\textbf{Double-Transpose Product:}
Performs matrix multiplication between the transpose of the left matrix 
and the transpose of the right matrix. The matrices may be of arbitrary
size, provided that the first dimension of the left matrix equals the
second dimension of the right matrix.
\begin{equation*}
[\textbf{Y}] = [\textbf{A}]^T [\textbf{B}]^T
\end{equation*}
\begin{code}
template< size_t dim1, size_t dim2, size_t dim3>
static void matrix_mult_trans_trans(const double (&left)[dim1][dim2],
                                    const double (&right)[dim3][dim1],
                                    double (&result)[dim2][dim3])
\end{code}

\textbf{Matrix Transformation:}
Applies a square transformation matrix to a square matrix defined in the
original frame, producing the transformed matrix in the new frame. The 
transformation matrix must represent the change of basis from the original 
frame to the new frame (e.g. $\textbf{T} = T_{B \rightarrow C}$ when 
$\textbf{A}$ is in frame B).
\begin{equation*}
[\textbf{Y}] = [\textbf{T}] [\textbf{A}] [\textbf{T}]^T
\end{equation*}
\begin{code}
template< size_t dim1, size_t dim2>
static void matrix_transformation(const double (&transform)[dim1][dim2],
                                  const double (&target)[dim2][dim2],
                                  double (&result)[dim1][dim1])
\end{code}

\textbf{Matrix Inverse Transformation:}
Applies the inverse of a square transformation matrix to a square matrix defined in the
original frame, producing the transformed matrix in the new frame. The 
transformation matrix must represent the change of basis from the new 
frame to the original frame (e.g. $\textbf{T} = T_{C \rightarrow B}$ when 
$\textbf{A}$ is in frame B).
\begin{equation*}
[\textbf{Y}] = [\textbf{T}]^T [\textbf{A}] [\textbf{T}]
\end{equation*}
\begin{code}
template< size_t dim1, size_t dim2>
static void matrix_inverse_transformation(const double (&transform)[dim1][dim2],
                                          const double (&target)[dim1][dim1],
                                          double (&result)[dim2][dim2])
\end{code}

\subsubsection{Matrix Inversion (via Cholesky Decomposition)}
The only matrix inversion currently provided is limited to inverting a
symmetric, positive-definite matrix via the Cholesky decomposition.
\begin{equation*}
[\textbf{Y}] = [\textbf{A}]^{-1}
\end{equation*}
This method has an extra optional flag indicating the desired response in
case the matrix inversion fails (e.g. if the
determinant is non-positive). By default (argument omitted) or if passed in as
\textit{true}, failure will terminate the simulation.
If passed in as \textit{false}, the method will:
\begin{itemize}
 \item populate the target matrix (\textbf{Y}) with a zero-matrix
 \item post an error message
 \item return a \textit{false} outcome to the calling method, allowing the
 calling method to attempt a recovery.
\end{itemize}

\begin{code}
template <size_t dimension>
static bool matrix_inv_using_cholesky(
            const double (in_array)[dimension][dimension],
            double (&result)[dimension][dimension],
            bool fail_on_error = true)
\end{code}

%******************************************************************************
%******************************************************************************
\subsection{Calculus Methods}


\subsubsection{Backward Difference}
A function that provides the backward difference when provided with a set of
(equally-spaced) historical values.
For example, when estimating the derivative by use of a second-order
backward-difference method, one would use:
\begin{equation*}
f'(x_n) = \frac{3 f(x_n) - 4f(x_{n-1}) + f(x_{n-2})} {2h}
\end{equation*}
where $h$ is the interval between adjacent values $x_i$.$h = x_j - x_{j-1}$

This function evaluates the component
\begin{equation*}
\frac{3 f(x_n) - 4f(x_{n-1}) + f(x_{n-2})} {2}
\end{equation*}
if provided with 3 values for the variable $x$

The function includes coefficients to support up to a 4th-order difference
(requiring 5 historical values).

\begin{code}
static double compute_backward_difference( const std::list<double> & history)
\end{code}

\subsubsection{Unit Vector Derivative}
A function that computes the derivative of a unit vector given the original
vector (not normalized) and its derivative.

\begin{code}
static void compute_unit_vector_derivative(const double vector[3],
                                           const double vector_derivative[3],
                                           double unit_vector_derivative[3])
\end{code}

%******************************************************************************
%******************************************************************************
\subsection{Geometry Methods}
\subsubsection{Ellipsoid Intersection}
The Ellipsoid-intersection capability is represented in its own
\texttt{EllipsoidIntersection} class. Given an ellipsoid defined by
\begin{equation*}
\left( \frac{x}{a} \right)^2 +
\left( \frac{y}{b} \right)^2 +
\left( \frac{z}{c} \right)^2= 1
\end{equation*}

and a line between two defined points, the
\texttt{EllipsoidIntersection::update()} method will compute the
locations at which the line intersects with the ellipsoid.

If it is only
necessary to know \textit{whether} the line intersects the ellipsoid,
passing a \texttt{false} argument (
\texttt{EllipsoidIntersection::update(false)}) will be more efficient and
the roots will not be computed.

The \texttt{EllipsoidIntersection::update()} method returns a boolean value
indicating whether the intersection exists.

\begin{code}
bool EllipsoidIntersection::update(bool compute_roots = true)
\end{code}

If \texttt{compute\_roots} is set to \texttt{false} or if the lines don't intersect
the ellipsoid, the method will maintain the last found intersection points. Therefore, 
it is important to reference whether \texttt{compute\_roots} or the return flag, 
\texttt{intersection} indicating whether the intersection actually occurs, are enabled,
before accessing the intersection points.

%******************************************************************************
%******************************************************************************
\subsection{\texttt{std::array<double, N>} Operator Overloads}
All operators are overloaded exclusively for double types, with a flexible,
templated size N.

\textbf{Insertion}

Outputs all elements of the \texttt{std::array} to an output stream in bracketed, 
comma-separated form.

Example: \texttt{std::cout << arr;} \qquad (prints \texttt{[1.0, 2.0, 3.0]})
\begin{code}
template<size_t N>
std::ostream& operator<<(std::ostream& out, const std::array<double, N>& rhs)
\end{code}

\textbf{Copying/Casting}

Copies/casts from C-style arrays to \texttt{std::array}.

Example: \texttt{copy = std\_copy(arr);}
\begin{code}
template<size_t N>
std::array<double, N> std_copy(const double (&c_arr)[N])
\end{code}

Note: Two additional casting options are implemented using 
\texttt{reinterpret\_cast} to treat C-style arrays as \texttt{std::array}. 
However, the use of \texttt{reinterpret\_cast} requires a coding standards 
waiver with justification, which has not been pursued since these casts
are not currently needed anywhere. Their implementations are documented
below for possible future reference:

\begin{code}
/* Cast the C-style array as an immutable std::array
result = as_std(c_arr) + other_c_arr;

This is guaranteed to be a well-formed cast according to the 2011 ISO C++
standard:
  "array<T, N> stores N elements of type T, so that size() == N is an
  invariant. The elements of an array are stored contiguously, meaning that
  if a is an array<T, N> then it obeys the identity &a[n] == &a[0] + n for
  all 0 <= n < N."
This remains unchanged in the 2014, 2017, and 2020 revisions.

While standard doesn't prevent a compiler from adding additional members
after the data, accessing past array bounds is undefined behavior anyways.

Furthermore, this cast and all of the overloads check the type of the
arguments. If you attempt to pass a double[N] or std::array<double, N> of
mismatched length, the code won't compile. These functions won't even compile
if you try passing a double*. The array lengths are guaranteed to match. */

template <size_t N>
const std::array<double, N>& as_std(const double (&c_arr)[N]) {
 static_assert(sizeof(c_arr) == sizeof(std::array<double, N>),
     "std::array<double, N> isn't layout-compatible with double[N]. Are you "
     "using a sane compiler?");
 return reinterpret_cast<const std::array<double, N>&>(c_arr);
}

/* Cast the C-style array as a mutable std::array
* as_std_mut(c_arr) += other_c_arr;
*
* This call signature has explicitly been appended with _mut to ensure that the
* programmer is deliberately mutating the array.
*
* See discussion of well-formedness of cast in the description of as_std. */

template <size_t N>
std::array<double, N>& as_std_mut(double (&c_arr)[N]) {
 static_assert(sizeof(c_arr) == sizeof(std::array<double, N>),
     "std::array<double, N> isn't layout-compatible with double[N]. Are you "
     "using a sane compiler?");
 return reinterpret_cast<std::array<double, N>&>(c_arr);
}
\end{code}

\textbf{Addition Assignment}

Assigns the result of element-wise addition to the left-hand operand.
The left-hand operand must be an \texttt{std::array}, and the right-hand 
operand may be either an \texttt{std::array} or a C-style array 
of the same size.

Example: \texttt{lhs += rhs;}
\begin{code}
template<size_t N>
std::array<double, N>& operator+=(std::array<double, N>& lhs, 
                                  const std::array<double, N>& rhs)

template<size_t N>
std::array<double, N>& operator+=(std::array<double, N>& lhs, 
                                  const double (&rhs)[N])
\end{code}

\textbf{Addition}

Returns a new array representing element-wise addition.  
Supports both lvalues and rvalues, with either two \texttt{std::array} objects or
a combination of a \texttt{std::array} and a C-style array, all of the same size.

Example: \texttt{result = lhs + rhs;}
\begin{code}
template<size_t N>
std::array<double, N> operator+(const std::array<double, N>& lhs, 
                                const std::array<double, N>& rhs)

template<size_t N>
std::array<double, N> operator+(std::array<double, N>&& lhs, 
                                const std::array<double, N>& rhs)

template<size_t N>
std::array<double, N> operator+(const std::array<double, N>& lhs, 
                                std::array<double, N>&& rhs)

template<size_t N>
std::array<double, N> operator+(std::array<double, N>&& lhs, 
                                std::array<double, N>&& rhs)

template<size_t N>
std::array<double, N> operator+(const std::array<double, N>& lhs, 
                                const double (&rhs)[N])

template<size_t N>
std::array<double, N> operator+(std::array<double, N>&& lhs, 
                                const double (&rhs)[N])

template<size_t N>
std::array<double, N> operator+(const double (&lhs)[N], 
                                const std::array<double, N>& rhs)

template<size_t N>
std::array<double, N> operator+(const double (&lhs)[N], 
                                std::array<double, N>&& rhs)
\end{code}

\textbf{Unary Minus}  

Returns a new array representing element-wise negation of the \texttt{std::array}.  
Supports both lvalues and rvalues, with \texttt{std::array}.

Example: \texttt{result = -arr;}  
\begin{code}
template<size_t N>
std::array<double, N> operator-(const std::array<double, N>& rhs)

template<size_t N>
std::array<double, N> operator-(std::array<double, N>&& rhs)
\end{code}

\textbf{Subtraction Assignment}  

Assigns the result of element-wise addition to the left-hand operand, 
subtracting the right-hand side from the left-hand side.
The left-hand operand must be an \texttt{std::array}, and the right-hand 
operand may be either an \texttt{std::array} or a C-style array 
of the same size.

Example: \texttt{lhs -= rhs;} 
\begin{code}
template<size_t N>
std::array<double, N>& operator-=(std::array<double, N>& lhs, 
                                  const std::array<double, N>& rhs)

template<size_t N>
std::array<double, N>& operator-=(std::array<double, N>& lhs, 
                                  const double (&rhs)[N])
\end{code}

\textbf{Subtraction}  

Returns a new array representing element-wise subtraction
of the right-hand side from the left-hand side.
Supports both lvalues and rvalues, with either two \texttt{std::array} objects or
a combination of a \texttt{std::array} and a C-style array, all of the same size.

Example: \texttt{result = lhs - rhs;}  
\begin{code}
template<size_t N>
std::array<double, N> operator-(const std::array<double, N>& lhs, 
                                const std::array<double, N>& rhs)

template<size_t N>
std::array<double, N> operator-(std::array<double, N>&& lhs, 
                                const std::array<double, N>& rhs)

template<size_t N>
std::array<double, N> operator-(const std::array<double, N>& lhs, 
                                std::array<double, N>&& rhs)

template<size_t N>
std::array<double, N> operator-(std::array<double, N>&& lhs, 
                                std::array<double, N>&& rhs)

template<size_t N>
std::array<double, N> operator-(const std::array<double, N>& lhs, 
                                const double (&rhs)[N])

template<size_t N>
std::array<double, N> operator-(std::array<double, N>&& lhs, 
                                const double (&rhs)[N])

template<size_t N>
std::array<double, N> operator-(const double (&lhs)[N], 
                                const std::array<double, N>& rhs)

template<size_t N>
std::array<double, N> operator-(const double (&lhs)[N], 
                                std::array<double, N>&& rhs)
\end{code}

\textbf{Scalar Multiplication Assignment}  

Assigns the result of element-wise scalar multiplication to the left-hand operand.
The left-hand operand must be an \texttt{std::array}, and the right-hand 
operand must be a \texttt{double}.

Example: \texttt{lhs *= rhs;} 
\begin{code}
template <size_t N>
std::array<double, N>& operator*=(std::array<double, N>& lhs, 
                                  const double& rhs)
\end{code}

\textbf{Scalar Multiplication}  

Returns a new array representing element-wise scalar multiplication.  
Supports both lvalues and rvalues of \texttt{std::array} with a \texttt{double}.

Example: \texttt{result = arr * 2.0;} or \texttt{result = 2.0 * arr;}  
\begin{code}
template<size_t N>
std::array<double, N> operator*(const std::array<double, N>& lhs, 
                                double rhs)
                                
template<size_t N>
std::array<double, N> operator*(std::array<double, N>&& lhs, 
                                double rhs)

template<size_t N>
std::array<double, N> operator*(double lhs, 
                                const std::array<double, N>& rhs)

template<size_t N>
std::array<double, N> operator*(double lhs, 
                                std::array<double, N>&& rhs)
\end{code}

\textbf{Multiplication Assignment}  

Assigns the result of element-wise multiplication to the left-hand operand.
The left-hand operand must be an \texttt{std::array}, and the right-hand 
operand may be either an \texttt{std::array} or a C-style array 
of the same size.

Example: \texttt{lhs *= rhs;}  
\begin{code}
template<size_t N>
std::array<double, N>& operator*=(std::array<double, N>& lhs, 
                                  const std::array<double, N>& rhs)

template<size_t N>
std::array<double, N>& operator*=(std::array<double, N>& lhs, 
                                  const double (&rhs)[N])
\end{code}

\textbf{Multiplication}  

Returns a new array representing element-wise multiplication.  
Supports both lvalues and rvalues, with either two \texttt{std::array} objects or
a combination of a \texttt{std::array} and a C-style array, all of the same size.

Example: \texttt{result = lhs * rhs;}  
\begin{code}
template<size_t N>
std::array<double, N> operator*(const std::array<double, N>& lhs, 
                                const std::array<double, N>& rhs)

template<size_t N>
std::array<double, N> operator*(std::array<double, N>&& lhs, 
                                const std::array<double, N>& rhs)

template<size_t N>
std::array<double, N> operator*(const std::array<double, N>& lhs, 
                                std::array<double, N>&& rhs)

template<size_t N>
std::array<double, N> operator*(std::array<double, N>&& lhs, 
                                std::array<double, N>&& rhs)

template<size_t N>
std::array<double, N> operator*(const std::array<double, N>& lhs, 
                                const double (&rhs)[N])

template<size_t N>
std::array<double, N> operator*(std::array<double, N>&& lhs, 
                                const double (&rhs)[N])

template<size_t N>
std::array<double, N> operator*(const double (&lhs)[N], 
                                const std::array<double, N>& rhs)

template<size_t N>
std::array<double, N> operator*(const double (&lhs)[N], 
                                std::array<double, N>&& rhs)
\end{code}

\textbf{Scalar Division Assignment}  

Assigns the result of element-wise scalar division to the left-hand operand.
The left-hand operand must be an \texttt{std::array}, and the right-hand 
operand must be a \texttt{double}. 
Division-by-zero is protected using \texttt{MathUtils::divide\_protected}.  

Example: \texttt{lhs /= 2.0;}  
\begin{code}
template<size_t N>
std::array<double, N>& operator/=(std::array<double, N>& lhs, 
                                  const double& rhs)
\end{code}

\textbf{Scalar Division}  

Returns a new array representing element-wise scalar division.  
Supports both lvalues and rvalues of \texttt{std::array} with a \texttt{double}.
Division-by-zero is protected using \texttt{MathUtils::divide\_protected}.  

Example: \texttt{result = arr / 2.0;} or \texttt{result = 2.0 / arr;}  
\begin{code}
template<size_t N>
std::array<double, N> operator/(const std::array<double, N>& lhs, 
                                const double& rhs)

template<size_t N>
std::array<double, N> operator/(std::array<double, N>&& lhs, 
                                const double& rhs)

template<size_t N>
std::array<double, N> operator/(const double& lhs, 
                                const std::array<double, N>& rhs)

template<size_t N>
std::array<double, N> operator/(const double& lhs, 
                                std::array<double, N>&& rhs)
\end{code}

\textbf{Division Assignment}  

Assigns the result of element-wise division to the left-hand operand.
The left-hand operand must be an \texttt{std::array}, and the right-hand 
operand may be either an \texttt{std::array} or a C-style array 
of the same size.

Example: \texttt{lhs /= rhs;}  
\begin{code}
template<size_t N>
std::array<double, N>& operator/=(std::array<double, N>& lhs, 
                                  const std::array<double, N>& rhs)

template<size_t N>
std::array<double, N>& operator/=(std::array<double, N>& lhs, 
                                  const double (&rhs)[N])
\end{code}

\textbf{Division}  

Returns a new array representing element-wise division.  
Supports both lvalues and rvalues, with either two \texttt{std::array} objects or
a combination of a \texttt{std::array} and a C-style array, all of the same size.

Example: \texttt{result = lhs / rhs;}  
\begin{code}
template<size_t N>
std::array<double, N> operator/(const std::array<double, N>& lhs, 
                                const std::array<double, N>& rhs)

template<size_t N>
std::array<double, N> operator/(std::array<double, N>&& lhs, 
                                const std::array<double, N>& rhs)

template<size_t N>
std::array<double, N> operator/(const std::array<double, N>& lhs, 
                                std::array<double, N>&& rhs)

template<size_t N>
std::array<double, N> operator/(std::array<double, N>&& lhs, 
                                std::array<double, N>&& rhs)

template<size_t N>
std::array<double, N> operator/(const std::array<double, N>& lhs, 
                                const double (&rhs)[N])

template<size_t N>
std::array<double, N> operator/(std::array<double, N>&& lhs, 
                                const double (&rhs)[N])

template<size_t N>
std::array<double, N> operator/(const double (&lhs)[N], 
                                const std::array<double, N>& rhs)

template<size_t N>
std::array<double, N> operator/(const double (&lhs)[N], 
                                std::array<double, N>&& rhs)
\end{code}

%******************************************************************************
%******************************************************************************
%******************************************************************************
\section{Mathematical Formulation}

\subsection{Frame Transformations}

\subsubsection{Inertial to LVLH}
The LVLH frame is a rotating coordinate system where the y-axis is normal to
the
orbit plane and opposite of the angular momentum vector, the z-axis
is along the position vector pointing toward the center of the Earth and
the x-axis
completes a right-handed system. In order for these axes to be well-defined,
it
is assumed that the position vector is nonzero and that the velocity vector is
neither radial nor zero. Assuming $\hat{\textbf{r}}$ and $\hat{\textbf{v}}$
are the unit position and velocity vectors in the inertial frame,
respectively, let
\begin{multline}
 \hat{\textbf{\i}} = \hat{\textbf{\j}} \times \hat{\textbf{k}}, \\
 \hat{\textbf{\j}} = \frac{\hat{\textbf{v}} \times \hat{\textbf{r}}}
{||\hat{\textbf{v}} \times \hat{\textbf{r}} ||} \\
 \hat{\textbf{k}} = -\hat{\textbf{r}} \hfill
 \label{eq:lvlh_main}
 \end{multline}
be representations of the LVLH coordinate axes expressed in the inertial
frame.
If the assumptions on the position and velocity vectors are violated, then
(\ref{eq:lvlh_main})
cannot be used to align the LVLH frame and alternative methods are needed.
Implementation of an alternative method allows the simulation to continue
without risking segmentation faults and generates an error message describing
the error source and the chosen method.

In the event of a zero position vector, the angular momentum and position
vectors are both not defined, and the LVLH frame is aligned with the inertial
frame. Hence, the transformation matrix will be the identity matrix.

Now, if the position vector is defined but the angular momentum vector is not,
then the velocity vector in the definition of the LVLH y-axis is replaced by
the
inertial x-axis. This substitution results in the inertial and LVLH x-axes
being
aligned in the same sense, i.e. their dot product will be nonnegative.
Assuming
$\hat{\textbf{e}}_x$ is the unit vector pointing in the direction of the
inertial x-axis, the equations
\begin{multline}
 \hat{\textbf{\i}} = \hat{\textbf{\j}} \times \hat{\textbf{k}} \\
 \hat{\textbf{\j}} = \frac{\hat{\textbf{e}}_x \times \hat{\textbf{r}}}
{||\hat{\textbf{e}}_x \times \hat{\textbf{r}} ||} \\
 \hat{\textbf{k}} = -\hat{\textbf{r}} \hfill
 \label{eq:lvlh_nox}
\end{multline}
can be used to define the LVLH coordinate axes when the position vector is
defined but the angular momentum vector is not. Note that the expression
for $\hat{\textbf{\i}}$ can also be expressed as:
\begin{equation}
 \hfill \hat{\textbf{\i}} =
\frac{\hat{\textbf{e}}_x -
\hat{\textbf{r}}
\left( \hat{\textbf{e}}_x \cdot \hat{\textbf{r}} \right)}
 {||\hat{\textbf{e}}_x \times \hat{\textbf{r}}||} \hfill
\end{equation}

by substituting for $\hat{\textbf{\j}}$ and $\hat{\textbf{k}}$.

If in addition, the position vector lies along the inertial x-axis, then
(\ref{eq:lvlh_nox}) cannot be used because
$\hat{\textbf{r}}= \pm \hat{\textbf{e}}_x$.
In this case, the LVLH y-axis is aligned in the same sense as the inertial
y-axis.
Assuming $\hat{\textbf{e}}_y$ and $\hat{\textbf{e}}_z$ are unit vectors
pointing
in the direction of the inertial y-axis and the inertial z-axis, respectively,
the equations
\begin{multline}
 \hat{\textbf{\i}}= \pm \hat{\textbf{e}}_z \\
 \hat{\textbf{\j}}=\hat{\textbf{e}}_y \\
 \hat{\textbf{k}}=-\hat{\textbf{r}}=\mp \hat{\textbf{e}}_x\hfill
 \label{eq:lvlh_x}
\end{multline}
can be used to define the LVLH coordinate axes.

Finally, the transformation matrix
\begin{equation}
 \hfill R = \begin{bmatrix} \hat{\textbf{\i}}^T \\
\hat{\textbf{\j}}^T \\
\hat{\textbf{k}}^T
\end{bmatrix} \hfill
 \label{eq:lvlh_mat}
\end{equation}
can be used to transform vectors in the inertial frame to the LVLH frame
with the \\ \texttt{generate\_inertial\_to\_lvlh} method.


\subsubsection{Inertial to UVW}
In the UVW frame, the w-axis is in the direction of the angular momentum
vector,
the u-axis is in the direction of the position vector and the v-axis
completes a right-handed system. As with the LVLH formulation, it is assumed
here that the position vector is nonzero and that the velocity vector is
neither
radial nor zero in order for the axes to be well-defined.
Assuming $\hat{\textbf{r}}$ and $\hat{\textbf{v}}$ are the unit position
and velocity vectors in the inertial frame, respectively, let:
\begin{multline}
 \hat{\textbf{\i}} = \hat{\textbf{r}} \\
 \hat{\textbf{\j}} = \hat{\textbf{k}} \times \hat{\textbf{\i}} \\
 \hat{\textbf{k}} = \frac{\hat{\textbf{r}} \times \hat{\textbf{v}}}
 {|| \hat{\textbf{r}} \times \hat{\textbf{v}}||}
 \hfill
 \label{eq:uvw_main}
\end{multline}
be representations of the UVW coordinate axes expressed in the inertial frame.
If the assumptions on the position and velocity vectors are violated,
then (\ref{eq:uvw_main}) cannot be used to align the UVW frame and alternative
methods are needed. As with the inertial to LVLH transformation, implementing
an alternative method allows the simulation to continue and generates an
error message indicating the source and chosen method.

In the event of a zero position vector, the angular momentum and position
vectors are both not defined, and the UVW frame is aligned with the inertial
frame.Hence, the transformation matrix will be the identity matrix.

Now, if position vector is defined but the angular momentum vector is not,
then the original definition of the UVW w-axis is not defined and is
instead aligned in the same sense as the inertial z-axis while remaining
orthogonal to the UVW u-axis.In this scenario, the equations
\begin{multline}
 \hat{\textbf{\i}} = \hat{\textbf{r}} \\
 \hat{\textbf{\j}} = \frac{\hat{\textbf{e}}_z \times \hat{\textbf{r}}}
{||\hat{\textbf{e}}_z \times \hat{\textbf{r}}||} \\
 \hat{\textbf{k}} = \hat{\textbf{\i}} \times \hat{\textbf{\j}} \hfill
 \label{eq:uvw_noz}
\end{multline}
can be used to define the UVW coordinate axes.
Note that the expression for $\hat{\textbf{k}}$ can also be expressed as
\begin{equation}
 \hfill \hat{\textbf{k}} =
 \frac{ \hat{\textbf{e}}_z -
\hat{\textbf{r}}
\left( \hat{\textbf{e}}_z \cdot \hat{\textbf{r}} \right)}
{||\hat{\textbf{e}}_z \times \hat{\textbf{r}}||} \hfill
\end{equation}

by substituting for $\hat{\textbf{\i}}$ and $\hat{\textbf{\j}}$.

If in addition, the position vector lies along the inertial z-axis, then
(\ref{eq:uvw_noz}) cannot be used because
$\hat{\textbf{r}}= \pm \hat{\textbf{e}}_z$.
The UVW w-axis will instead be aligned in the same sense as the inertial
x-axis, which involves replacing $\hat{\textbf{e}}_z$ with
$\hat{\textbf{e}}_x$
in the definition of the UVW v-axis in (\ref{eq:uvw_noz}). After
simplification, this results in the equations
\begin{multline}
 \hat{\textbf{\i}} = \pm \hat{\textbf{e}}_z \\
 \hat{\textbf{\j}} = \mp \hat{\textbf{e}}_y \\
 \hat{\textbf{k}} = \hat{\textbf{e}}_x \hfill
 \label{eq:uvw_z}
\end{multline}
for the definition of the UVW coordinate axes.

With these definitions, the transformation matrix
\begin{equation}
 \hfill R = \begin{bmatrix} \hat{\textbf{\i}}^T \\
\hat{\textbf{\j}}^T \\
\hat{\textbf{k}}^T
\end{bmatrix} \hfill
 \label{eq:uvw_mat}
\end{equation}
can be used to transform vectors in the inertial frame to the UVW frame
with the \\ \texttt{generate\_inertial\_to\_uvw} method.


\subsubsection{Inertial to Reference}
In this transformation it is assumed that the x-axis of the user-specified
reference-frame is given. Hence, $\hat{\textbf{\i}}$ is already specified and
is the unit vector in the direction of the given x-axis vector. In order for
these axes to be well-defined, it is assumed that the x-axis and position
vectors are both nonzero and not aligned with each other. Assuming
$\hat{\textbf{r}}$ is the unit position vector, let:

\begin{multline}
 \hat{\textbf{\j}} = \frac{\hat{\textbf{\i}} \times \hat{\textbf{r}}}
{||\hat{\textbf{\i}} \times \hat{\textbf{r}}||} \\
 \hat{\textbf{k}} = \hat{\textbf{\i}} \times \hat{\textbf{\j}} \hfill
 \label{eq:ref_main}
\end{multline}
complete the definition of the reference-frame coordinate axes.
If the assumptions on the x-axis and position vectors are violated,
then (\ref{eq:ref_main}) cannot be used to align the frame and
alternative methods are needed. As with the previous
transformations, implementing an alternative method allows the simulation to
continue and generates an error message indicating the source and chosen method.

If the position and reference vectors are aligned with each other, then
their cross product returns zero and the expression for $\hat{\textbf{\j}}$
is invalid. If these vectors are \textit{not} aligned with the inertial z-axis,
then we can substitute the inertial z-axis in place of the position vector,
and:
\begin{multline}
 \hat{\textbf{\j}} = \frac{\hat{\textbf{\i}} \times \hat{\textbf{e}}_z}
{||\hat{\textbf{\i}} \times \hat{\textbf{e}}_z||} \\
 \hat{\textbf{k}} = \hat{\textbf{\i}} \times \hat{\textbf{\j}} \hfill
 \label{eq:ref_noz}
\end{multline}
complete the definition of the reference-frame coordinate axes.

In this case, the new reference-frame z-axis ($\hat{\textbf{k}}$) is aligned in
generally the opposite sense as the inertial z-axis (i.e. it will have a
negative z-axis component when expressed in the inertial frame).
Note that the expression for $\hat{\textbf{k}}$ can also be expressed as
\begin{equation*}
 \hat{\textbf{k}} = \frac{ \hat{\textbf{r}}
 \left( \hat{\textbf{r}} \cdot
\hat{\textbf{e}}_z
 \right) -
 \hat{\textbf{e}}_z
}
{||\hat{\textbf{r}} \times \hat{\textbf{e}}_z||}
\hfill
\end{equation*}
by substituting for $\hat{\textbf{\i}}$ and $\hat{\textbf{\j}}$.

If the position and reference vectors are aligned with each other \textit{and}
the inertial z-axis, then the reference-frame
y-axis is instead aligned with the inertial y-axis. In this case, the
reference-frame axes become:
\begin{multline}
 \hat{\textbf{\i}} = \pm \hat{\textbf{e}}_z \\
 \hat{\textbf{\j}} = \hat{\textbf{e}}_y \\
 \hat{\textbf{k}}= \mp \hat{\textbf{e}}_x \hfill
 \label{eq:ref_z}
\end{multline}

Finally, in the event that either vector is zero, the reference-frame is
undefined in a non-recoverable way; it is aligned with the inertial frame:
\begin{multline}
 \hat{\textbf{\i}} = \hat{\textbf{e}}_x \\
 \hat{\textbf{\j}} = \hat{\textbf{e}}_y \\
 \hat{\textbf{k}}= \hat{\textbf{e}}_z \hfill
\end{multline}
and the resulting transformation matrix will be the identity matrix.

The transformation matrix
\begin{equation}
 \hfill R = \begin{bmatrix}
 \hat{\textbf{\i}}^T \\
 \hat{\textbf{\j}}^T \\
 \hat{\textbf{k}}^T
\end{bmatrix} \hfill
 \label{eq:ref_mat}
\end{equation}
can then be used to transform vectors in the inertial frame to the
user-specified reference-frame 
with the \texttt{generate\_inrtl\_to\_reference} method.


\subsubsection{Inertial to VNC}
The VNC frame is a rotating coordinate system where the x-axis points along
the velocity vector, the y-axis is in the direction of the angular momentum
vector, and the z-axis is conormal to complete the right-handed system. In 
order for these axes to be well-defined, it is assumed that the position
vector is nonzero and that the velocity vector is neither radial nor zero.
Assuming $\hat{\textbf{r}}$ and $\hat{\textbf{v}}$
are the unit position and velocity vectors in the inertial frame,
respectively, let
\begin{multline}
 \hat{\textbf{\i}} = \hat{\textbf{v}} \\
 \hat{\textbf{\j}} = \frac{\hat{\textbf{r}} \times \hat{\textbf{v}}}
{||\hat{\textbf{r}} \times \hat{\textbf{v}} ||} \\
 \hat{\textbf{k}} = \hat{\textbf{\i}} \times \hat{\textbf{\j}} \hfill
 \label{eq:vnc_main}
 \end{multline}
complete the definition of the VNC frame coordinate axes.
If the assumptions on the position and velocity vectors are violated, then
(\ref{eq:vnc_main}) cannot be used to align the VNC frame and alternative methods are needed.
As with the previous transformations, implementing an alternative method allows the simulation
to continue and generates an error message indicating the source and chosen method.

In the event of a zero velocity vector, the VNC x-axis wouldn't be defined, and the VNC frame
is aligned with the inertial frame. Hence, the transformation matrix will be the identity 
matrix. 

Similarly, if the position vector is zero or if the velocity vector is aligned with the 
position vector, the angular momentum vector won't be defined, and the VNC frame will be 
aligned with the inertial frame again.

Finally, the transformation matrix
\begin{equation}
 \hfill R = \begin{bmatrix}
 \hat{\textbf{\i}}^T \\
 \hat{\textbf{\j}}^T \\
 \hat{\textbf{k}}^T
\end{bmatrix} \hfill
 \label{eq:vnc_mat}
\end{equation}
can then be used to transform vectors in the inertial frame to the
VNC frame with the \\ \texttt{generate\_inrtl\_to\_vnc} method.

\subsubsection{Cartesian Planet-fixed to ENU Transformation Matrix}
In the \texttt{generate\_T\_pfix\_to\_enu} transformation, we start
with the position vector
$\vec r = [x, y,z] = x \hat i + y \hat j + z \hat k$.
The \textit{Up} ($\hat u$), \textit{East} ($\hat e$), and \textit{North}
($\hat n$) vectors can be written in terms of the position vector:
\begin{align*}
\hat{u} & = \hat{r}\\
\hat{e} & = \frac{(\hat{k} \times \hat{r})}{||(\hat{k} \times \hat{r})||} \\
\hat{n} & = \hat{u} \times \hat{e}
\end{align*}

Note that $\hat e$ must therefore be perpendicular to the
polar axis (equivalently, parallel to the equatorial plane) and $\hat n$
comprises two components -- one parallel to the equatorial plane pointing to
the polar-axis, and one parallel to the polar axis.

The transformation matrix is then:
\begin{equation}
 T = \begin{bmatrix} \hat e ^T \\
 \hat n ^T \\
 \hat u ^T
 \end{bmatrix}
\end{equation}

In the event of a zero position vector, or a position vector along the polar
axis ($\vec r = [0,0,z]$), the ENU frame is not defined.
\textbf{For a zero position
vector, the transformation matrix will be assigned the zero matrix}.
For a polar
position, $\hat u$ will still be $\hat r$, $[0,0, \pm 1]$ depending on whether
the position is over the north or south pole,
$\hat e$ will be aligned with $\hat i$, and $\hat n$
completes. The resulting transformation matrix will be:
\begin{equation*}
 T = \begin{bmatrix} 1 &0 &0 \\
 0 &\pm 1 &0 \\
 0 &0 &\pm 1
 \end{bmatrix}
\end{equation*}

Note -- the behavior of the method for a zero-position input may be
re-considered at some point in the future.
Whether it is preferable to assign an identity matrix or a zero
matrix is debatable.

\subsubsection{ENU to Cartesian Planet-fixed Left-Transformation Quaternion}
The primary difference between the \texttt{generate\_Q\_enu\_to\_pfix} method and 
the previous \\ \texttt{generate\_T\_pfix\_to\_enu} method for
generating the transformation matrix between ENU and planet-fixed is in the
inputs -- the previous method used a cartesian 3-vector, this uses a
latitude-longitude pair. The other differences -- that one is ENU-to-Pfix
and the other Pfix-to-ENU, and that one produces a matrix while the other
produces a quaternion -- are insignificant; the direction is easily reversed,
and there are algorithms available to convert between transformation
matrices and their equivalent quaternions.

The derivation of the generation from the latitude-longitude pair is more
easily followed with visuals.

In the figure below, the Cartesian planet-fixed frame origin is placed
at planet-center and labeled \texttt{[x,y,z]}; hereafter the axes of
frames in this figure are colored in order [blue, green, red].
A point of interest is identified with an orange star. The
longitude $\lambda$ and latitude $\phi$ of this point are shown. In the upper
right, the orientation of the ENU frame at that point of interest is shown.

\begin{figure}[h!]
 \label{fig:lon_lat}
 \begin{center}
 \includegraphics[scale=0.5]{Figs/lon-lat.png}
 % lon-lat.png: 554x586 px, 95dpi, 14.81x15.67 cm, bb=0 0 420 444
 \caption{Generation of ENU from longitude-latitude}
 \end{center}
\end{figure}



We want the transformation from \textit{ENU} to \textit{Pfix}, but it is
visually simpler to start be generating the reverse transformation,
from \textit{Pfix} to \textit{ENU}.
In the lower left, we start with frame $S_1$ with an orientation matching that
of the original Cartesian \textit{Pfix} frame.
This frame is then rotated about the z-axis (of $S_1$) through the
angle $\lambda$ to
create frame $S_2$; the x-axis of $S_2$ is aligned with the vector
joining the planet-center to the projection of the point of interest onto the
equator and the z-axis is still aligned with the polar axis.
Frame $S_2$ is then rotated by angle $\phi$ about the negative-y-axis
(of $S_2$) to create the \textit{UEN} (Up-East-North) frame;
the x-axis of the \textit{UEN} frame
is aligned with the position vector of the point of interest.
The \textit{UEN} frame is
then reordered, swapping axes to create the \textit{ENU} frame.

The left-rotation quaternion, used to rotate a vector by $\lambda$ about the
z-axis is
\begin{equation*}
 Q_{LR, \lambda, z} =
 \begin{bmatrix}
 cos \frac{\lambda}{2} \\
 sin \frac{\lambda}{2}
 \begin{bmatrix}
 0 \\ 0 \\ 1
 \end{bmatrix}
 \end{bmatrix} \\
\end{equation*}
Equivalently, this is $Q_{RT, \lambda, z}$,
the right-transformation quaternion used to transform between two frames
using the same$\lambda$ rotation about the z-axis. For our purposes, this is
$Q_{RT,(1\rightarrow2)}$, used to transform the representation of a vector from
using the $S_1$ frame basis vectors (and therefore the \textit{Pfix} frame
basis vectors) into using the $S_2$ frame basis vectors.

Similarly, the right-transformation quaternion
from the $S_2$ frame into the \textit{UEN} frame is:
\begin{equation*}
 Q_{RT,(2\rightarrow UEN)} =
 \begin{bmatrix}
 cos \frac{\phi}{2} \\
 sin \frac{\phi}{2}
 \begin{bmatrix}
 0 \\ -1 \\ 0
 \end{bmatrix}
 \end{bmatrix} \\
\end{equation*}

Finally, we want to map
\begin{equation*}
 \begin{bmatrix} a \\ b \\ c \end{bmatrix}_{UEN}
 \rightarrow
 \begin{bmatrix} b \\ c \\ a \end{bmatrix}_{ENU}
\end{equation*}

The right-transformation quaternion for this mapping is:
\begin{equation*}
 Q_{RT,(UEN \rightarrow ENU)} =
 \begin{bmatrix}
 0.5 \\
 \begin{bmatrix}
 0.5 \\ 0.5 \\ 0.5
 \end{bmatrix}
 \end{bmatrix}
\end{equation*}

Now, compounding these quaternions (to the right, they are right-transformation
quaternions), we have:
\begin{align*}
 Q_{RT, (Pfix \rightarrow ENU)} &=
 ~Q_{RT, (1 \rightarrow ENU)}~ =
 ~Q_{RT, (1 \rightarrow 2)}\otimes
 Q_{RT, (2 \rightarrow UEN)} \otimes
 Q_{RT, (UEN \rightarrow ENU)} \\
 \\
 &= \begin{bmatrix}
 cos \frac{\lambda}{2} \\
 sin \frac{\lambda}{2}
 \begin{bmatrix}
 0 \\ 0 \\ 1
 \end{bmatrix}
 \end{bmatrix} \otimes
 \begin{bmatrix}
 cos \frac{\phi}{2} \\
 sin \frac{\phi}{2}
 \begin{bmatrix}
 0 \\ -1 \\ 0
 \end{bmatrix}
 \end{bmatrix} \otimes
 \begin{bmatrix}
 0.5 \\
 \begin{bmatrix}
 0.5 \\ 0.5 \\ 0.5
 \end{bmatrix}
 \end{bmatrix}
\end{align*}


The left-transformation quaternion $Q_{LT, (ENU \rightarrow Pfix)}$ is identical
to this result (a left-transformation is conjugate to a right-transformation
quaternion; an inverse quaternion is conjugate to its original; conjugating
twice returns the original). Thus:
\begin{equation}
Q_{LT, (ENU \rightarrow Pfix)} =
 \begin{bmatrix}
 cos \frac{\lambda}{2} \\
 sin \frac{\lambda}{2}
 \begin{bmatrix}
 0 \\ 0 \\ 1
 \end{bmatrix}
 \end{bmatrix} \otimes
 \begin{bmatrix}
 cos \frac{\phi}{2} \\
 sin \frac{\phi}{2}
 \begin{bmatrix}
 0 \\ -1 \\ 0
 \end{bmatrix}
 \end{bmatrix} \otimes
 \begin{bmatrix}
 0.5 \\
 \begin{bmatrix}
 0.5 \\ 0.5 \\ 0.5
 \end{bmatrix}
 \end{bmatrix}
\end{equation}

%******************************************************************************
%******************************************************************************
\subsection{Protection Methods}

The protection methods utilize standard functions when the input is
within domain and apply pre-defined values when the input is out of domain.

%******************************************************************************
%******************************************************************************
\subsection{Numerical Methods}

\subsubsection{is\_near\_equal}

This method evaluates whether two floating-point
values are identically equal by comparing their difference against an
automatically-computed threshold value representing a comparatively very small
number.

The threshold value is equal to some (user-assignable) multiple of the
unit-of-least-precision (\textit{ULP}) associated with one of the two
values.

\begin{quote}
Aside -- The unit-of-least-precision is the difference between two consecutive
numbers, as represented in the computer architecture. This \textit{ULP}
value varies by data type and by the value of the variable (e.g. for type
\texttt{double}, it is typically somewhere close to the 16th significant
digit). It is typically identified by multiplying the value with
\texttt{std::numeric\_limits<T>::epsilon()} where \texttt{T} is the template
parameter identifying the data type.
\end{quote}

The first two arguments to the method are the two numbers being compared.
The third argument (i.e. \texttt{ulp}) is optional and represents
the multiple of the identified \textit{ULP} that will define the threshold
for comparing the difference between the two values. If not provided, it
defaults to one-half (0.5) meaning that the two values will be considered
identical if their difference is less than one-half of the identified
\textit{ULP}.

The two values are thus considered equal if the absolute-value of the
difference between them satisfies the following conditions:
\begin{itemize}
\item when both values are non-zero and either is a normal number, they
are considered equal when the absolute-value of the difference
between them \textbf{is less than} the product of:
\begin{itemize}
\item the \texttt{ulp} argument (default 0.5) and
\item the \textit{ULP} of the value with
the larger absolute-value, i.e. the product of:
\begin{itemize}
\item \texttt{std::numeric\_limits<T>::epsilon()} and
\item \texttt{std::max( std::abs(val1), std::abs(val2))}
\end{itemize}
\end{itemize}
\item when \textit{either} value is zero, or when both values are
sub-normal numbers (note -- this happens very rarely),
the two
values are considered equal when the absolute-value of the difference
between them \textbf{is not greater than} the product of:
\begin{itemize}
\item the \texttt{ulp} argument (default 0.5) and
\item the \textit{ULP} of the smallest value normally available in
the specified data type, i.e. the product of:
\begin{itemize}
\item \texttt{std::numeric\_limits<T>::epsilon()} and
\item \texttt{std::numeric\_limits<T>::min()}
\end{itemize}
\end{itemize}
\end{itemize}
\begin{quote}
Aside -- a normal number is one that has magnitude no less than \newline
\texttt{std::numeric\_limits<T>::min()}; a sub-normal number is one that
has magnitude \newline less than
\texttt{std::numeric\_limits<T>::min()} but is still non-zero.
\end{quote}

In these expressions, \texttt{<T>} is the template parameter providing the
distinction between data types; it will be either \texttt{float} or
\texttt{double}.

\subsubsection{has\_changed\_from}

This method determines whether a variable has changed
from a previously stored value. It compares two values of the same generic type and returns
\texttt{true} if they differ. 

The implementation simply checks if one value is 
greater than the other (i.e. $(v_1 < v_2) \lor (v_2 < v_1)$). 
This formulation ensures compatibility with types that may not define the
\texttt{!=} operator.

Mathematically, this would be written as:
\[
\text{has\_changed\_from}(v_1, v_2) =
\begin{cases}
\text{true}, & \text{if } v_1 \ne v_2 \\
\text{false}, & \text{otherwise}
\end{cases}
\]

\subsubsection{is\_equal}

This method tests equality between two values of the same generic type, and 
returns \texttt{true} is they do not differ. 
This test of equality is best suited for comparing previously assigned 
floating-point values, not independently computed results. The 
\texttt{is\_within\_abs\_tolerance} method is a better choice for comparing against
independent values.

For most types, it directly applies the equality operator (i.e. $v_1 == v_2$).

However, for floating-point types (\texttt{float}, \texttt{double}), this
method uses a specialized comparison based on negating the previously described
\texttt{is\_near\_equal} logic to logically test: 
$\lnot(v_1 < v_2) \land \lnot(v_2 < v_1)$.

Mathematically, this would be written as:
\[
\text{is\_equal}(v_1, v_2) =
\begin{cases}
\text{true}, & \text{if } v_1 = v_2 \\
\text{false}, & \text{otherwise}
\end{cases}
\]

\subsubsection{is\_within\_abs\_tolerance}
This method extends the \texttt{is\_near\_equal} capability 
by allowing an absolute tolerance to be specified directly, rather 
than a multiple of the unit-of-least-precision. It compares two values, \textit{x} 
and \textit{y} using a provided absolute threshold, \textit{tol}. Although templated 
for generic types, all arguments must be of the same type. It returns \texttt{true} if
the values are within the specified absolute tolerance, 
according to the following logic:
\begin{enumerate}
\item if $tol < 0$, return false and broadcast an error
\item else if $y>x$, return {$x+tol \geqslant y$}
\item else, return {$y+tol \geqslant x$}
\end{enumerate}

This is mathematically equivalent to:
\[
|x - y| \le tol \quad \text{for } tol \ge 0
\]
but is implemented without direct subtraction to avoid issues with unsigned
types.

For example, testing whether -2 and +1 are within 3 of each other:
\begin{multline*}
x=-2 \quad y=+1 \quad tol=3 \\\\
\text{1.} \quad tol < 0 \quad \texttt{false} \\
\text{2.} \quad y > x \quad \texttt{true} \\
\qquad -2 + 3 \geqslant 1 \quad \texttt{return true} \hfill
\end{multline*}

\subsubsection{is\_within\_rel\_tolerance}
This method extends the \texttt{is\_within\_abs\_tolerance} capability by 
determining whether a value lies within a specified \textit{relative}
tolerance of an expected reference value. It compares some \textit{value} to be within
some \textit{relative tolerance} of the \textit{expected} value. 
The two values to compare are templated for generic types, but the 
\textit{relative tolerance} must be of type \texttt{double}. According to the below logic,
it returns \texttt{true} if the values are witihin the specified relative tolerance.

The implementation simply converts the relative tolerance to 
an absolute tolerance by scaling it by the magnitude of the expected value, 
and delegates the absolute tolerance check to \texttt{is\_within\_abs\_tolerance}.
This approach will again ensure consistent behavior for unsigned types, but will 
still allow generic data types to be relatively compared to the tolerance.

Mathematically:
\[
|v - v_{\text{expected}}| \le |v_{\text{expected}}| \times tol
\]
where $tol$ is a non-negative real number.

If $tol < 0$, the method returns false and may broadcast an error. 

\subsubsection{sign}
This method simply returns the sign of a given value. It evaluates to
\texttt{1} if the value is positive, \texttt{-1} if negative, and \texttt{0} if
zero. The input value can be any generic data type that supports comparison with
a "zero" value, as the implementation simply evaluates the condition:
$(x > 0) - (x < 0)$

Mathematically:
\[
\text{sign}(x) =
\begin{cases}
1, & \text{if } x > 0 \\
0, & \text{if } x = 0 \\
 -1, & \text{if } x < 0
\end{cases}
\]

\subsubsection{Polynomial}
The \texttt{polynomial} method is mathematically simple.It receives:
\begin{itemize}
 \item \textit{\textbf{x}}, the value of the variable
 \item \textit{\textbf{coeffs}}, an STL-vector of polynomial coefficients
\end{itemize}
The algorithm starts by setting two values:
\begin{itemize}
 \item \textit{x\_to\_i = 1}
 \item\textit{sum = 0}
\end{itemize}
Processing through the coefficients one at a time, from front to back:
\begin{itemize}
 \item the variable \textit{sum} is incremented by the product of the
 current coefficient and the current value of \textit{x\_to\_i}
 \item the variable \textit{x\_to\_i} is multiplied by the value of $x$
\end{itemize}
which results in the simple polynomial summation:
\begin{equation}
 f(x) = \sum_{i=0}^{N} c_i x^i
\end{equation}
where $c_i$ is the $i^{th}$ component of the \textit{coeffs} vector


\subsubsection{Quadratic Solver}
The \texttt{QuadraticSolver::update()} method takes specified values
for the three coefficients of
\begin{equation*}
ax^2 + bx + c = 0
\end{equation*}

and applies the standard ``quadratic formula'':
\begin{equation*}
x = \frac{ -b \pm \sqrt{b^2 - 4 ac}}{2a}
\end{equation*}

If the discriminant, $b^2 - 4 ac$ is negative, the roots are not computed and the method returns
\texttt{false}. Otherwise, the method returns \texttt{true}, and the roots are stored in class 
members \texttt{root1} and \texttt{root2}.

\paragraph{Overflow Handling}
If $b^2$ overflows (i.e., evaluates to \texttt{inf}), the coefficients are normalized by dividing 
$a$ and $c$ by $b$, and setting $b = 1$ to prevent undefined behavior.

\paragraph{Numerical Stability}
In cases where the discriminant is large due to the $b^2$ term (i.e. $b^2 >> 4ac$), the standard formula
becomes numerically unstable due to \textit{catastrophic cancellation}. This occurs when subtracting
two nearly equal values of $b$.
\begin{equation*}
x = \frac{ -b \pm \sqrt{b^2}}{2a} = \frac{ -b \pm b}{2a} = \left \{\frac{ -b - b}{2a}, \frac{ -b + b}{2a} \right \} = \left \{-\frac{b}{a}, 0 \right \}
\end{equation*}
In this case, the smaller root (near zero) is lost due to precision errors, while the larger
root remains accurate.

To resolve this, a first-order Taylor series expansion is used when $\frac{4ac}{b^2} << 1$. The 
quadratic formula is reformulated as:
\begin{equation*}
x = \frac{ -b \pm \sqrt{b^2 - 4 ac}}{2a} = \frac{-b}{2a} \left (1 \pm \sqrt{1 - \frac{4ac}{b^2}} \right )
\end{equation*}

Expanding the square root term using a first-order approximation:
\begin{equation*}
\sqrt{1 - \epsilon} \approx 1 - \frac{1}{2}\epsilon \quad \text{where } \epsilon = \frac{4ac}{b^2}
\end{equation*}

Simplifying to:
\begin{equation*}
x \approx -\frac{b}{a} \left\{1 - \frac{ac}{b^2}, \frac{ac}{b^2}\right\}
\end{equation*}

Using the $\frac{4ac}{b^2} << 1$, produces the roots:
\begin{equation*}
x \approx \left\{-\frac{b}{a}, -\frac{c}{b}\right\}
\end{equation*}
This formulation provides a more stable estimate of both roots when the $b^2$ coefficient is 
much larger than that of $a$ or $c$.

\paragraph{Root Assignment}
\begin{itemize}
  \item If only one root can be safely computed (e.g., due to overflow risk in $b/a$ or $c/b$), it 
  is assigned to both \texttt{root1} and \texttt{root2}.
  \item If both roots are computable, the smaller root (typically $-c/b$) is assigned to \texttt{root1}, 
  and the larger root (typically $-b/a$) to \texttt{root2}.
  \item If the discriminant doesn't produce any numerical instability, the standard quadratic formula is
  used directly: the root computed using the ``$-b - \sqrt{b^2 - 4ac}$'' term is assigned to \texttt{root1}, 
  and the root using ``$-b + \sqrt{b^2 - 4ac}$'' is assigned to \texttt{root2}.
\end{itemize}

This formulation ensures robust root computation across a range of coefficient values.

%******************************************************************************
%******************************************************************************
\subsection{Matrix and Vector Methods}

\subsubsection{Cholesky Decomposition}
The \texttt{cholesky\_decomposition} method provides the decomposition of
some square symmetric matrix
$A$ into its ``square-root'' matrix $S$ such that:
\begin{equation*}
A = S S^T
\end{equation*}
The matrix $S$ has its lower-left components (on and below the diagonal)
populated and the upper-right components set to 0.

\paragraph{Nomenclature}
$S_{ij}$ is the element of matrix $S$ in row $i$ and column $j$.

\paragraph{Positive-definite $A$}
The requirements on $A$ is that it be symmetric (or contain all zeros
above the diagonal) and positive semi-definite.Just as the square-root
of a negative number is not real, the decomposition of a matrix with
negative determinant is not real.The decomposition of a matrix with a
zero determinant is ambiguous, this is covered in the next section; for
now we assume that the determinant is positive.

The algorithm proceeds by computing the values of $S$ one column at a time,
moving down within each column from the diagonal element to the bottom of
the matrix.The diagonal element is computed based on the diagonal element
from $A$ and all elements to the left of the diagonal element in $S$:
\begin{equation}
 S_{ii} = \sqrt{A_{ii} - \sum_{j=0}^{i-1} S_{ij}^2}
\end{equation}

The off-diagonal elements $S_{ji}$ are computed from:
\begin{itemize}
 \item the corresponding value from $A$, $A_{ji}$
 \item the values of $S$ to the left of the current element,
 $S_{jk}$ for $k < i$
 \item the value of the diagonal element of $S$ in the same column $i$
 (so above the current element), $S_{ii}$
 \item the values of $S$ to the left of that diagonal element,
 $S_{ik}$ for $k < i$
\end{itemize}
\begin{equation}
 S_{ji} = \frac{A_{ji} - \sum_{k=0}^{i-1} S_{ik} * S_{jk}}{S_{ii}}
\end{equation}


\paragraph{Zero determinant A}
When the matrix $A$ has a zero determinant, the evaluation of $S_{ii}$ at
some point will return $0.0$, which would cause a divide-by-zero error in
the computation of the elements below the diagonal. In this case, all
values in column $i$ ($S_{ji},\forall j$) are assigned to $0.0$ and the
computation continues with the next column. Note that this results in one
of many potential solutions for $S$.

\subsubsection{Position-Velocity Covariance Matrix Transformation}

The \texttt{transform\_pv\_matrix} method transforms a $6 \times 6$ 
position-velocity (PV) covariance matrix from one coordinate frame (A) 
to another (B) given the $3 \times 3$ transformation matrix.

\paragraph{Nomenclature}
\begin{itemize}
\item Vectors and matrices are denoted in bold (e.g. $\mathbf{X}$).
\item Superscripts are used to provide the reference frame in which a vector
or matrix is expressed (e.g. $\mathbf{X}^B$ describes a vector expressed
in the $B$ frame).
\item $\mathbf{P_{XX}}$ represents the covariance matrix of a random vector, 
$\mathbf{X}$, which quantifies the spread or uncertainty of the 
state about its mean.
\item $\mathbb{E}[\cdot]$ denotes the expected value operator and $\boldsymbol{\mu_{X}}$ 
denotes the expected value (mean) of a random vector $\mathbf{X}$.
\end{itemize}

\paragraph{Background} 
For arbitrary position and velocity vectors:
\[
\mathbf{x} = 
\begin{bmatrix}
x \\ y \\ z
\end{bmatrix}, \qquad
\mathbf{v} = 
\begin{bmatrix}
v_x \\ v_y \\ v_z
\end{bmatrix}
\]
the combined position-velocity state vector can be expressed as:
\[
\mathbf{X} = 
\begin{bmatrix}
\mathbf{x} \\ \mathbf{v}
\end{bmatrix}
\]

The corresponding $6 \times 6$ covariance matrix of this state is:
\[
\mathbf{P}_{\text{PV}} =
\left[
\begin{array}{c|c}
\begin{array}{ccc}
P_{x,x} & P_{x,y} & P_{x,z} \\
P_{y,x} & P_{y,y} & P_{y,z} \\
P_{z,x} & P_{z,y} & P_{z,z}
\end{array}
&
\begin{array}{ccc}
P_{x,v_x} & P_{x,v_y} & P_{x,v_z} \\
P_{y,v_x} & P_{y,v_y} & P_{y,v_z} \\
P_{z,v_x} & P_{z,v_y} & P_{z,v_z}
\end{array}
\\ \hline
\begin{array}{ccc}
P_{v_x,x} & P_{v_x,y} & P_{v_x,z} \\
P_{v_y,x} & P_{v_y,y} & P_{v_y,z} \\
P_{v_z,x} & P_{v_z,y} & P_{v_z,z}
\end{array}
&
\begin{array}{ccc}
P_{v_x,v_x} & P_{v_x,v_y} & P_{v_x,v_z} \\
P_{v_y,v_x} & P_{v_y,v_y} & P_{v_y,v_z} \\
P_{v_z,v_x} & P_{v_z,v_y} & P_{v_z,v_z}
\end{array}
\end{array}
\right]
= 
\begin{bmatrix}
\mathbf{P_{xx}} & \mathbf{P_{xv}} \\
\mathbf{P_{vx}} & \mathbf{P_{vv}}
\end{bmatrix}
\]
where:
\begin{itemize}
\item $\mathbf{P_{xx}}$ is the $3 \times 3$ position covariance submatrix,
\item $\mathbf{P_{vv}}$ is the $3 \times 3$ velocity covariance submatrix,
\item $\mathbf{P_{xv}}$ and $\mathbf{P_{vx}}$ represent the cross-covariances between position and velocity.
\end{itemize}

\paragraph{Covariance Formulation}
Given a transformation from coordinate frame $A$ to $B$, 
represented by an orthonormal $3 \times 3$ transformation matrix 
$\mathbf{T}_{B/A}$ where $\mathbf{T}_{B/A}^{-1} = \mathbf{T}_{B/A}^T$, 
the position and velocity vectors expressed in frame $A$ can be 
transformed to the $B$ frame as:
\[
\mathbf{x}^B = \mathbf{T}_{B/A} \mathbf{x}^A, \quad
\mathbf{v}^B = \mathbf{T}_{B/A} \mathbf{v}^A
\]
The covariance matrix follows the general covariance formula:
\[
\mathbf{P_{XX}} = \text{cov}(\mathbf{X},\mathbf{X}) = \mathbb{E}[(\mathbf{X}-\mu_\mathbf{X})(\mathbf{X}-\mu_\mathbf{X})^T]
\]

This can be broken down into each submatrix for simplification. For instance, 
for the position covariance submatrix expressed in frame $B$:
\begin{align*}
\mathbf{P}_{xx}^B 
&= \mathbb{E}[(\mathbf{x}^B - \mu_x^B)(\mathbf{x}^B - \mu_x^B)^T] \\
&= \mathbb{E}[(\mathbf{T}_{B/A}\mathbf{x}^A - \mathbf{T}_{B/A}\mu_x^A)(\mathbf{T}_{B/A}\mathbf{x}^A - \mathbf{T}_{B/A}\mu_x^A)^T] \\
&= \mathbb{E}[\mathbf{T}_{B/A}(\mathbf{x}^A - \mu_x^A)(\mathbf{T}_{B/A}(\mathbf{x}^A - \mu_x^A))^T] \\
&= \mathbb{E}[\mathbf{T}_{B/A}(\mathbf{x}^A - \mu_x^A)(\mathbf{x}^A - \mu_x^A)^T\mathbf{T}_{B/A}^T] \\
&= \mathbf{T}_{B/A}\mathbb{E}[(\mathbf{x}^A - \mu_x^A)(\mathbf{x}^A - \mu_x^A)^T]\mathbf{T}_{B/A}^T \\
&= \mathbf{T}_{B/A}\,\mathbf{P}_{xx}^A\,\mathbf{T}_{B/A}^T
\end{align*}

By the same logic, all the submatrices transform as:
\begin{equation}
\begin{aligned}
\mathbf{P}_{xx}^B &= \mathbf{T}_{B/A}\,\mathbf{P}_{xx}^A\,\mathbf{T}_{B/A}^T, \\
\mathbf{P}_{xv}^B &= \mathbf{T}_{B/A}\,\mathbf{P}_{xv}^A\,\mathbf{T}_{B/A}^T, \\
\mathbf{P}_{vx}^B &= \mathbf{T}_{B/A}\,\mathbf{P}_{vx}^A\,\mathbf{T}_{B/A}^T, \\
\mathbf{P}_{vv}^B &= \mathbf{T}_{B/A}\,\mathbf{P}_{vv}^A\,\mathbf{T}_{B/A}^T
\end{aligned}
\end{equation}

The implementation of \texttt{transform\_pv\_matrix} performs these 
transformations independently for each $3\times 3$ submatrix and reassembles
 them into a single transformed $6\times 6$ covariance matrix:
\[
\mathbf{P}_{\text{PV}}^B =
\begin{bmatrix}
\mathbf{P}_{xx}^B & \mathbf{P}_{xv}^B \\
\mathbf{P}_{vx}^B & \mathbf{P}_{vv}^B
\end{bmatrix}
\]

\subsubsection{Correlation Matrix}

The \texttt{extract\_correlation\_coefficients} method computes the correlation
matrix corresponding to a given covariance matrix. The correlation matrix 
expresses the normalized linear relationships between each pair of variables
in the covariance matrix, with values ranging from -1 (perfect negative 
correlation) to +1 (perfect positive correlation).

This method accepts a square covariance matrix of arbitrary size, where each 
element is a \texttt{double}. It assumes the covariance matrix is symmetric and 
that at least the lower diagonal is populated. The output is a correlation matrix
of the same size, where each correlation coefficient, $\rho_{ij}$ is computed as:

\begin{equation}
\rho_{ij} = \frac{P_{ij}}{\sqrt{P_{ii}} \cdot \sqrt{P_{jj}}}
\label{eqn:corr_coeff}
\end{equation}
where $P_{ij}$ is the covariance between variables $i$ and $j$, and 
$P_{ii}$ and $P_{jj}$ are the variances of variables $i$ and $j$, respectively.

The method implementation performs the following:
\begin{enumerate}
\item Extracts the square roots of the diagonal elements (variances). If 
any diagonal element is negative, the function returns \texttt{false} 
and logs an error, as variables cannot be negatively correlated with
themselves. 
\item Furthermore, if a diagonal element is zero, it checks that the corresponding 
row and column are also zero (i.e., make sure the variables don't have 
covariance without variance). 
If this condition is met, the square root is set to 1.0 to avoid 
division by zero, and a message is logged to inform this.
\item Computes the correlation coefficients for all off-diagonal elements 
using the normalized covariance formula from (\ref{eqn:corr_coeff}).
\item Sets the diagonal of the correlation matrix to 1.0. It is important to
note that a variable with zero variance would produce an
undefined correlation coefficient according to (\ref{eqn:corr_coeff}), 
but its correlation coefficient will be set to 1.0 here.
\end{enumerate}

The method returns \texttt{true} if the correlation matrix is successfully 
computed, and \texttt{false} if the input covariance matrix is invalid.

\subsubsection {Matrix and Vector Algebra}
The templated vector and matrix methods follow standard rules for matrix
algebra.

\paragraph{Nomenclature}
\begin{itemize}
\item Subscripts on a vector denote the element at that position in the vector 
following a zero-based index (e.g. $X_0$ represents the first element in a 
vector $X$). Similarly, subscripts on a matrix denote the element at that row 
and column, respectively (e.g. $X_{ij}$ represent the element in row $i$ 
and column $j$ in a matrix $X$).
\item Any formulation that uses variable indices is assumed to apply to all valid
indices of that vector or matrix (e.g. $X_{ij} = ... \quad \forall i,j$).
\end{itemize}

\textbf{Zero Vector}

The \texttt{zero\_vector} is populated as:
\begin{equation*}
X_i = 0
\end{equation*}

\textbf{Unit Vector}

The \texttt{unit\_vector} of a vector $X$ is the vector:
\begin{equation*}
\hat{X}_i = 
\begin{cases}
\frac{X_i}{||X||}, & \text{if } ||X|| > 0 \\
0, & \text{otherwise}
\end{cases}
\end{equation*}
where
\begin{equation*}
||X||=\sqrt{\sum_i X_i^2}
\end{equation*}
which is computed from \texttt{vec\_mag}

\textbf{Scalar Product of Two Vectors}

The \texttt{vector\_scalar\_product} of two vectors $X$ and $Y$ is
a scalar:
\begin{equation*}
X \cdot Y = \sum_i X_i Y_i
\end{equation*}

\textbf{Cross Product of Two Vectors}

The \texttt{vector\_cross\_product} of two 3D vectors $X$ and $Y$ is
the vector:
\begin{equation*} 
Z = X \times Y = 
\begin{bmatrix} 
X_1 Y_2 - X_2 Y_1 \\
X_2 Y_0 - X_0 Y_2 \\
X_0 Y_1 - X_1 Y_0 
\end{bmatrix} 
\end{equation*}

\textbf{Vector Magnitude}

The \texttt{vec\_mag} of a vector $X$ is a scalar:
\begin{equation*} 
||X||=\sqrt{\sum_i X_i^2}
\end{equation*}
which is computed from taking the square-root of \texttt{vec\_mag\_sq}
with $X$ as the input vector. 

\textbf{Vector Magnitude Squared}

The \texttt{vec\_mag\_sq} of a vector $X$ is a scalar:
\begin{equation*} 
||X||^2=\sum_i X_i^2
\end{equation*}
which is computed from \texttt{vector\_scalar\_product} with $X$ as
both input vectors.

\textbf{Vector X-Y Projection Magnitude}

The \texttt{vec\_mag\_xy} of a vector $X$ is a scalar:
\begin{equation*} 
||X_{xy}||=\sqrt{X_0^2 + X_1^2}
\end{equation*}

\textbf{Vector X-Z Projection Magnitude}

The \texttt{vec\_mag\_xz} of a vector $X$ is a scalar:
\begin{equation*} 
||X_{xz}||=\sqrt{X_0^2 + X_2^2}
\end{equation*}

\textbf{Vector Y-Z Projection Magnitude}

The \texttt{vec\_mag\_yz} of a vector $X$ is a scalar:
\begin{equation*} 
||X_{yz}||=\sqrt{X_1^2 + X_2^2}
\end{equation*}

\textbf{Vector Element-wise Square-Root}

The \texttt{vector\_elementwise\_sqrt} of a vector $X$ is the vector:
\begin{equation*} 
Y_i = \sqrt{X_i}
\end{equation*}

\textbf{Vector Addition/Difference}

The \texttt{sum/diff} of two vectors $X$ and $Y$ ($X \pm Y$) is the vector:
\begin{equation*} 
Z_i = X_i \pm Y_i
\end{equation*}

\textbf{Column-Major Vector to Matrix}

The \texttt{col\_maj\_vec\_to\_matrix} of a vector $X$ with length
$N*M$ into a matrix of size $N \times M$ arranged in column-major format is:
\begin{equation*} 
Y_{ij} = X_{i + Nj}
\end{equation*}

\textbf{Row-Major Vector to Matrix}

The \texttt{row\_maj\_vec\_to\_matrix} of a vector $X$ with length
$N*M$ into a matrix of size $N \times M$ arranged in column-major format is:
\begin{equation*} 
Y_{ij} = X_{Mi + j}
\end{equation*}

\textbf{Zero Matrix}

The \texttt{zero\_matrix} is populated as:
\begin{equation*}
X_{ij} = 0
\end{equation*}

\textbf{Matrix Scale}

The \texttt{matrix\_scale} of a matrix $X$ by a scalar $C$ ($C*[X]$) is the matrix:
\begin{equation*}
Z_{ij} = C * X_{ij}
\end{equation*}

\textbf{Matrix Increment/Decrement}

The \texttt{matrix\_incr/matrix\_decr} of two matrices $X$ and $Y$ ($[X] \pm [Y]$) is
the matrix:
\begin{equation*}
Z_{ij} = X_{ij} \pm Y_{ij}
\end{equation*}

\textbf{Matrix Copy}

The \texttt{matrix\_copy} of a matrix $X$ is the matrix:
\begin{equation*}
Y_{ij} = X_{ij}
\end{equation*}

\textbf{Extract Submatrix}

The \texttt{matrix\_copy\_submatrix\_out} method extracts a submatrix of size
$N_2 \times M_2$ from a larger source matrix $X$ of size 
$N_1 \times M_1$ beginning at the upper-left index $(i_0,j_0)$ in $X$ to produce
the submatrix:
\begin{equation*}
Y_{ij} = 
\begin{cases}
X_{i+i_0, j+j_0}, & \text{if } i + i_0 < N_1 \text{ and } j + j_0 < M_1 \\
0, & \text{otherwise}
\end{cases}
\end{equation*}

Note that in the implementation, the loop indices iterate over the source matrix, $X$,
and the submatrix $Y$ is populated using offsets relative to the source matrix 
indices. This differs from the formulation above, which expresses the mapping 
in terms of the submatrix indices for easier understanding.

\textbf{Insert Submatrix}

The \texttt{matrix\_copy\_submatrix\_in} method inserts a submatrix $X$ of size
$N_1 \times M_1$ into a larger target matrix $Y$ of size 
$N_2 \times M_2$ beginning at the upper-left index $(i_0,j_0)$ in $Y$ to produce
the new target matrix:
\begin{equation*} 
Y_{ij} = 
\begin{cases} 
X_{i - i_0, j - j_0}, & \text{if } i \in [i_0, i_0 + N_1),~ j \in [j_0, j_0 + M_1) \\ 
Y_{ij}, & \text{otherwise} 
\end{cases} 
\end{equation*}

\textbf{Matrix Transpose}

The \texttt{matrix\_trans} of a matrix $X$ of size $N \times M$ ($[X]^T$)
is the matrix of size $M \times N$:
\begin{equation*}
Y_{ij} = X_{ji}
\end{equation*}

\textbf{Basic Matrix Product}

The \texttt{matrix\_mult} of two matrices $X$ and $Y$ ($[X][Y]$) of sizes $N \times M$ and
$M \times R$, respectively, produces the matrix of size $N \times R$:
\begin{equation*}
Z_{ij} = \sum_k^M X_{ik} \cdot Y_{kj}
\end{equation*}
The product is undefined if the matrices are not appropriately
dimensioned. Subsequently, the method call cannot be compiled if this condition 
is not met.

\textbf{Left-Transpose Product}

The \texttt{matrix\_mult\_left\_trans} of two matrices $X$ and $Y$ ($[X]^T[Y]$) 
of sizes $N \times M$ and
$N \times R$, respectively produces the matrix of size $M \times R$. 
This implementation consists of the following:
\begin{enumerate}
\item Transposes the left matrix $[X]$ using \texttt{matrix\_trans}.
\item Multiply the transposed matrix $[X]^T$ with the right matrix $[Y]$ using 
\texttt{matrix\_mult}.
\end{enumerate}
Similar to \texttt{matrix\_mult}, this method will ensure matrix dimensions are
compatible during compile time.

\textbf{Right-Transpose Product}

The \texttt{matrix\_mult\_right\_trans} of two matrices $X$ and $Y$ ($[X][Y]^T$) 
of sizes $N \times M$ and
$R \times M$, respectively produces the matrix of size $N \times R$:
\begin{equation*}
Z_{ij} = \sum_k^M X_{ik} \cdot Y_{jk}
\end{equation*}
using the \texttt{vector\_scalar\_product} method to calculate calculate the dot 
product. It is important to note that this method can use the dot product directly
unlike the other matrix multiplication variants, because these would require 
accessing columns instead of rows. As the columns aren't contiguous in memory already
like the rows, they would need to be copied over to another vector/array to use 
\texttt{vector\_scalar\_product}, which would be less efficient.

Similar to \texttt{matrix\_mult}, this method will ensure matrix dimensions are
compatible during compile time.

\textbf{Double-Transpose Product}

The \texttt{matrix\_mult\_trans\_trans} of two matrices $X$ and $Y$ ($[X]^T[Y]^T$) 
of sizes $N \times M$ and
$R \times N$, respectively produces the matrix of size $M \times R$. 
This implementation consists of the following:
\begin{enumerate}
\item Transposes the left matrix $[X]$ using \texttt{matrix\_trans}.
\item Multiply the transposed matrix $[X]^T$ with the transpose 
of the right matrix $[Y]^T$ using 
\texttt{matrix\_mult\_right\_trans}.
\end{enumerate}
Similar to \texttt{matrix\_mult}, this method will ensure matrix dimensions are
compatible during compile time.

\textbf{Matrix Transformation}

For matrices associated with specific basis vectors (e.g. an inertia tensor
specified using body-axes, or a covariance specified using LVLH axes), it is
necessary to transform the matrix when changing reference frames. 
To apply an orthonormal transformation matrix $[T]$ to a matrix $[X]$
using \texttt{matrix\_transformation}, the result is:
\begin{equation}
[Y] = [T][X][T]^T
\label{eqn:matrix_transform}
\end{equation}
For example, the angular momentum, $\vec L = [I] \vec \omega$ can be written
in frames $A$ or $B$.
\begin{align*}
\vec L_A &= [I]_A ~\vec \omega_A \\
\vec L_B &= [I]_B ~\vec \omega_B
\end{align*}

We can transform the vectors
\begin{align*}
\vec L_B&= [T]_{A \rightarrow B}~ \vec L_A \\
\vec \omega_B &= [T]_{A \rightarrow B}~ \vec \omega_A \hfill
\end{align*}

Consequently:
\begin{align*}
\vec L_B &= [T]_{A \rightarrow B}~ \vec L_A \\
 &= [T]_{A \rightarrow B}~ ([I]_A~ \vec \omega_A) \\
 &= [T]_{A \rightarrow B}~ ([I]_A ~
([T]^T_{A \rightarrow B}~~ \vec \omega_B)) \\
 &= ([T]_{A \rightarrow B}~ [I]_A~ [T]^T_{A \rightarrow B})~~
\vec \omega_B
\end{align*}
and
\begin{align*}
[I]_B= ([T]_{A \rightarrow B}~ [I]_A~ [T]^T_{A \rightarrow B})
\end{align*}

This implementation consists of the following:
\begin{enumerate}
\item Pre-multiply the transformation matrix by the matrix ($[T][X]$) using 
\texttt{matrix\_mult}.
\item Post-multiply that result by the transpose of the transformation matrix
($[T][X][T]^T$) using \texttt{matrix\_mult\_right\_trans}.
\end{enumerate}

\textbf{Matrix Inverse Transformation}

The \texttt{matric\_inverse\_transformation} method functions
the same as \texttt{matrix\_transformation}, but with the transformation being
applied as the inverse of the provided transformation matrix.
As the method assumes the transformation matrix is orthonormal, the transpose
of the matrix is equivalent to the inverse ($T^T = T^{-1}$).
Therefore, the transpose of the transformation matrix can simply substitude into
(\ref{eqn:matrix_transform}) to form:
\begin{equation*}
[Y] = [T]^T[X][T]
\end{equation*}
This implementation consists of the following:
\begin{enumerate}
\item Pre-multiply the transpose of the transformation matrix by the
matrix ($[T]^T[X]$) using \texttt{matrix\_mult\_left\_trans}.
\item Post-multiply that result by the transformation matrix
($[T]^T[X][T]$) using \texttt{matrix\_mult}.
\end{enumerate}

\subsubsection{Matrix Inversion (via Cholesky Decomposition)}

Taking the inverse of a matrix can be numerically expensive when using a
completely generic algorithm, especially so as the size of the matrix
increases. For symmetric matrices with a positive determinant, we can simplify
the process by utilizing the Cholesky decomposition already implemented in
this model. The following logic is implemented in the 
\texttt{matrix\_inv\_using\_cholesky} method.

Starting with a symmetric positive-definite matrix $[A]$, we use the
Cholesky decomposition to build the ``square-root'' half-matrix $[B]$ in the
lower-left diagonal.


\begin{equation*}
[A] = \begin{bmatrix} A_{00} & A_{10} & A_{20} & ... \\
A_{10} & A_{11} & A_{21} & ... \\
A_{20} & A_{21} & A_{22} & ... \\
 ...
\end{bmatrix}
\end{equation*}

\begin{equation*}
[B] = \begin{bmatrix} B_{00} &0 &0 & ... \\
B_{10} & B_{11} &0 & ... \\
B_{20} & B_{21} & B_{22} & ... \\
 ...
\end{bmatrix}
\end{equation*}

with:
\begin{equation*}
[A] = [B] [B]^T
\end{equation*}

Note that order is important here:
\begin{equation*}
[A] \neq [B]^T B]
\end{equation*}

The inverse of $[A]$:
\begin{equation*}
[A]^{-1} = ([B] [B]^T)^{-1} = ([B]^T)^{-1} [B]^{-1} = ([B]^{-1})^T
[B]^{-1} = C^T C
\end{equation*}
where $[C] = [B]^{-1}$.

This uses the property that the inverse of a
transpose matrix and the transpose of the inverse matrix are identical. For
demonstration of this property, consider:
\begin{align*}
([X]^T)^{-1}[X]^T = I = I^T = ([X] [X]^{-1})^T = {[X]^{-1}}^T [X]^T
~~~\Rightarrow~~~ ([X]^T)^{-1}=([X]^{-1})^T
\end{align*}

It then remains to find the inverse of the half-matrix $[B]$, which is
significantly easier than finding the inverse of a full-matrix.

Consider the product of $[B]$ with its inverse $[C]$:
\begin{equation*}
\begin{bmatrix} B_{00} &0 &0 & ... \\
B_{10} & B_{11} &0 & ... \\
B_{20} & B_{21} & B_{22} & ... \\
 ...
\end{bmatrix}
\begin{bmatrix} C_{00} & C_{01} & C_{02} & ... \\
C_{10} & C_{11} & C_{12} & ... \\
C_{20} & C_{21} & C_{22} & ... \\
 ...
\end{bmatrix}
= [I]
\end{equation*}
Starting in the top right and working down, then to the left:
\begin{align*}
 &\sum_i B_{0i} C_{i n} = B_{00} C_{0n} + 0 + ... + 0
 &=I_{0n} = 0
 &\Rightarrow C_{0n} = 0 \\
 &\sum_i B_{1i} C_{in)} = 0 + B_{11} C_{1n} + 0 + ... + 0
 &=I_{1n} = 0
 &\Rightarrow C_{1n} = 0 \\
 ... \\
 &\sum_i B_{0i} C_{i (n-1)} = B_{00} C_{0(n-1)} + 0 + ... + 0
 &=I_{0(n-1)} = 0
 &\Rightarrow C_{0(n-1)} = 0
\end{align*}
Thus $[C]$ is also a half-matrix, limited to the lower diagonal with $C_{ji}
= 0 ~~\forall j<i$.

The product of row $i$ of $[B]$ and column $i$ of $[C]$ simplify to $B_{ii}
C_{ii} = 1$ and the diagonal elements of $[C]$ consequently become
\begin{equation*}
C_{ii} = \frac{1}{B_{ii}}
\end{equation*}

The remaining off-diagonal elements can be expressed as an iterative sum,
obtained from the product of row $j$ of $[B]$ and column $i$ of $[C]$ with
$i<j$. Removing the products involving a zero in either $[B]$ and $[C]$
simplifies:
\begin{equation*}
\sum_{k=0}^{n} B_{jk}C_{ki} =
\sum_{k=i}^{j} B_{jk}C_{ki} = 0
\end{equation*}

Separating the last term of this sum ($k=j$) provides the expression for
$C_{ji}$ as an iterative sum:
\begin{equation*}
C_{ji} = -\frac{1}{B_{jj}} \sum_{k=i}^{j-1} B_{jk}C_{ki} ~~\forall j>i
\end{equation*}

Thus for each column in $[C]$, we can work down the column, setting the
values above the diagonal element to zero, setting the diagonal element as
the reciprocal of the same element from $[B]$, and sequentially generating
the values below the diagonal.

Then, as identified above, the inverse of $[A]$ is
\begin{equation}
[A]^{-1} = C^T C
\end{equation}


%******************************************************************************
%******************************************************************************
\subsection{Calculus Methods}

\subsubsection{Backward Difference}
The \texttt{compute\_backward\_difference} method 
uses a STL-list of historical points in determining the
backward difference, using a set of coefficients determined by the number
of historical points available. Note that this does not compute a
differential, $\frac {dy}{dx}$, it only computes the difference $dy$.
It is assumed that the historical points are equally spaced by some
distance $dx$ so the differential may be obtained by dividing the
result $dy$ by that spacing.

\begin{itemize}
\item Two points provided:
\begin{equation*}
df(x_n) = f(x_n) - f(x_{n-1})
\end{equation*}
\item Three points provided:
\begin{equation*}
df(x_n) = \frac{3 f(x_n) - 4 f(x_{n-1}) + f(x_{n-2})}{2}
\end{equation*}
\item Four points provided:
\begin{equation*}
df(x_n) = \frac{11 f(x_n) - 18 f(x_{n-1}) + 9 f(x_{n-2}) - 2f(x_{n-3})}{6}
\end{equation*}
\item Five (or more) points provided:
 \begin{equation*}
 df(x_n) = \frac{ 25 f(x_n) - 48 f(x_{n-1}) + 36 f(x_{n-2}) -
16 f(x_{n-3}) + 3 f(x_{n-4})}
{ 12}
 \end{equation*}
\end{itemize}

\subsubsection{Unit Vector Derivative}
The \texttt{compute\_unit\_vector\_derivative} method provides the derivative 
of some unit vector representation of
some vector $x$:
\begin{equation*}
 \dot {\hat{x}} = \frac{d}{dt} \left( \frac{ \vec x}
 { \rVert \vec x \rVert}
 \right)
\end{equation*}

Note this is not the same as the unit vector of the derivative,
$\hat {\dot{x}} = \frac{ \dot{\vec{x}}}
 { \rVert \dot{\vec{x}} \rVert}$
but it does require
knowledge of the derivative of the vector, $\dot{\vec x}$.
The unit vector is defined as:
\begin{equation*}
 \hat{x} = \frac{\vec{x}}{\sqrt{ \vec{x} \cdot \vec x}}
\end{equation*}
and therefore its derivative is:
\begin{equation*}
\begin{split}
\dot{\hat{x}} & = \frac{\dot{\vec{x}}}{\sqrt{ \vec x \cdot \vec x}} +
\frac{\vec x \left(-\frac{1}{2} \right)
\left(2 \vec{x} \cdot \dot{\vec{x}}\right)}
 {(\vec{x} \cdot \vec x)^{3/2}} \\
& = \frac{1}{\sqrt{ \vec x \cdot \vec x}}
\left(\dot{\vec x} - \vec x
 \left(\frac{ \vec x \cdot \dot{\vec x}}
{\vec x \cdot \vec x}
 \right)
\right)
\end{split}
\end{equation*}
As an illustrative example of the difference between $\dot{\hat x}$ and
$\dot{\vec x}$, consider the case in which $\vec x$ and $\dot{\vec x}$ are
aligned (e.g. $\vec x$ represents the position of a body moving radially
away from the frame origin). The
unit vector $\hat x$ does not change during this motion because $\vec x$
is not changing direction, so $\dot{\hat x} = \vec 0$ regardless of the
magnitude of $\dot{\vec x}$.

%******************************************************************************
%******************************************************************************
\subsection{Geometry Methods}
\subsubsection{Ellipsoid Intersection}

In the \texttt{EllipsoidIntersection::update()} method, we evaluate whether
a defined line intersects with a defined ellipsoid. It is assumed that the
ellipsoid is centered on the frame origin and the two end points are
specified in that same frame.

Consider two vectors defining the end-points of the line:
\begin{equation*}
\vec r_1 = [ x_1, y_1, z_1]
\vec r_2 = [ x_2, y_2, z_2]
\end{equation*}

and an ellipsoid defined by semi-axes ${a, b, c}$ such that:
\begin{equation*}
\left( \frac{x}{a} \right)^2 +
\left( \frac{y}{b} \right)^2 +
\left( \frac{z}{c} \right)^2= 1
\end{equation*}

We can parametrize the line with parameter $s$ such that every point on the
line can be represented as:
\begin{equation*}
\vec r = \vec r_1 + s( \vec r_2 - \vec r_1) = \vec r_1 (1-s) + \vec r_2 (s)
\end{equation*}

If we can identify a value s such that $\vec r$ satisfies the equation of
the ellipsoid, then the line must intersect with the ellipsoid. That is,
does a real solution exist to the following equation in s:
\begin{equation}\label{eqn:ellipsoid_intersect}
\left( \frac{x_1(1-s) + x_2 s}{a} \right)^2 +
\left( \frac{y_1(1-s) + y_2 s}{b} \right)^2 +
\left( \frac{z_1(1-s) + z_2 s}{c} \right)^2= 1
\end{equation}

The three terms can be expanded and simplified following the same pattern;
only the first term is shown here:
\begin{align*}
\left( \frac{x_1(1-s) + x_2 s}{a} \right)^2 & =
\frac{1}{a^2} \left( x_1^2(1-s)^2 + 2 x_1 x_2 s (1-s) + x_2 s^2 \right) \\
{} & = \frac{1}{a^2} \left( (x_2 - x_1)^2 s^2 + 2 x_1 (x_2 - x_1) + x_1^2 \right) \\
{} & = \left( \frac{x_2 - x_1}{a} \right)^2 s^2 +
 2 \left( \frac{x_1}{a} \right) \left( \frac{x_2 - x_1}{a} \right) s +
 \left( \frac{x_1}{a} \right)^2
\end{align*}

The coefficients of this quadratic are functions of
\begin{equation*}
\frac{x_2 - x_1}{a}
\end{equation*}
and
\begin{equation*}
\frac{x_1}{a}
\end{equation*}

Across all three axes, we can take these components:
\begin{align*}
\vec q & = \left[ \frac{x_2 - x_1}{a},
\frac{y_2 - y_1}{b},
\frac{z_2 - z_1}{c} \right] \\
\vec p & = \left[ \frac{x_1}{a},
\frac{y_1}{b},
\frac{z_1}{c} \right]
\end{align*}

then (\ref{eqn:ellipsoid_intersect}) becomes:
\begin{equation}
( \vec q \cdot \vec q) s^2 + 2 (\vec p \cdot \vec q) s + (\vec p \cdot
\vec p - 1) = 0
\end{equation}

which can be solved as a quadratic equation for the two values of s.

If there are no roots to the equation, the line does not intersect the
ellipsoid. Otherwise, the two points of intersection are computed as:
\begin{equation}
\vec r_A = \vec r_1 + s_A( \vec r_2 - \vec r_1) \\
\vec r_B = \vec r_1 + s_B( \vec r_2 - \vec r_1)
\end{equation}


%******************************************************************************
%******************************************************************************
%******************************************************************************
%******************************************************************************
\chapter{User's Guide}

\section{Frame Transformations}
This model contains utility functions to generate transformation matrices for
desired coordinate transformations. These utilities can be used in other
models by either including \texttt{math\_utils.hh} in the header file of
the appropriate model or in the source-code file of the appropriate model
and if using Trick, as a library dependency to the same source-code file.
The code below shows an example implementations of the methods in the
\textit{MathUtils} class that can be used once the header file has
been included.

\begin{code}
double position_inrtl[3];
double velocity_inrtl[3];
// Populate position_inrtl and velocity_intrl

double T_inrtl_lvlh[3][3];
MathUtils::generate_inertial_to_lvlh( position_inrtl,
velocity_inrtl,
T_inrtl_lvlh);
\end{code}

\begin{code}
double position_inrtl[3];
double velocity_inrtl[3];
// Populate position_inrtl and velocity_intrl

double T_inrtl_uvw[3][3];
MathUtils::generate_inertial_to_uvw( position_inrtl,
velocity_inrtl,
T_inrtl_uvw);
\end{code}

\begin{code}
double position_inrtl[3];
double x_axis_inrtl[3];
// Populate position_inrtl and x_axis_intrl

double T_inrtl_reference[3][3];
MathUtils::generate_inrtl_to_reference( x_axis_inrtl,
position_inrtl,
T_inrtl_reference);
\end{code}

\begin{code}
double position_inrtl[3];
double velocity_inrtl[3];
// Populate position_inrtl and velocity_intrl

double T_inrtl_vnc[3][3];
MathUtils::generate_inertial_to_vnc( position_inrtl,
velocity_inrtl,
T_inrtl_vnc);
\end{code}

\begin{code}
double position_pcpf[3];
// Populate position_pcpf

double T_pfix_enu[3][3];
MathUtils::generate_T_pfix_to_enu( position_pcpf,
T_pfix_enu);
\end{code}

\begin{code}
// Assign longitude and latitude pair
double longitude = ...;
double latitude = ...;

jeod::Quaternion Q_enu_pfix;
MathUtils::generate_Q_enu_to_pfix( longitude,
latitude,
Q_enu_pfix);
\end{code}


%******************************************************************************
%******************************************************************************
\section{Protection Methods}

Functions such as \texttt{sqrt\_protected}, \texttt{asin\_protected},
\texttt{acos\_protected}, etc. are simply called like their non-protected counterparts.
For example:
\begin{code}
double x = 1.1;
double y = MathUtils::acos_protected(x)
\end{code}
will result in $y=0.0$.

The \texttt{divide\_protected} is slightly different because the numerator and
denominator are provided as separate arguments.Two additional optional
arguments then specify the value to output if the division should fail (default
0.0), and a flag to indicate whether the failure of this division is critical
(default false).Obviously, if the 4th argument is specified as true, then the
value of the 3rd argument is moot because execution will be terminated.
\begin{code}
double x, y = 0.0;
z = MathUtils::divide_protected( x, y, 5.0)
\end{code}
will result in $z = 5.0$ whereas
\begin{code}
double x, y = 0.0;
z = MathUtils::divide_protected( x, y, 5.0, true)
\end{code}
will terminate the execution.

The same fail-value options may also be applied to \texttt{log\_protected} and
\texttt{log10\_protected}, e.g.
\begin{code}
double x = -1.0;
y = MathUtils::log_protected(x);
\end{code}
will return 0.0, whereas:
\begin{code}
double x = -1.0;
y = MathUtils::log_protected(x, 0.1);
\end{code}
will return 0.1, and
\begin{code}
double x = -1.0;
double y = MathUtils::log_protected(x, 0.0, true);
\end{code}
will terminate the execution.

%******************************************************************************
%******************************************************************************
\section{Numerical Methods}

\subsection{is\_near\_equal}
For example, to test whether two values are equal to some unit
in last place (ulp):

\begin{code}
double x = asin_protected(2.0); // pi/2
double y = acos_protected(-2.0) // pi
if (MathUtils::is_near_equal(x*2.0, y)) { // ulp defaults to 0.5
...
}
\end{code}

\subsection{has\_changed\_from}
For example, to test whether a value has changed from a previously stored value
(e.g. like across iterations or timesteps):
\begin{code}
double prev_x = 1.0;
double curr_x = prev_x + delta_x; // some update to the value
if (MathUtils::has_changed_from(prev_x, curr_x)) {
...
}
\end{code}

\subsection{is\_equal}
For example, to test whether a value has not changed from a previously stored value
and is still equal (e.g. like across iterations or timesteps):
\begin{code}
double prev_x = 1.0;
double curr_x = prev_x; // no change in value
if (MathUtils::has_changed_from(prev_x, curr_x)) {
...
}
\end{code}

\subsection{is\_within\_abs\_tolerance}
For example, to test whether two values are within a specified absolute
tolerance of one another:
\begin{code}
double value_1 = MathUtils::acos_protected(x);
double value_2 = M_PI;
double absolute_tolerance = 0.1;
if (MathUtils::is_within_abs_tolerance(value_1,
 value_2,
 absolute_tolerance)) {
...
}
\end{code}
This will return \texttt{true} for values of x close to or less than -1, where the
arc-cosine result is close to or equal to $\pi$ by an absolute difference
of 0.1.

Note that if the absolute tolerance is negative, this method will
automatically return \texttt{false}.

\subsection{is\_within\_rel\_tolerance}
For example, to test whether a value is within a specified relative
tolerance of an expected value:
\begin{code}
double value = MathUtils::acos_protected(x);
double expected_value = M_PI;
double relative_tolerance = 0.01; // 1% tolerance of expected_value
if (MathUtils::is_within_rel_tolerance(value,
 expected_value,
 relative_tolerance)) {
...
}
\end{code}
This will return \texttt{true} for values of x close to or less than -1, where the
arc-cosine result is close to or equal to $\pi$ within a relative difference
of 0.01 (1\%).

Note:
\begin{itemize}
\item A relative tolerance compares how close two values are as a 
percetange of the expected value.
\item If the relative tolerance is negative, this method will
automatically return \texttt{false}.
\end{itemize}

\subsection{sign}
For example, to determine what sign a value has:
\begin{code}
double x = -5;
if (MathUtils::sign(x) == -1) { // returns -1 for negative values
...
}
\end{code}

\subsection{Polynomial}
The \texttt{polynomial} method requires a populated STL-vector and a value
for the variable:
\begin{code}
double x = 1.1;
std::vector<double> coeffs{1.0, 0.0, 2.1, -1.0};
double y = MathUtils::polynomial(x, coeffs);
\end{code}
evaluates $y = 1.0 + 0.0x + 2.1 x^2 -1.0 x^3$ at $x=1.1$


\subsection{Quadratic Solver}

The quadratic solver requires values for the three coefficients of the quadratic equation:
\begin{equation*}
ax^2 + bx + c = 0
\end{equation*}

A typical usage involves setting the coefficients manually and then calling the 
\texttt{update()} method:
\begin{code}
QuadraticSolver q_solve;
q_solve.a = 1.0;
q_solve.b = 0.0;
q_solve.c = -9.0;
\end{code}

Then the \texttt{QuadraticSolver::update()} method is called. The call takes
an optional argument indicating whether the roots are to be computed, or
simply identified as existing (i.e., whether they are real). This optional argument
defaults to \texttt{true}.

\texttt{QuadraticSolver::update(false)} with the settings above will simply
return \texttt{true} -- the equation $x^2 - 9 = 0$ has real roots.

\texttt{QuadraticSolver::update()} (or \texttt{QuadraticSolver::update(true)})
with the settings above will also return \texttt{true} and will populate the member variables
\texttt{root1, root2} as ${-3, +3}$ respectively.

Alternatively, the solver can be initialized directly with the coefficients using the constructor:
\begin{code}
QuadraticSolver q_solve(1.0, 0.0, -9.0);
\end{code}

This constructor takes an optional fourth argument \texttt{compute\_roots} 
(defaulting to \texttt{true}). If set to \texttt{true}, the roots are computed immediately upon
construction using the same logic as \texttt{update()}. This provides a convenient one-line
instantiation for solving quadratic equations when the coefficients are known at construction time.

%******************************************************************************
%******************************************************************************
\section{Matrix and Vector Methods}

\subsection{Cholesky Decomposition}
The \texttt{cholesky\_decomposition} method is called with 4 or 5 arguments:
\begin{code}
double A[6][6];
// Populate A

double sqrt_A[6][6];
MathUtils::cholesky_decomposition( 'identifying string',
&A[0][0],
&sqrt_A[0][0],
6);
\end{code}
or, to process only the first 3x3 part of the matrix:
\begin{code}
double A[6][6];
// Populate A

double sqrt_A[6][6];
if (MathUtils::cholesky_decomposition( 'identifying string',
 &A[0][0],
 &sqrt_A[0][0],
 6,
 3)) {
...
}
\end{code}

\subsection{Position-Velocity Covariance Matrix Transformation}
For example, to transform a position-velocity covariance from one frame to 
another using the \texttt{transform\_pv\_matrix} method:
\begin{code}
double T[3][3]; // transformation matrix from frame A to frame B
double cov_A[6][6]; // position-velocity covariance matrix in frame A
// Populate T and A

double cov_B[6][6]; // position-velocity covariance matrix in frame B
MathUtils::transform_pv_matrix(T,
 cov_A,
 cov_B);
\end{code}

\subsection{Correlation Matrix}
For example, to extract a correlation matrix from a covariance matrix
using the \\\texttt{extract\_correlation\_coefficients} method requires a
square covariance matrix that is at least populated on the lower diagonal:
\begin{code}
double cov[3][3];
// Populate cov

double corr[3][3];
if (MathUtils::extract_correlation_coefficients(cov,
corr)) {
...
}
\end{code}

\subsection{Matrix and Vector Algebra}
See Section \ref{sec:spec:matrix_methods} for signatures of these methods, e.g.:
\begin{code}
double A[6][3];
double B[6][3]
// Populate A and B

double Y[6][6];
MathUtils::matrix_mult_right_trans<dim1, dim2, dim3>( A,
B,
Y);
\end{code}

Alternative Signatures: 

For the template-based matrix methods, the compiler identifies the 
matrix sizes from context. For example:
\begin{code}
matrix_copy(const double (&A)[dim1][dim2],
double (&B)[dim1][dim2])
\end{code}
Instead of the full declaration \texttt{matrix\_mult\_right\_trans<dim1,dim2>},
which includes the template parameter specification \texttt{<dim1,dim2>}, but
these can be omitted as illustrated in the alternative signature.

\textbf{Note: the automatic dimensioning applies to any of the other
template-based matrix methods too}

\subsection {Matrix Inversion (via Cholesky Decomposition)}
This method takes two or three arguments, the third argument will default to
false if not provided.
\begin{code}
double X[9][9];
// Populate X

double X_inv[9][9];
if (MathUtils::matrix_inv_using_cholesky( X,
X_inv)) {
...
}
\end{code}

If the third argument is provided and is true:
\begin{code}
if (MathUtils::matrix_inv_using_cholesky( X,
X_inv,
true)) {
...
}
\end{code}
then a failure to obtain an inverse matrix will result in simulation
termination. Otherwise, failure will result in an error message an the
elements of \textit{X\_inv} will be assigned to zero. Failure to obtain an
inverse will result if the Cholesky decomposition fails (due to having a
negative determinant) or if any of the diagonal elements of the
decomposition are zero (due to having a zero determinant).

Of particular note, inputting a non-symmetric matrix \textbf{will not result
in a failure flag}. The Cholesky decomposition algorithm will generate an
error message indicating that the upper diagonal elements will be ignored,
then it will continue to decompose a matrix generated from the values on and
below the diagonal. The matrix produced by this method will then be
the inverse of this generated matrix, not of the original matrix.This
implementation is designed to prevent minor numerical differences in
symmetry from blocking the generation of the decomposed and inverse
matrices. \textbf{IT IS ASSUMED THAT THE INPUT MATRIX IS SYMMETRIC}.


%******************************************************************************
%******************************************************************************
\section{Calculus Methods}

\subsection{Backward Difference}
The \texttt{compute\_backward\_difference} method takes a STL-list of a 
variable's history and returns
a value representing the backward-differenced representation of the derivative.
The difference-value does not include the step size; to obtain a derivative,
divide the difference by the step size.
The STL-list is to be arranged with the most recent values to the front.

\begin{code}
double x0;
double x1;
std::list<double> x_history;
x_history.push_front( x0);
x_history.push_front( x1);
double x_diff =
MathUtils::compute_backward_difference( x_history);
double x_dot = x_diff / (t1 - t0);
\end{code}
Alternatively, to make that last line safer:
\begin{code}
double x_dot = MathUtils::divide_protected( x_diff,
t1 - t0);
\end{code}


\subsection{Unit Vector Derivative}
The \texttt{compute\_unit\_vector\_derivative} method takes 3 arguments:
\begin {enumerate}
 \item a 3-vector $\vec x$; this argument is an input
 \item a 3-vector representing the derivative of $\vec x$, $\dot{\vec x}$;
 this argument is an input
 \item a 3-vector representing the derivative of the unit-vector
 representation of $\vec x$, $\dot{ \hat x}$; this argument is an ouput.
\end {enumerate}
\begin{code}
double position[3] = ...;
double velocity[3] = ...;
double d_pos_hat[3];
MathUtils::compute_unit_vector_derivative( position,
 velocity,
 d_pos_hat);
\end{code}
For example,
\begin{code}
double position[3] = {1.0, 2.0, 3.0};
double velocity[3] = {0.1, 0.2, 0.3};
double d_pos_hat[3];
MathUtils::compute_unit_vector_derivative( position,
 velocity,
 d_pos_hat);
\end{code}
will result in $d\_pos\_hat = \{0.0, 0.0, 0.0\}$ because the velocity is
aligned with
the position vector, so the direction of the position vector does not change.


%******************************************************************************
%******************************************************************************
\section{Geometry Methods}
\subsection{Ellipsoid Intersection}
The ellipsoid intersection is used with its own instance of the 
\texttt{EllipsoidIntersection} class:
\begin{code}
#include "cml/models/utilities/math_utils/include/ellipsoid_intersection.hh"
...
EllipsoidIntersection ellipsoid_intersect;
\end{code}

The constructor for the model takes 5 arguments:
\begin{enumerate}
\item A 3-vector specifying a reference to the position vector used for one
end of the line being tested. This position vector must be in a frame
concentric with the ellipsoid.
\item A 3-vector specifying a reference to the position vector used for the
other end of the line being tested.
\item The x-axis semi-axis of the ellipsoid. Note that the order in which
the semi-axes are specified is important. This first one shall be associated
with the values for the first axis in the position vectors above.
\item The y-axis semi-axis of the ellipsoid.
\item The z-axis semi-axis of the ellipsoid.
\end{enumerate}

\begin{code}
ellipsoid_intersect ( end1,
end2,
x_sma,
y_sma,
z_sma)
\end{code}

To test whether the line between the two reference points penetrates the
ellipsoid, call
\begin{code}
bool line_intersection = ellipsoid_intersect.update(false);
\end{code}

To test whether the line between the two reference points penetrates the
ellipsoid \textit{and} compute the points of intersection, call
\begin{code}
bool line_intersection = ellipsoid_intersect.update();
\end{code}

After the update method has executed, its results can be accessed:
\begin{code}
bool line_intersection = ellipsoid_intersect.intersection;
double x_intersection_1 = ellipsoid_intersect.root1[0];
double y_intersection_1 = ellipsoid_intersect.root1[1];
double z_intersection_1 = ellipsoid_intersect.root1[2];
double x_intersection_2 = ellipsoid_intersect.root2[0];
double y_intersection_2 = ellipsoid_intersect.root2[1];
double z_intersection_2 = ellipsoid_intersect.root2[2];
\end{code}


%******************************************************************************
%******************************************************************************
%******************************************************************************
%******************************************************************************
\chapter{Inspection, Verification and Validation}
The model was inspected to ensure compliance with coding standards and
requirement satisfaction was determined through unit-simulation testing.

\section{Code Coverage}
\begin{figure}[h!]
\begin{center}
\label{fig:code_coverage}
\includegraphics[width=\linewidth]{Figs/code_coverage.png}
\end{center}
\end{figure}

\subsection{Exceptions}
The only exceptions to the code coverage are deprecated alias methods such as
\texttt{sqrt}. 

Additionally, some templated methods that are not directly executed are ignored
by the \textit{GCOVR} coverage report. These methods are either deprecated aliases
or templated variations of other methods that are already covered.

\section{Frame Transformations}

Tests evaluating frame transformations were performed in
\textit{SIM\_unit\_math\_utils}.Various scenarios based on vector inputs
were executed; the vector inputs are shown in Table \ref{tab:test} and were
applied to the following runs:
\begin{itemize}
 \item \textit{RUN\_lvlh}
 \item \textit{RUN\_uvw}
 \item \textit{RUN\_reference}
 \item \textit{RUN\_vnc}
 \item \textit{RUN\_pfix\_enu}
\end{itemize}

Additionally, a set of longitude/latitude pairs were used to evaluate the
ENU-to-Pfix Quaternion generator in \textit{RUN\_enu\_pfix}.

Results obtained through these tests were as expected and are detailed below.

The input arguments in the transformation-matrix generators are
three-element vectors for all four options. These vectors can be
labeled \textbf{a} and \textbf{b}. For the inertial to LVLH, UVW,
VNC methods, \textbf{a} and \textbf{b} represent position and velocity,
respectively. The Pfix-to-ENU method only uses one vector (\textbf{a}) that
represents position. For the inertial to reference method, \textbf{a} and
\textbf{b} represent
the x-axis in inertial space and position, respectively. Note that while
all cases were chosen to extract a specific response from the model, they
are not all relevant for each method. The results of each relevant case for
each method in the model are presented in the following sections.

\begin{table}[ht!]
\caption{Test scenarios investigated.}
\centering
\begin{tabular}{c c c c c c c}
\hline \hline
 & \multicolumn{3}{c}{\textbf{a}} & \multicolumn{3}{c}{\textbf{b}} \\
 Case & $a_x$ & $a_y$ & $a_z$ & $b_x$ & $b_y$ & $b_z$ \\ \hline
1 & 0 & 0 & 0 & 10 & 10 & 10 \\
2 & 200 & 100 & 0 & 0 & 0 & 0 \\
3 & 100 & 0 & 0 & 0 & 0 & 0 \\
4 & 0 & 0 & 200 & 0 & 0 & 0 \\
5 & 200 & 100 & 300 & 10 & 30 & 20 \\
6 & 200 & 100 & 300 & 200 & 100 & 300 \\
7 & 0 & 0 & 10 & 0 & 0 & 10 \\ \hline \hline
\end{tabular}
\label{tab:test}
\end{table}

Since the frame transformation functions for inertial-to-LVLH, UVW, reference, 
Pfix-to-ENU, and ENU-to-Pfix accept C-style arrays as values rather than 
references (unlike the inertial-to-VNC transformation), any null pointer 
passed to these functions should trigger an error and terminate the operation. 
The function \textit{RUN\_transform\_NULL} is designed to test each parameter with null 
pointers to ensure that such cases are properly detected and handled.

\subsection{Results for Inertial to LVLH Transformation}
The expected results of calling \texttt{generate\_inertial\_to\_lvlh} for the
specified cases are discussed below. Note that \textbf{a} and \textbf{b} from
Table \ref{tab:test} represent the position and velocity vectors, respectively.
\begin{enumerate}
 \item $\textbf{r} = [0 \quad 0 \quad 0]^T, \qquad
\textbf{v} = [10 \quad 10 \quad 10]^T$ \\
 The magnitude of the position vector is zero, so an error message
 indicating the same is expected. The LVLH frame will be aligned with
 the inertial frame and the resulting transformation matrix will be
 the identity matrix.

 \item $\textbf{r} = [200 \quad 100 \quad 0]^T, \qquad
\textbf{v} = [0 \quad 0 \quad 0]^T$ \\
 The magnitude of the velocity vector is zero and the position vector
 does not lie along the inertial x-axis. The LVLH x-axis will be
 aligned close to the inertial x-axis and a warning indicating the
 same is expected. The resulting transformation matrix, found using
 (\ref{eq:lvlh_nox}) and (\ref{eq:lvlh_mat}), is
 \begin{equation*}
\hfill R = \begin{bmatrix} 0.4472 & -0.8944 & 0 \\
 0&0& 1 \\
-0.8944 & -0.4472 & 0
 \end{bmatrix} \hfill
 \end{equation*}

 \item $\textbf{r} = [100 \quad 0 \quad 0]^T, \qquad
\textbf{v} = [0 \quad 0 \quad 0]^T$ \\
 The magnitude of the velocity vector is zero and the position vector
 lies along the inertial x-axis. The LVLH y-axis will be aligned close
 to the inertial y-axis and a warning indicating the same is expected.
 The resulting transformation matrix, found using (\ref{eq:lvlh_x})
 and (\ref{eq:lvlh_mat}), is
 \begin{equation*}
 \hfill R = \begin{bmatrix} 0 & 0 & 1 \\
0 & 1 & 0 \\
-1 & 0 & 0
\end{bmatrix} \hfill
 \end{equation*}

 \item $\textbf{r} = [0 \quad 0 \quad 200]^T, \qquad
\textbf{v} = [0 \quad 0 \quad 0]^T$ \\
 The magnitude of the velocity vector is zero and the position vector
 does not lie along the inertial x-axis. The LVLH x-axis will be
 aligned close to the inertial x-axis and a warning indicating the
 same is expected. The resulting transformation matrix, found using
 (\ref{eq:lvlh_nox}) and (\ref{eq:lvlh_mat}), is
 \begin{equation*}
\hfill R = \begin{bmatrix} 1 & 0 & 0 \\
 0 & -1 & 0 \\
 0 & 0 & -1
 \end{bmatrix} \hfill
 \end{equation*}

 \item $\textbf{r} = [200 \quad 100 \quad 300]^T, \qquad
\textbf{v} = [10\quad 30\quad 20]^T$ \\
 The position and velocity vectors satisfy the assumptions that
 the position is nonzero and the velocity is neither radial nor nonzero.
 The resulting transformation matrix, found using (\ref{eq:lvlh_main})
 and (\ref{eq:lvlh_mat}), is
 \begin{equation*}
\hfill R = \begin{bmatrix} -0.2469 & 0.9567 & -0.1543 \\
0.8083 & 0.1155 & -0.5774 \\
 -0.5345 & -0.2673 & -0.8018
 \end{bmatrix} \hfill
 \end{equation*}
\end{enumerate}


\subsection{Results for Inertial to UVW Transformation}
The expected results of calling \texttt{generate\_inertial\_to\_uvw} for
the specified cases are discussed below. Note that $\textbf{a}$ and
$\textbf{b}$ from Table \ref{tab:test} represent the position and
velocity vectors, respectively.

\begin{enumerate}
 \item $\textbf{r} = [0 \quad 0 \quad 0]^T, \qquad
 \textbf{v} = [10 \quad 10 \quad 10]^T$ \\
 The magnitude of the position vector is zero, so an error message
 indicating the same is expected. The UVW frame will be aligned with
 the inertial frame and the resulting transformation matrix will be
 the identity matrix.

 \item $\textbf{r} = [200 \quad 100 \quad 0]^T, \qquad
 \textbf{v} = [0 \quad 0 \quad 0]^T$ \\
 The magnitude of the velocity vector is zero and the position vector
 does not lie along the inertial z-axis. The UVW w-axis will be
 aligned close to the inertial z-axis and a warning indicating the
 same is expected. The resulting transformation matrix, found using
 (\ref{eq:uvw_noz}) and (\ref{eq:uvw_mat}), is
 \begin{equation*}
\hfill R = \begin{bmatrix} 0.8944 & 0.4472 & 0 \\
-0.4472 & 0.8944 & 0 \\
 0& 0& 1
 \end{bmatrix} \hfill
 \end{equation*}

 \item $\textbf{r} = [100 \quad 0 \quad 0]^T, \qquad
\textbf{v} = [0 \quad 0 \quad 0]^T$ \\
 The magnitude of the velocity vector is zero and the position
 vector does not lie along the inertial z-axis. The UVW w-axis
 will be aligned close to the inertial z-axis and a warning
 indicating the same is expected. The resulting transformation matrix,
 found using (\ref{eq:uvw_noz}) and (\ref{eq:uvw_mat}), is the
 identity matrix.

 \item $\textbf{r} = [0 \quad 0 \quad 200]^T, \qquad
\textbf{v} = [0 \quad 0 \quad 0]^T$ \\
 The magnitude of the velocity vector is zero and the position vector
 lies along the inertial z-axis. The UVW w-axis will be aligned close
 to the inertial x-axis and a warning indicating the same is expected.
 The resulting transformation matrix, found using (\ref{eq:uvw_z})
 and (\ref{eq:uvw_mat}), is
 \begin{equation*}
\hfill R = \begin{bmatrix} 0 &0 & 1 \\
 0 & -1 & 0 \\
 1 &0 & 0
 \end{bmatrix} \hfill
 \end{equation*}

\item $\textbf{r} = [200 \quad 100 \quad 300]^T, \qquad
 \textbf{v} = [10 \quad 30 \quad 20]^T$ \\
 The position and velocity vectors satisfy the assumptions that the
 position is nonzero and the velocity is neither radial nor nonzero.
 The resulting transformation matrix, found using
 (\ref{eq:uvw_main}) and (\ref{eq:uvw_mat}), is
 \begin{equation*}
\hfill R = \begin{bmatrix} 0.5345 & 0.2673 & 0.8018 \\
-0.2469 & 0.9567 & -0.1543 \\
-0.8083 & -0.1155 & 0.5774
 \end{bmatrix} \hfill
 \end{equation*}
\end{enumerate}

\subsection{Results for Inertial to Reference Transformation}
The expected results of calling \texttt{generate\_inrtl\_reference} for
the specified cases are discussed below. Note that the \textbf{a} and
\textbf{b} from Table \ref{tab:test} represent the direction of the
user-specified reference-frame's x-axis and the location of its origin,
respectively.
\begin{enumerate}
 \item $\textbf{x} = [0 \quad 0 \quad 0]^T, \qquad
\textbf{r} = [10 \quad 10 \quad 10]^T$ \\
 The magnitude of the x-axis vector is zero, so an error message
 indicating the same is expected. The reference-frame will be aligned
 with the inertial frame and the resulting transformation matrix will
 be the identity matrix.a

 \setcounter{enumi}{3}
 \item $\textbf{x} = [0 \quad 0 \quad 200]^T, \qquad
\textbf{r} = [0 \quad 0 \quad 0]^T$ \\
 The magnitude of the position vector is zero, so an error message
 indicating the same is expected. The reference-frame will be
 aligned with the inertial frame and the resulting transformation
 matrix will be the identity matrix.

 \item $\textbf{x} = [200 \quad 100 \quad 300]^T, \qquad
\textbf{r} = [10 \quad 30 \quad 20]^T$ \\
 The x-axis and position vectors satisfy the assumptions and the
 resulting transformation matrix found using (\ref{eq:ref_main})
 and (\ref{eq:ref_mat}) is
 \begin{equation*}
\hfill R = \begin{bmatrix} 0.5345 & 0.2673 & 0.8018 \\
-0.8083 & -0.1155 & 0.5774 \\
 0.2469 & -0.9567 & 0.1543
 \end{bmatrix} \hfill
 \end{equation*}

\item $\textbf{x} = [200 \quad 100 \quad 300]^T, \qquad
 \textbf{r} = [200 \quad 100 \quad 300]^T$ \\
 The x-axis and position vectors are aligned and the position vector
 is not aligned with inertial z-axis. The reference-frame z-axis will
 be aligned with the inertial z-axis and the resulting transformation
 matrix found using (\ref{eq:ref_main}) and (\ref{eq:ref_noz}) is
 \begin{equation*}
\hfill R = \begin{bmatrix} 0.5345 & 0.2673& 0.8018 \\
 0.4472 & -0.8944 & 0 \\
 0.7171 & 0.3586& -0.5976
 \end{bmatrix} \hfill
 \end{equation*}

\item $\textbf{x} = [0 \quad 0 \quad 10]^T, \qquad
 \textbf{r} = [0 \quad 0 \quad 10]^T$ \\
 The x-axis and position vectors are aligned and the position vector
 lies along the inertial z-axis. The reference-frame y-axis will be
 aligned with the inertial y-axis and the resulting transformation
 matrix found using (\ref{eq:ref_main}) and (\ref{eq:ref_z}) is
 \begin{equation*}
\hfill R = \begin{bmatrix} 0 & 0 & 1 \\
 0 & 1 & 0 \\
 -1 & 0 & 0
 \end{bmatrix} \hfill
 \end{equation*}
\end{enumerate}

\subsection{Results for Inertial to VNC Transformation}
The expected results of calling \texttt{generate\_inertial\_to\_vnc} for the
specified cases are discussed below. Note that \textbf{a} and \textbf{b} from
Table \ref{tab:test} represent the position and velocity vectors, respectively.
\begin{enumerate}
 \item $\textbf{r} = [0 \quad 0 \quad 0]^T, \qquad
\textbf{v} = [10 \quad 10 \quad 10]^T$ \\
 The magnitude of the position vector is zero, so an error message
 indicating the same is expected. The VNC frame will be aligned with
 the inertial frame and the resulting transformation matrix will be
 the identity matrix.

 \item $\textbf{r} = [200 \quad 100 \quad 0]^T, \qquad
\textbf{v} = [0 \quad 0 \quad 0]^T$ \\
 The magnitude of the velocity vector is zero, so an error message
 indicating the same is expected. The VNC frame will be aligned with
 the inertial frame and the resulting transformation matrix will be
 the identity matrix.

\setcounter{enumi}{4}
 \item $\textbf{r} = [200 \quad 100 \quad 300]^T, \qquad
\textbf{v} = [10\quad 30\quad 20]^T$ \\
 The position and velocity vectors satisfy the assumptions that
 the position is nonzero and the velocity is neither radial nor nonzero.
 The resulting transformation matrix, found using (\ref{eq:vnc_main})
 and (\ref{eq:vnc_mat}), is
 \begin{equation*}
\hfill R = \begin{bmatrix}0.2673 &0.8018 &0.5345\\
 -0.8083 & -0.1155 &0.5774 \\
0.5246 & -0.5864 &0.6172 
 \end{bmatrix} \hfill
 \end{equation*}

\item $\textbf{r} = [200 \quad 100 \quad 300]^T, \qquad
\textbf{v} = [200 \quad 100 \quad 300]^T$ \\
The position and velocity vectors satisfy the assumption that
both vectors are nonzero. However, the position and velocity vectors 
are aligned, as the velocity is radial. The VNC frame 
will be aligned with the inertial frame and the resulting transformation
matrix will be the identity matrix.
\end{enumerate}

\subsection{Results for Pfix to ENU Transformation}
The expected results of calling \texttt{generate\_pfix\_enu} for
the specified cases are discussed below. Note that the \textbf{a}
vector from Table \ref{tab:test} represents the position of the origin
of the ENU frame; the \textbf{b} vector is not used in this method.
\begin{enumerate}
 \item $\textbf{x} = [0 \quad 0 \quad 0]^T$ \\
 With a 0 position vector, the matrix output will be the zero-matrix.

 \item $\textbf{r} = [200 \quad 100 \quad 0]^T$ \\
 With the position vector being in the equatorial plane, the
 \textit{north}-axis will be aligned with the Pfix z-axis.
 The position unit-vector and \textit{up}-axis are
 \begin{equation*}
 \hat u = \hat r = \begin{bmatrix} \frac{2}{\sqrt{5}} \\
 \frac{1}{\sqrt{5}} \\
 0
 \end{bmatrix} \hfill
 \end{equation*}
 and the \textit{east} axis completes.
 \begin{equation*}
\hfill T = \begin{bmatrix}
 \text{[east]}^T \\
 \text{[north]}^T \\
 \text{[up]}^T
 \end{bmatrix}
 =
 \begin{bmatrix}
 -\frac{1}{\sqrt5} & \frac{2}{\sqrt5} & 0 \\
 0 &0 & 1 \\
\frac{2}{\sqrt5} & \frac{1}{\sqrt5} & 0
 \end{bmatrix}
 =
 \begin{bmatrix}
 -0.4472 & 0.8944 & 0 \\
 0 & 0& 1 \\
0.8944 & 0.4472 & 0
 \end{bmatrix} \hfill
 \end{equation*}

 \item $\textbf{r} = [100 \quad 0 \quad 0]^T$ \\
 As with the previous case, the position vector being in the
 equatorial plane means the
 \textit{north}-axis will be aligned with the Pfix z-axis.
 In this case,
 the position unit-vector and \textit{up}-axis are
 aligned with the Pfix x-axis.
 \begin{equation*}
\hfill T = \begin{bmatrix}
0 & 1 & 0 \\
0 & 0 & 1 \\
1 & 0 & 0
 \end{bmatrix} \hfill
 \end{equation*}

 \item $\textbf{r} = [0 \quad 0 \quad 200]^T$ \\
 With the position vector over the North-pole,
 \textit{up} is aligned with the z-axis,
 \textit{east} is undefined and is forcibly aligned with the x-axis,
 and \textit{north} completes, therefore aligned with the y-axis. The
 transformation matrix is identity.

 \item $\textbf{r} = [200 \quad 100 \quad 300]^T$ \\
 With this position,
 \textit{up} is oriented along the $[2~~1 ~~ 3]$ vector,
 \textit{east} (which always has a zero z-component) is parallel
 to $[-1~~ 2 ~~0]$,
 and \textit{north} completes.
 \begin{equation*}
\hfill T =
\begin{bmatrix}
 -\frac{1}{\sqrt5}&\frac{2}{\sqrt5} & 0 \\
 -\frac{6}{\sqrt{70}} & -\frac{3}{\sqrt{70}} & \frac{5}{\sqrt{70}} \\
\frac{2}{\sqrt{14}} &\frac{1}{\sqrt{14}} & \frac{3}{\sqrt{14}}
\end{bmatrix}
=
\begin{bmatrix}
-0.4472 & 0.8944 & 0 \\
-0.7171 &-0.3586 & 0.5976 \\
 0.5345 & 0.2673 & 0.8018
\end{bmatrix} \hfill
 \end{equation*}
\end{enumerate}

\subsection{Results for ENU-to-Pfix Quaternion}
To verify the results of the quaternion-generator, we use a two-level
approach with an independent spreadsheet for a selection of
latitude-longitude pairs:
\begin{itemize}
\item the interim quaternions $Q_{LT, S1 \rightarrow S2}$ and
$Q_{LT, S2 \rightarrow UEN}$ are generated for inspection, and
the overall quaternion $Q_{LT, Pfix \rightarrow ENU}$ generated for
comparison with the result of the model.
\item The resulting quaternion is applied to the basis vectors
representing \textit{east}, \textit{north}, and \textit{up} in the
ENU frame to create the representations of those vectors in the Pfix
frame. These representations are evaluated using geometry to confirm
that, for the given longitude and latitude, the resulting east,
north, and up directions are as expected.
\end {itemize}

In the figure below, a selection of the test points are
illlustrated with the east, north, and up (blue, green, red respectively)
axes shown for each illustrated test point.
\begin{figure}[h!]
\begin{center}
\label{fig:quat_POI}
\includegraphics[scale=0.5]{Figs/quat_POI.png}
\caption{Selection of the Test Points for Verification of Q\_enu\_pfix}
\end{center}
\end{figure}

\newpage
\begin{enumerate}
\item \textbf{latitude=0, longitude=0} The point of interest lies on the
Pfix x-axis: \\
\textit{east} is directed along $[0~,~1~,~0]$ \\
\textit{north} is directed along $[0~,~0~,~1]$ \\
\textit{up} is directed along $[1~,~0~,~0]$ \\
The resulting quaternion is $
\begin{bmatrix} \frac{1}{2} \\
\begin{bmatrix} \frac{1}{2} \\
\frac{1}{2} \\
\frac{1}{2}
\end{bmatrix}
\end{bmatrix}$
\begin{figure}[h!]
\begin{center}
\includegraphics[scale=0.45]{Figs/quat_case1.png}
\end{center}
\end{figure}

\newpage
\item \textbf{latitude=0, longitude=90} The point of interest lies on the
Pfix y-axis: \\
\textit{east} is directed along $[-1~,~0~,~0]$ \\
\textit{north} is directed along $[0~,~0~,~1]$ \\
\textit{up} is directed along $[0~,~1~,~0]$ \\
The resulting quaternion is$
\begin{bmatrix} 0 \\
\begin{bmatrix} 0 \\
\frac{1}{\sqrt{2}} \\
\frac{1}{\sqrt{2}}
\end{bmatrix}
\end{bmatrix}$
\begin{figure}[h!]
\begin{center}
\includegraphics[scale=0.45]{Figs/quat_case2.png}
\end{center}
\end{figure}

\newpage
\item \textbf{latitude=89.99999, longitude=0} The point of interest lies
effectively on the Pfix z-axis with a position slightly offset
along the +x-axis from the z-axis; this avoids the ambiguity of being
exactly on the axis.\\
\textit{east} is directed along $[0~,~1~,~0]$ \\
\textit{north} is directed along $[-1~,~0~,~0]$ \\
\textit{up} is directed along $[0~,~0~,~1]$ \\
The resulting quaternion is $
\begin{bmatrix} \frac{1}{\sqrt{2}} \\
\begin{bmatrix} 0 \\
0 \\
\frac{1}{\sqrt{2}}
\end{bmatrix}
\end{bmatrix}$
\begin{figure}[h!]
\begin{center}
\includegraphics[scale=0.45]{Figs/quat_case3.png}
\end{center}
\end{figure}

\newpage
\item \textbf{latitude=89.99999, longitude=90} The point of interest lies
effectively on the Pfix z-axis with a position slightly offset
along the +y-axis from the z-axis, again to avoid the ambiguity of
being exactly on the axis. This test illustrates the ambiguity, with
a very small ``spatial displacement'' around the pole, the ENU frame
changes substantially as north now points along the -y-axis instead
of the -x-axis.\\
\textit{east} is directed along $[-1~,~0~,~0]$ \\
\textit{north} is directed along $[0~,~-1~,~0]$ \\
\textit{up} is directed along $[0~,~0~,~1]$ \\
The resulting quaternion is $
\begin{bmatrix} 0 \\
\begin{bmatrix} 0 \\
0 \\
1
\end{bmatrix}
\end{bmatrix}$
\begin{figure}[h!]
\begin{center}
\includegraphics[scale=0.45]{Figs/quat_case4.png}
\end{center}
\end{figure}

\newpage
\item \textbf{latitude=90.00001, longitude=90} The point of interest lies
effectively on the Pfix z-axis with a position slightly offset
from the z-axis, again to avoid the ambiguity of being
exactly on the axis. In this test case, we evaluate the model's
internal handling of latitudes greater than 90 degrees. Note that
while this is non-physical, a polar (or near-polar) trajectory may
well produce a latitude greater than 90 degrees due to numerical
rounding as it crosses the pole.
The model should recognize this value and interpret it as
a latitude of a little less than 90 degreees with
longitude offset by 180 degrees, effectively producing a latitude
of 89.99999 degrees and longitude of 270 degrees. \\
\textit{east} is directed along $[1~,~0~,~0]$ \\
\textit{north} is directed along $[0~,~1~,~0]$ \\
\textit{up} is directed along $[0~,~0~,~1]$ \\
The resulting quaternion is the identity:$
\begin{bmatrix} 1 \\
\begin{bmatrix} 0 \\
0 \\
0
\end{bmatrix}
\end{bmatrix}$ \\
Note that the model actually produces the negative of these values,
but recognizing that a quaternion is operationally equivalent to its
negative, we confirm that these are the same values.
\begin{figure}[h!]
\begin{center}
\includegraphics[scale=0.45]{Figs/quat_case5.png}
\end{center}
\end{figure}

\newpage
\item \textbf{latitude=-90.00001, longitude=90} This test investigates the
model behavior with out-of-domain latitude values at the South pole.
The model should correct these lat/lon values to -89.99999 degrees
latitude, 270 degrees longitude. \\
\textit{east} is directed along $[1~,~0~,~0]$ \\
\textit{north} is directed along $[0~,~-1~,~0]$ \\
\textit{up} is directed along $[0~,~0~,~-1]$ \\
The resulting quaternion is$
\begin{bmatrix} 0 \\
\begin{bmatrix} -1 \\
0 \\
0
\end{bmatrix}
\end{bmatrix}$
\begin{figure}[h!]
\begin{center}
\includegraphics[scale=0.45]{Figs/quat_case6.png}
\end{center}
\end{figure}

For the remaining cases, we will not be including the spreadsheet
data in the document. However, the spreadsheet is committed at
\textit{Figs/Q\_enu\_pfix.ods} and has been used to verify that the
model outputs match with the spreadsheet prediction, and that the
spreadsheet prediction produces the expected geometry.

\item \textbf{latitude=0, longitude=30} This test starts a set with a
position located off the basis vectors.
\textit{east} is$[-\frac{1}{2}~,~\frac{\sqrt{3}}{2}~,~0]$ \\
\textit{north} is $[0~,~0~,~1]$ \\
\textit{up} is$[\frac{\sqrt{3}}{2}~,~\frac{1}{2}~,~0]$ \\
The resulting quaternion is$
\begin{bmatrix} 0.35 \\
\begin{bmatrix} 0.35 \\
0.61 \\
0.61
\end{bmatrix}
\end{bmatrix}$
\item \textbf{latitude=30, longitude=0} \\
\textit{east} is$[0~,~1~,~0]$ \\
\textit{north} is $[-\frac{1}{2}~,~0~,~\frac{\sqrt{3}}{2}]$ \\
\textit{up} is$[\frac{\sqrt3}{2}~,~0~,~\frac{1}{2}]$ \\
The resulting quaternion is$
\begin{bmatrix} 0.61 \\
\begin{bmatrix} 0.35 \\
0.35 \\
0.61
\end{bmatrix}
\end{bmatrix}$

\item \textbf{latitude=30, longitude=30}\\
\textit{east} is$[-\frac{1}{2}~,~\frac{\sqrt{3}}{2}~,~0]$ \\
\textit{north} is $[0~,~0~,~1]$ \\
\textit{north} is $[-\frac{\sqrt3}{4}~,~-\frac{1}{4}~,
 ~\frac{\sqrt3}{2}]$ \\
\textit{up} is$[\frac{3}{4}~,~\frac{\sqrt{3}}{2}~,~\frac{1}{2}]$ \\
The resulting quaternion is$
\begin{bmatrix} 0.43 \\
\begin{bmatrix} 0.25 \\
0.43 \\
0.75
\end{bmatrix}
\end{bmatrix}$

\item \textbf{latitude=0, longitude=45}\\
\textit{east} is$[-\frac{1}{\sqrt2}~,~\frac{1}{\sqrt2}~,~0]$ \\
\textit{north} is $[0~,~0~,~1]$ \\
\textit{up} is$[\frac{1}{\sqrt2}~,~\frac{1}{\sqrt2}~,~0]$ \\
The resulting quaternion is$
\begin{bmatrix} 0.27 \\
\begin{bmatrix} 0.27 \\
0.65 \\
0.65
\end{bmatrix}
\end{bmatrix}$

\item \textbf{latitude=45, longitude=0}\\
\textit{east} is$[0~,~1~,~0]$ \\
\textit{north} is$[-\frac{1}{\sqrt2}~,~0~,~\frac{1}{\sqrt2}]$ \\
\textit{up} is$[\frac{1}{\sqrt2}~,~0~,~\frac{1}{\sqrt2}]$ \\
The resulting quaternion is$
\begin{bmatrix} 0.65 \\
\begin{bmatrix} 0.27 \\
0.27 \\
0.65
\end{bmatrix}
\end{bmatrix}$

\item \textbf{latitude=0, longitude=60}\\
\textit{east} is$[-\frac{\sqrt3}{2}~,~\frac{1}{2}~,~0]$ \\
\textit{north} is $[0~,~0~,~1]$ \\
\textit{up} is$[\frac{1}{2}~,~\frac{\sqrt3}{2}~,~0]$ \\
The resulting quaternion is$
\begin{bmatrix} 0.18 \\
\begin{bmatrix} 0.18 \\
0.68 \\
0.68
\end{bmatrix}
\end{bmatrix}$

\item \textbf{latitude=60, longitude=0}\\
\textit{east} is$[0~,~1~,~0]$ \\
\textit{north} is$[-\frac{\sqrt3}{2}~,~0~,~\frac{1}{2}]$ \\
\textit{up} is$[\frac{1}{2}~,~0~,~\frac{sqrt3}{2}]$ \\
The resulting quaternion is$
\begin{bmatrix} 0.68 \\
\begin{bmatrix} 0.18 \\
0.18 \\
0.68
\end{bmatrix}
\end{bmatrix}$

\item \textbf{latitude=-60, longitude=0} This test confirms the symmetry
with the southern hemisphere.\\
\textit{east} is$[0~,~1~,~0]$ \\
\textit{north} is$[\frac{\sqrt3}{2}~,~0~,~\frac{1}{2}]$ \\
\textit{up} is$[\frac{1}{2}~,~0~,~-\frac{sqrt3}{2}]$ \\
The resulting quaternion is$
\begin{bmatrix} 0.18 \\
\begin{bmatrix} 0.68 \\
0.68 \\
0.18
\end{bmatrix}
\end{bmatrix}$

\item \textbf{latitude=-30, longitude=120} \\
\textit{east} is$[-\frac{\sqrt3}{2}~,~-\frac{1}{2}~,~0]$ \\
\textit{north} is $[-\frac{1}{4}~,~\frac{\sqrt3}{4}~,
~\frac{\sqrt3}{2}]$ \\
\textit{up} is$[-\frac{\sqrt3}{4}~,~\frac{3}{4}~,~-\frac{1}{2}]$ \\
The resulting quaternion is$
\begin{bmatrix} -0.13 \\
\begin{bmatrix} -0.22 \\
0.84 \\
0.48
\end{bmatrix}
\end{bmatrix}$

\item \textbf{latitude=0, longitude=240} This test starts a set of 3 tests
with identical positions specified by different longitude
numerical representations.\\
\textit{east} is$[\frac{\sqrt3}{2}~,~-\frac{1}{2}~,~0]$ \\
\textit{north} is $[0~,~0~,~1]$ \\
\textit{up} is$[-\frac{1}{2}~,~-\frac{\sqrt3}{2}~,~0]$ \\
The resulting quaternion is$
\begin{bmatrix} -0.68 \\
\begin{bmatrix} -0.68 \\
0.18 \\
0.18
\end{bmatrix}
\end{bmatrix}$

\item \textbf{latitude=0, longitude=-120}
This test produces the same results as the previous test.

\item \textbf{latitude=0, longitude=600}
This test produces the same results as the previous test.

\end{enumerate}



\subsection{Compound Transformations}
The methods in this model can also be used to perform compound
transformations to facilitate conversions between the LVLH, UVW and
user-defined reference-frame.These transformations will be demonstrated
by an example using the $\textbf{x}$ and $\textbf{r}$ vectors defined in
the common Case 5, shared in the verifications of the individual
transformations detailed above; this case satisfies the requirements
for nominal operations for all three inertial-to-designated-frame
transformations.

\subsubsection{LVLH to UVW}
In order to transform vectors in the LVLH frame to the UVW frame, the vectors in
the LVLH frame can first be converted to the inertial frame through the method
\texttt{generate\_inertial\_to\_lvlh}. This can be followed by the method
\texttt{generate\_inertial\_to\_uvw} to convert from inertial to UVW.
For example,
if $\textbf{r} = [200 \quad 100 \quad 300]^T$ and
$\textbf{v} = [10 \quad 30 \quad 20]^T$,
then the matrices
\begin{equation*}
 \hfill R_{I\rightarrow lvlh} = \begin{bmatrix} -0.2469 & 0.9567 & -0.1543 \\
 0.8083 & 0.1155 & -0.5774 \\
-0.5345 & -0.2673 & -0.8018
\end{bmatrix} \hfill
\end{equation*}
and
\begin{equation*}
 \hfill R_{I\rightarrow uvw} = \begin{bmatrix} 0.5345 & 0.2673 & 0.8018 \\
-0.2469 & 0.9567 & -0.1543 \\
-0.8083 & -0.1155 & 0.5774
 \end{bmatrix} \hfill
\end{equation*}
can used to transform vectors in the inertial frame to the LVLH and UVW frames,
respectively. Note that these are the same matrices shown for Case 5 in the
respective previous sections. The transformation matrix
\begin{align*}
 \hfill R_{lvlh\rightarrow uvw} &=
R_{I\rightarrow uvw} R_{lvlh\rightarrow I} =
R_{I\rightarrow uvw} R_{I\rightarrow lvlh}^T \hfill \\
 \hfill &= \begin{bmatrix} 0.5345 & 0.2673 & 0.8018 \\
-0.2469 & 0.9567 & -0.1543 \\
-0.8083 & -0.1155 & 0.5774
 \end{bmatrix}
 \begin{bmatrix} -0.2469 & 0.9567 & -0.1543 \\
0.8083 & 0.1155 & -0.5774 \\
 -0.5345 & -0.2673 & -0.8018
 \end{bmatrix}^T\hfill \\
 \hfill &= \begin{bmatrix} 0 & 0 & -1 \\
 1 & 0 & 0 \\
 0 & -1 & 0
 \end{bmatrix} \hfill
\end{align*}
can then be used to convert vectors in the LVLH frame to the UVW frame.

\section{Protection Methods}
Verification of these methods is found in runs within
\textit{SIM\_unit\_math\_utils\_private}, utilizing the following tests:
\begin{itemize}
 \item ERROR\_sqrt\_protected
 \item FAIL\_div\_protected
 \item FAIL\_log\_protected
 \item FAIL\_log10\_protected
 \item RUN\_acos\_protected
 \item RUN\_asin\_protected
 \item RUN\_div\_protected
 \item RUN\_log\_protected
 \item RUN\_log10\_protected
 \item RUN\_sqrt\_protected
\end{itemize}

Each run executes the respective protection method, and the results are
compared against the expected result using the
\texttt{MathUtils::is\_near\_equal(...)} utility.

\subsection{Division-by-zero}

\subsubsection{RUN\_div\_protected}
The table below lists the tests used to evaluate the response of
\texttt{MathUtils::divide\_protected}.
\begin{center}
\begin{tabular}{|l||l|l|l|l|}
\hline
 \textbf{Data Type} & \textbf{Dividend} & \textbf{Divisor} &
 \textbf{Recovery} & \textbf {Expected Value} \\
\hline \hline
double & 15.0 & 5.0& None (0.0) & 3.0 \\
double &1.0 & 0.0& 70.0 & 70.0
\textsuperscript{*} \\
double & $1 \times 10^{-300}$ & $1 \times 10^{300}$& 2.0& 0.0 \\
double & $1 \times 10^{38}$ & $1 \times 10^{-38}$ & 35.0 & $1 \times 10^{76}$\\
double & $1 \times 10^{300}$& $1 \times 10^{-300}$ & 35.0 & 35.0
\textsuperscript{*} \\
\hline
float& 15.0 & 5.0& None (0.0) & 3.0 \\
float&1.0 & 0.0& 3.14 & 3.14
\textsuperscript{*}\\
float& $1 \times 10^{-38}$& $1 \times 10^{38}$ & 2.0& 0.0 \\
float& $1 \times 10^{38}$ & $1 \times 10^{-38}$& 5.45678& 5.45678
\textsuperscript{*}\\
\hline \hline
\end{tabular}
\end{center}

\textsuperscript{*} These runs produce the expected warning regarding div-0 or
overflow.

\begin{figure}[H]
 \begin{center}
 \includegraphics[scale=0.7]{Figs/div_protect.png}
 \caption{Screen Output from RUN\_div\_protected}
 \end{center}
\end{figure}

Note that when using type \texttt{double} the operation with $10^{38}$
succeeds, but it fails with an overflow when using data-type \texttt{float}.

\clearpage


\subsubsection{FAIL\_div\_protected}
This run executes \texttt{MathUtils::divide\_protected} with a divide-by-zero
operation and the termination flag set to true. The simulation is terminated as
expected.


\subsection{Arc-cosine}
The table below lists the tests in \textit{RUN\_acos\_protected} used to evaluate the response of \\
\texttt{MathUtils::acos\_protected}.
\begin{center}
\begin{tabular}{|l||l|l|}
\hline
 \textbf{Data Type} & \textbf{Input} & \textbf{Expected Value} \\
\hline \hline
double &0.877582561890373 & 0.5 \textsuperscript{*} \\
double & -1.1 & $\pi$ \\
double & +1.1 & 0.0 \\
\hline
float&0.877582561890373 & 0.5 \\
float& -1.1 & $\pi$ \\
float& +1.1 & 0.0 \\
\hline \hline
\end{tabular}
\end{center}
\textsuperscript{*} Because of the sensitivity of the function to the input
values, this comparaison agrees only to within 10 ULP when using data-type
\texttt{double}.


\subsection{Arc-sine}
The table below lists the tests in \textit{RUN\_asin\_protected} used to evaluate the response of \\
\texttt{MathUtils::asin\_protected}.
\begin{center}
\begin{tabular}{|l||l|l|}
\hline
 \textbf{Data Type} & \textbf{Input} & \textbf{Expected Value} \\
\hline \hline
double &0.479425538604203 & 0.5 \\
double & -1.1 & -$\frac{\pi}{2}$ \\
double & +1.1 &$\frac{\pi}{2}$ \\
\hline
float&0.877582561890373 & 0.5 \\
float& -1.1 & -$\frac{\pi}{2}$ \\
float& +1.1 &$\frac{\pi}{2}$ \\
\hline \hline
\end{tabular}
\end{center}

Results match with expectation.

\subsection{Natural Logarithm}
\subsubsection{RUN\_log\_protected}
The table below lists the tests used to evaluate the response of
\texttt{MathUtils::log\_protected}
\begin{center}
\begin{tabular}{|l||l|l|}
\hline
 \textbf{Data Type} & \textbf{Input} & \textbf{Expected Value} \\
\hline \hline
double &1.0 & 0.0 \\
double & $e$& 1.0 \\
double & $1 \times 10^{-308}$ & -709.1962086421661 \\
double & $1 \times 10^{308}$&709.1962086421661 \\
\hline
float&1.0 & 0.0 \\
float& $e$& 1.0 \textsuperscript{*}\\
float& $e$& 0.99999994 \\
float& $1 \times 10^{-38}$ & -87.49823 \\
float& $1 \times 10^{38}$&87.49823 \\
\hline \hline
\end{tabular}
\end{center}

\textsuperscript{*} When using \texttt{float}, the natural log of $e$ is 1.0 to
within 1 ULP, but not within the default 0.5 ULP.

\subsubsection{FAIL\_log\_protected}
This case tests the application of the \texttt{MathUtils::log\_protected} 
to an invalid input
(0.0), with the terminate-on-failure flag set to true. The simulation
terminates.


\subsection{Base-10 Logarithm}
\subsubsection{RUN\_log10\_protected}
The table below lists the tests used to evaluate the response of
\texttt{MathUtils::log10\_protected}
\begin{center}
\begin{tabular}{|l||l|l|}
\hline
 \textbf{Data Type} & \textbf{Input} & \textbf{Expected Value} \\
\hline \hline
double &1.0 & 0.0 \\
double & 10.0 & 1.0 \\
double & $1 \times 10^{-308}$ & -308 \\
double & $1 \times 10^{308}$&308 \\
\hline
float&1.0 & 0.0 \\
float& 10.0 & 1.0 \\
float& $1 \times 10^{-38}$& -38 \\
float& $1 \times 10^{38}$ &38 \\
\hline \hline
\end{tabular}
\end{center}

\subsubsection{FAIL\_log10\_protected}
This case tests the application of the \texttt{MathUtils::log10\_protected}
to an invalid input
(0.0), with the terminate-on-failure flag set to true. The simulation
terminates.

\subsection{Square-Root}
\subsubsection{RUN\_sqrt\_protected}
The table below lists the tests used to evaluate the response of
\texttt{MathUtils::sqrt\_protected}
\begin{center}
\begin{tabular}{|l||l|l|}
\hline
 \textbf{Data Type} & \textbf{Input} & \textbf{Expected Value} \\
\hline \hline
double &16.0& 4.0 \\
double & 2.0& 1.414213562373095 \textsuperscript{1} \\
double & 0.0& 0.0 \\
double & -16.0& 0.0 \textsuperscript{2} \\
\hline
float&16.0& 4.0 \\
float& 2.0& 1.4142135 \\
\hline \hline
\end{tabular}
\end{center}

\textsuperscript{1} Because of the sensitivity of the function to the input
values, this comparison agrees only to within 1.0 ULP when using data-type
\texttt{double}.

\textsuperscript{2} As the input is invalid and therefore invalid, it will broadcast 
an error and return 0.0.

\section{Numerical Methods}
Verification of these methods is found in runs within
\textit{SIM\_unit\_math\_utils\_private}, utilizing the following tests:
\begin{itemize}
 \item ERROR\_is\_within\_abs\_tolerance
 \item ERROR\_is\_within\_rel\_tolerance
 \item FAIL\_polynomial
 \item RUN\_has\_changed\_from
 \item RUN\_is\_equal
 \item RUN\_is\_near\_equal
 \item RUN\_is\_within\_abs\_tolerance
 \item RUN\_is\_within\_rel\_tolerance
 \item RUN\_polynomial
 \item RUN\_sign
\end{itemize}

\subsection{is\_near\_equal}
\subsubsection{RUN\_is\_near\_equal}

The table below lists the tests used in \\ \textit{RUN\_is\_near\_equal} to evaluate
the response of \texttt{MathUtils::is\_near\_equal}
\begin{center}
\begin{tabular}{|l||l|l|l|l|}
\hline
 \textbf{Data Type} & \textbf{Left} & \textbf{Right} &
 \textbf{ulp} & \textbf {Expected Result} \\
\hline \hline
double & 3.45 & 3.45 & None (0.5) & true \\
double & 1.0000000000000001 & 1.0& None (0.5) & true \\
double & 1.000000000000001& 1.0& None (0.5) & false \\
double & 5.4321 & 5.4322 & 1.0& false \\
double & 5.4321 & 5.4322 & $10^{7}$ & false \textsuperscript{*} \\
double & 5.4321 & 5.4322 & $10^{15}$& true \\
float &3.45 & 3.45 & None (0.5) & true \\
float &1.00000001 & 1.0& None (0.5) & true \\
float &1.0000001& 1.0& None (0.5) & false \\
float &5.4321 & 5.4322 & 1.0& false \\
float &5.4321 & 5.4322 & $10^{7}$ & true \textsuperscript{*}\\
\hline \hline
\end{tabular}
\end{center}
\textsuperscript{*} When using data-type \texttt{double}, the evaluation
$ 5.4321 == 5.4322$ to within $10^7 ULP$ returns false, but with data-type
\texttt{float} it returns true because the value of ULP is very much
larger with \texttt{float} than with \texttt{double}.

Results are as expected.

\subsection{has\_changed\_from}
\subsubsection{RUN\_has\_changed\_from}

The table below lists the tests in \\ \textit{RUN\_has\_changed\_from} used to evaluate
the response of \texttt{MathUtils::has\_changed\_from}
\begin{center}
\begin{tabular}{|l||l|l|l|l|}
\hline
 \textbf{Data Type} & \textbf{Left} & \textbf{Right} 
 & \textbf{Expected Result} \\
\hline \hline
double & 3.45 & 4.56 & true \\
double & 4.56 & 3.45 & true \\
double & 3.45 & 3.45 & false \\
float& 3.45 & 4.56 & true \\
float& 4.56 & 3.45 & true \\
float& 3.45 & 3.45 & false \\
\hline \hline
\end{tabular}
\end{center}

Results are as expected.

\subsection{is\_equal}
The table below lists the tests in \\ \textit{RUN\_is\_equal} used to evaluate the response of
\texttt{MathUtils::is\_equal}
\begin{center}
\begin{tabular}{|l||l|l|l|l|}
\hline
 \textbf{Data Type} & \textbf{Left} & \textbf{Right} 
 & \textbf{Expected Result} \\
\hline \hline
int& 3 & 3 & true \\
int& 3 & 4 & false \\
bool & 1 & 1 & true \\
bool & 0 & 1 & false \\
double & 3.45& 3.45& true \\
double & 3.45& 3.46& false \\
float& 3.45& 3.45& true \\
float& 3.45& 3.46& false \\
\hline \hline
\end{tabular}
\end{center}

Results are as expected.

\subsection{is\_within\_abs\_tolerance}
\subsubsection{RUN\_is\_within\_abs\_tolerance}
The table below lists the tests used to evaluate the response of \\
\texttt{MathUtils::is\_within\_abs\_tolerance}
\begin{center}
\begin{tabular}{|l||l|l|l|l|}
\hline
 \textbf{Data Type} & \textbf{Left} & \textbf{Right} &
 \textbf{Absolute Tolerance} & \textbf {Expected Value} \\
\hline \hline
double & 3.45 & 3.46 & 0.01 & true \\
double & 3.45 & 3.46 & 0.0099 & false \\
int& 3 & -2 & 5 & true \\
unsigned int & 3 & -2 & 5 & false \textsuperscript{1} \\
float & -2.0 & -2.01 & 0.0099 & false \\
double & -2.0 & -2.01 & -1.0 & false \textsuperscript{2} \\
bool &true & true & true & true \textsuperscript{*}\\
bool &true & true & false & true \textsuperscript{*}\\
bool &true & false & false & false \textsuperscript{*}\\
\hline \hline
\end{tabular}
\end{center}

\textsuperscript{1} The unsigned int (-2) wraps around to a large positive
integer.

\textsuperscript{2} The range is negative, this is automatically false, and 
broadcasts an error.

\textsuperscript{*} These tests don't make much sense and broadcasts an error
cause of it, but confirms that the model does handle boolean types.

Results are as expected.

\subsection{is\_within\_rel\_tolerance}
\subsubsection{RUN\_is\_within\_rel\_tolerance}
The table below lists the tests used to evaluate the response of \\
\texttt{MathUtils::is\_within\_rel\_tolerance}, checking 
whether the \textbf{Left} value lies within a range defined by multiplying
the \textbf{Relative Tolerance} with the \textbf{Right} value.
\begin{center}
\begin{tabular}{|l||l|l|l|l|}
\hline
 \textbf{Data Type} & \textbf{Left} & \textbf{Right} &
 \textbf{Relative Tolerance} & \textbf {Expected Value} \\
\hline \hline
double & 1.65 & 1.50 & 0.1 & true \\
double & 1.65 & 1.50 & 0.09999 & false \\
int& 2 & -2 & 2.0 & true \\
unsigned int & 2 & -2 & 2.0 & false \textsuperscript{1} \\
float & -1.65 & -1.50 & 0.1 & true \\
float & -2.0 & -2.01 & -1.0 & false \textsuperscript{2} \\
bool &true & true & true & true \textsuperscript{*}\\
bool &true & true & false & true \textsuperscript{*}\\
bool &true & false & false & false \textsuperscript{*}\\
\hline \hline
\end{tabular}
\end{center}

\textsuperscript{1} The unsigned int (-2) wraps around to a large positive
integer.

\textsuperscript{2} The range is negative, this is automatically false, and 
broadcasts an error.

\textsuperscript{*} These tests don't make much sense and broadcasts an error
cause of it, but confirms that the model does handle boolean types.

Results are as expected.

\subsection{sign}
The table below lists the tests in \textit{RUN\_sign} used to evaluate the response of
\texttt{MathUtils::sign}.
\begin{center}
\begin{tabular}{|l||l|l|l|l|}
\hline
 \textbf{Data Type} & \textbf{Value} & \textbf{Expected Value} \\
\hline \hline
int &5 & +1 \\
int &0 &0 \\
int & -5 & -1 \\
\hline \hline 
\end{tabular}
\end{center}

Results are as expected.


\subsection{Polynomial}
\subsubsection{RUN\_polynomial}
The table below lists the tests used to evaluate the response of
\texttt{MathUtils::polynomial}
\begin{center}
\begin{tabular}{|l||l|l|l|}
\hline
 \textbf{Coefficients} & \textbf{Variable} &
 \textbf{Recovery} & \textbf {Expected Value} \\
\hline \hline
\{3,4,5\} & 2.0& None (0.0) & 31.0 \\
\{3,4,5\} & $1 \times 10^{308}$& 75.7575& 75.7575
\textsuperscript{1} \\
\{0,4,5\} &$-1 \times 10^{-308}$ & None (0.0) &
$-4 \times 10^{-308}$ \\
\{0,4,5,-1,-1\} & -2.0 & None (0.0) & 4\\
\hline \hline
\end{tabular}
\end{center}
\textsuperscript{1} This evaluation overflows resulting in an error message being printed and the value assigned to the recovery value.

Results are as expected.

\subsubsection{FAIL\_polynomial}
This test sets the terminate-on-error flag; when the polynomial evaluation resulting in the overflow in \textit{RUN\_polynomial} is 
repeated with this flag, the simulation is terminated.

\subsection{Quadratic Solver}
\subsubsection{RUN\_quadratic\_solver}
The table below lists the tests used to evaluate the response of \\
\texttt{QuadraticSolver::update()}, computing the roots of a quadratic
equation.
\begin{center}
\begin{tabular}{|l|l|l||l|l|}
\hline
\textbf{Case} & \textbf{Form} & \makecell{\textbf{Quadratic} \\ \textbf{Equation}} & 
\makecell{\textbf{Expected} \\ \textbf{Roots}} & \makecell{\textbf{Roots} \\ \textbf{Exist}} \\
\hline \hline
Different Roots                & $ax^2+bx+c=0$ & $x^2-3x+2=0$ & $x=\{1,2\}$ & true \\
Same Root                      & $ax^2+bx+c=0$ & $x^2+2x+1=0$ & $x=\{-1\}$ & true \\
Linear Equation ($a=0$)        & $bx+c=0$ & $2x+1=0$ & $x=\{-0.5\}$ & true \\ 
$b=0$ with Real Roots          & $ax^2+c=0$ & $x^2-1=0$ & $x=\pm 1$ & true \\ 
$b=0$ with Imaginary Roots     & $ax^2+c=0$ & $x^2+1=0$ & $x=0$ & false \\ 
$c=0$                          & $ax^2+bx=0$ & $x^2+3x=0$ & $x=\{0,-3\}$ & true \\
No Solution ($a,b=0$)          & $c=0$ & $1=0$ & $x=0$ & false \\
$a,c=0$                        & $bx=0$ & $3x=0$ & $x=0$ & true \\
$b,c=0$                        & $ax^2=0$ & $3x^2=0$ & $x=0$ & true \\
Infinite Solutions ($a,b,c=0$) & $0=0$ & $0=0$ & $x=0$ & false \\
$b^2$ Overflows ($b >> 0$)     & $ax^2+bx+c=0$ & $x^2+(1e308)x+1$ & $x=\{-1e\text{-}308, -1e308\}$ & true \\
\hline \hline
\end{tabular}
\end{center}

Results are as expected.

\section{Matrix Methods}

\subsection{Cholesky Decomposition}

Cholesky Decomposition is found in
\textit{SIM\_cholesky\_decomposition}, with the following runs:
\begin{itemize}
 \item RUN\_01\_baseline
 \item RUN\_02\_3x3
 \item RUN\_03\_6x6
 \item RUN\_04\_negative
 \item RUN\_05\_legit\_zero
 \item RUN\_06\_illegitimate\_zero
 \item RUN\_07\_asymmetric
 \item RUN\_08\_four\_arg
 \item RUN\_09\_invalid\_size
 \item RUN\_10\_specA
 \item RUN\_11\_in\_array\_null
 \item RUN\_12\_out\_array\_null
 \item RUN\_13\_non\_symmetric\_array
\end{itemize}
For these tests, the following sequence of events if followed:
\begin{enumerate}
\item A 10x10 matrix is built and populated in and below the diagonal with
random numbers between 0 and 3. The random numbers have 4 significant
digits. The random seed is assignable so the generated matrix is
repeatable for a given test, but may vary between tests.
\item The diagonal half-matrix is multiplied by its transpose to generate a
symmetric matrix with positive determinant.
This matrix is identified as $[A]$.
\item The input file is processed; this allows manipulation of the values
of $[A]$ to produce $[A']$
\item the Cholesky decomposition is applied to the modified matrix $[A']$ to
create the decomposed matrix $[C]$, also referred to as the ``square-
root'' of $[A']$. $[C]$ is populated only in the lower diagonal.
\item The half-matrix $[C]$ is multiplied by its transpose to generate the
full matrix $[B] = [C][C]^T$. This matrix is compared against
$[A']$. If the
decomposition functioned properly, the two matrices $[A']$ and $[B]$
should be identical.
\end{enumerate}

\subsubsection{RUN\_01\_baseline}
This test simply executes the process as described above, with no modification
made to $[A]$ (so $[A'] = [A])$.

\subsubsection{RUN\_02\_3x3}
This test investigates the model capability to work with a sub-matrix in
computing the decomposition. In this case, $[A']$ is again left equal to
$[A]$; the submatrix of interest is the 3x3 matrix (from the upper left of
$[A'] = [A]$). The decomposed matrix contains only 6 non-zero elements, being
the lower diagonal of the submatrix

\begin{equation*}
 \begin{bmatrix}
 C_{00} & 0& 0& 0 & ... \\
 C_{10} & C_{11} & 0& 0 & ... \\
 C_{20} & C_{21} & C_{22} & 0 & ... \\
 0&0 & 0& 0 & ... \\
 ...
 \end{bmatrix}
\end{equation*}
The reconstructed matrix $[B]$ then contains nine non-zero elements populating
the 3x3 submatrix in the upper left of $[A']$.
These nine elements match the corresponding values in $[A] and [A']$.

\subsubsection{RUN\_03\_6x6}
This test further tests the submatrix capability, this time working with a 6x6
submatrix.

\subsubsection{RUN\_04\_negative}
In this test, the matrix $[A']$ is modified from $[A]$ in step 3
by resetting one value
to produce a matrix with a negative determinant. For the default matrix
generated from random numbers with the prescribed seed, setting
$A'_{33} = 0.0$ achieves the goal.

For this test, the decomposition fails, an error is printed and the processing
is terminated in step 4. Generation of the matrix $[C]$ is not completed.

\subsubsection{RUN\_05\_legit\_zero}
For this run, we test the case for decomposing a matrix with a
zero- determinant. The modified matrix $[A']$ is generated with all
values in the third column and third row being set to 0.0.

A warning is posted that the decomposition is not unique, but the decomposed
matrix $[C]$ is successfully used to reconstruct the modified matrix $[B] =
[A']$.

\subsubsection{RUN\_06\_illegitimate\_zero}
This test works much like \textit{RUN\_04\_negative}; this time the component
$A'_{00}$ is assigned to zero, again producing a negative determinant. As
with \textit{RUN\_04\_negative}, an error is printed and the processing
is terminated in step 4. Generation of the matrix $[C]$ is not completed.

\subsubsection{RUN\_07\_asymmetric}
For this test, we change an off-diagonal element in $[A']$ (setting
$A'_{14} = 0.0$, making $[A']$ asymmetric. Because the Cholesky decomposition
assumes symmetry, the overwritten value in the upper diagonal is lost; when
comparing $[B]$ and $[A']$, the element at $\{14\}$ does not match. However,
$[B]$ does match with $[A]$, as should be expected because the values in the
lower diagonal (the values used in the decomposition) were not modified.

Note that in this case, the asymmetry was against a value 0.0; this does not
trigger any error messages because values of 0.0 could be simply ``not
populated'', which is not an error. See \textit{RUN\_13} for an asymmetry with
a non-zero value.

\subsubsection{RUN\_08\_four\_arg}
In the runs to this point, we have been using the 5-argument version of
\newline
\textit{MathUtils::cholesky\_decomposition(...)}, including the specification
of the sub-matrix size (which, except for \textit{RUN\_02} and
\textit{RUN\_03}, had a value of
10 to force processing of the entire matrix). In this test, we use the
4-argument version of the method. Thus the
entire matrix is processed, so this test is effectively equivalent to
\textit{RUN\_01}.

\subsubsection{RUN\_09\_invalid\_size}
In this test, we switch back to the 5-argument call and specify the submatrix
size to be $14 \times 14$. This is larger than the available matrix; the
processing
is terminated in step 4. Generation of the matrix $[C]$ is not completed.


\subsubsection{RUN\_10\_specA}
For this test, the modified matrix $[A']$ is populated entirely in the
input file, with the top $4 \times 4$ matrix elements populated with non-zero
values and the rest populated with zeroes. The processing is instructed to use
this 4x4 submatrix. The resulting matrix is successfully decomposed and
reconstructed.

\subsubsection{RUN\_11\_in\_array\_null}
Because the values passed to the method are C-style arrays, the actual
arguments are pointers to the array memory space. As such, the arguments may
be passed as NULL. This test serves as a check that if the pointer addressing
the input array is NULL, the decomposition will cleanly error out. Processing
is terminated in step 4. Generation of the matrix $[C]$ is not completed.


\subsubsection{RUN\_12\_out\_array\_null}
This test partners with \textit{RUN\_11}; this time the pointer to the output
array is passed in as NULL. As with \textit{RUN\_11}, processing
is terminated in step 4. Generation of the matrix $[C]$ is not completed.

\subsubsection{RUN\_13\_non\_symmetric\_array}
This test works in conjunction with \textit{RUN\_07}, again modifying $[A']$
to make it asymmetric. This time, the modified values are non-zero and the
internal checking for symmetry produces an error message. For this test, we
make two values asymmetric to confirm that only one error message is sent
regardless of how many elements are not symmetric.

\subsection{Position-Velocity Covariance Matrix Transformation}
The following tests in \textit{RUN\_pv\_covariance\_matrix} evaluate the output of the 
\texttt{transform\_pv\_matrix} method by using a transformation matrix from 
the original frame to a new frame, denoted as $T$, and a position-velocity 
covariance matrix $P$ defined in the original frame. The expected covariance 
matrix defined in the new frame for each test case are described below.

\begin{enumerate}
\item{
A zero covariance matrix remains unchanged under any transformation,
as expected.
\begin{equation*}
T = 
\begin{bmatrix}
1 & 0 & 0 \\
0 & 0 & 1 \\
0 & -1 & 0
\end{bmatrix}, \quad
P = 
\begin{bmatrix}
0 & 0 & 0 & 0 & 0 & 0 \\
0 & 0 & 0 & 0 & 0 & 0 \\
0 & 0 & 0 & 0 & 0 & 0 \\
0 & 0 & 0 & 0 & 0 & 0 \\
0 & 0 & 0 & 0 & 0 & 0 \\
0 & 0 & 0 & 0 & 0 & 0
\end{bmatrix} \Longrightarrow
P = 
\begin{bmatrix}
0 & 0 & 0 & 0 & 0 & 0 \\
0 & 0 & 0 & 0 & 0 & 0 \\
0 & 0 & 0 & 0 & 0 & 0 \\
0 & 0 & 0 & 0 & 0 & 0 \\
0 & 0 & 0 & 0 & 0 & 0 \\
0 & 0 & 0 & 0 & 0 & 0
\end{bmatrix}
\end{equation*}}
\item{
An isotropic diagonal covariance matrix remains unchanged under rotation. 
This is expected because all variances are equal, and the transformation 
simply rotates the coordinate axes. Since the distribution is symmetric in 
all directions, rotating the frame does not alter the covariance matrix.
\begin{equation*}
T = 
\begin{bmatrix}
1 & 0 & 0 \\
0 & 0 & 1 \\
0 & -1 & 0
\end{bmatrix}, \quad
P = 
\begin{bmatrix}
5 & 0 & 0 & 0 & 0 & 0 \\
0 & 5 & 0 & 0 & 0 & 0 \\
0 & 0 & 5 & 0 & 0 & 0 \\
0 & 0 & 0 & 5 & 0 & 0 \\
0 & 0 & 0 & 0 & 5 & 0 \\
0 & 0 & 0 & 0 & 0 & 5
\end{bmatrix} \Longrightarrow
P = 
\begin{bmatrix}
5 & 0 & 0 & 0 & 0 & 0 \\
0 & 5 & 0 & 0 & 0 & 0 \\
0 & 0 & 5 & 0 & 0 & 0 \\
0 & 0 & 0 & 5 & 0 & 0 \\
0 & 0 & 0 & 0 & 5 & 0 \\
0 & 0 & 0 & 0 & 0 & 5
\end{bmatrix}
\end{equation*}}
\item{
The values in the covariance matrix were intentionally chosen to be 
visually traceable: the diagonal entries 1 through 6 correspond to the 
state components $x$, $y$, $z$, $v_x$, $v_y$, and $v_z$, respectively. 
Off-diagonal terms such as 12, 23, or 45 represent covariances between 
those state components (e.g., 12 corresponds to the covariance between 
$x$ and $y$). This pattern makes it easier to verify how each element 
is affected by the transformation. Note that the same covariance matrix $P$ 
is used in the following two test cases as well.

The transformation matrix $T$ represents a 90° rotation about the x-axis, 
where the new $y$-axis aligns with the original $z$-axis, and the new 
$z$-axis aligns with the negative of the original $y$-axis. This explains 
the pattern observed in the transformed covariance matrix: elements 
involving $y$ and $z$ are swapped, and appropriate signs are flipped 
to reflect the axis inversion. The result confirms that the transformation 
correctly reorients the covariance matrix according to a simple 90° 
rotation around the x-axis.
\begin{equation*}
T = 
\begin{bmatrix}
1 & 0 & 0 \\
0 & 0 & 1 \\
0 & -1 & 0
\end{bmatrix}, \quad
P = 
\begin{bmatrix}
 1 & 12 & 13 & 14 & 15 & 16 \\
12 &2 & 23 & 24 & 25 & 26 \\
13 & 23 &3 & 34 & 35 & 36 \\
14 & 24 & 34 &4 & 45 & 46 \\
15 & 25 & 35 & 45 &5 & 56 \\
16 & 26 & 36 & 46 & 56 &6
\end{bmatrix} \Longrightarrow
P = 
\begin{bmatrix}
 1&13& -12&14&16& -15 \\
13& 3& -23&34&36& -35 \\
 -12& -23& 2& -24& -26&25 \\
14&34& -24& 4&46& -45 \\
16&36& -26&46& 6& -56 \\
 -15& -35&25& -45& -56& 5
\end{bmatrix}
\end{equation*}}
\item{
The transformation matrix $T$ represents a 90° rotation about the y-axis, 
where the new $z$-axis aligns with the original $x$-axis, and the new 
$x$-axis aligns with the negative of the original $z$-axis. This explains 
the pattern observed in the transformed covariance matrix: elements 
involving $x$ and $z$ are swapped, and appropriate signs are flipped 
to reflect the axis inversion. The result confirms that the transformation 
correctly reorients the covariance matrix according to a simple 90° 
rotation around the y-axis.
\begin{equation*}
T = 
\begin{bmatrix}
0 & 0 & -1 \\
0 & 1 &0 \\
1 & 0 & 0
\end{bmatrix}, \quad
P = 
\begin{bmatrix}
 1 & 12 & 13 & 14 & 15 & 16 \\
12 &2 & 23 & 24 & 25 & 26 \\
13 & 23 &3 & 34 & 35 & 36 \\
14 & 24 & 34 &4 & 45 & 46 \\
15 & 25 & 35 & 45 &5 & 56 \\
16 & 26 & 36 & 46 & 56 &6
\end{bmatrix} \Longrightarrow
P = 
\begin{bmatrix}
 3& -23 & -13 &36 & -35 & -34 \\
 -23& 2 &12 & -26 &25 &24 \\
 -13&12 & 1 & -16 &15 &14 \\
36& -26 & -16 & 6 & -56 & -46 \\
 -35&25 &15 & -56 & 5 &45 \\
 -34&24 &14 & -46 &45 & 4
\end{bmatrix}
\end{equation*}}
\item{
The transformation matrix $T$ represents a 90° rotation about the z-axis, 
where the new $x$-axis aligns with the original $y$-axis, and the new 
$y$-axis aligns with the negative of the original $x$-axis. This explains 
the pattern observed in the transformed covariance matrix: elements 
involving $x$ and $y$ are swapped, and appropriate signs are flipped 
to reflect the axis inversion. The result confirms that the transformation 
correctly reorients the covariance matrix according to a simple 90° 
rotation around the z-axis.
\begin{equation*}
T = 
\begin{bmatrix}
0 & 1 & 0 \\
-1 & 0 & 0 \\
0 & 0 & 1
\end{bmatrix}, \quad
P = 
\begin{bmatrix}
 1 & 12 & 13 & 14 & 15 & 16 \\
12 &2 & 23 & 24 & 25 & 26 \\
13 & 23 &3 & 34 & 35 & 36 \\
14 & 24 & 34 &4 & 45 & 46 \\
15 & 25 & 35 & 45 &5 & 56 \\
16 & 26 & 36 & 46 & 56 &6
\end{bmatrix} \Longrightarrow
P = 
\begin{bmatrix}
 2& -12 &23 &25 & -24 &26 \\
 -12& 1 & -13 & -15 &14 & -16 \\
23& -13 & 3 &35 & -34 &36 \\
25& -15 &35 & 5 & -45 &56 \\
 -24&14 & -34 & -45 & 4 & -46 \\
26& -16 &36 &56 & -46 & 6
\end{bmatrix}
\end{equation*}}
\item{
This verification case uses a set of arbitrarily selected position-velocity state 
``measurements" in an inertial frame to compute an initial $6 \times 6$ covariance 
matrix. This matrix serves as the input to the method, to output a transformed 
covariance matrix expressed in a different reference frame. To 
validate the result, each of the original inertial state vectors are individually
transformed into the reference frame, and a new covariance matrix is computed 
directly from these transformed states. Therefore, this covariance matrix in 
the reference frame can be directly compared to the output of the method.

The following computations to setup the verification were performed 
using the provided Octave script, 
\textit{SIM\_unit\_math\_utils/pv\_covariance\_analysis.m}. The script contains 10 
arbitrarily chosen position-velocity state vectors in the inertial frame and 
calculates the following covariance matrix in the inertial frame as input to the method:
\begin{equation*}
P = 
\begin{bmatrix}
8333.33	 & 8500	& -5166.67 &	11.9444	 & 8	& -5.72222 \\
8500	 & 9098.89	& -5572.22 &	12.6122	 & 7.93333	& -6.11444 \\
-5166.67 & -5572.22	& 3494.44	 & -7.73889	 & -4.83333	& 3.83889 \\
11.9444	 & 12.6122	& -7.73889 &	0.0177656	 & 0.0112778	& -0.00858556 \\
8	 & 7.93333	& -4.83333 &	0.0112778	 & 0.00786667	& -0.00542222 \\
-5.72222 & -6.11444	& 3.83889	 & -0.00858556 &	-0.00542222	&0.00429
\end{bmatrix}
\end{equation*}

The following transformation matrix corresponding to a 3-2-1 Euler rotation sequence with 
arbitrarily chosen angles 321\textdegree, 243\textdegree, and 56\textdegree\text{ }
is applied to the above covariance matrix in the inertial frame:
\begin{equation*}
T = 
\begin{bmatrix}
-0.3528	& 0.2857	& 0.891 \\
-0.2221	& 0.8994	& -0.3764 \\
-0.9089	& -0.3307	& -0.2539
\end{bmatrix}
\end{equation*}

The expected covariance matrix
in the reference frame is computed by transforming each state vector using this
rotation and then calculating the covariance from the transformed set of states:
\begin{equation*}
P = 
\begin{bmatrix}
3252.05	 & -4947.23	&5574.43 &4.54908	 & -4.09574	 &7.49346 \\
-4947.23 &	7779.14	& -8494.81 & -7.09806	 &6.01336	 & -11.4802 \\
5574.43	 & -8494.81	&9895.48 &7.70011	 & -7.24646	 &13.1542 \\
4.54908	 & -7.09806	&7.70011 &0.00662308 & -0.00552239 & 	0.010546\\
-4.09574 & 	6.01336	& -7.24646 & -0.00552239 & 	0.00557713 & -0.00959254\\
7.49346	 & -11.4802	&13.1542 &0.010546	 & -0.00959254 &0.017722
\end{bmatrix}
\end{equation*}
}
\item{
This verification case builds directly on the previous one by using the resulting 
covariance matrix in the reference frame as the new input. The goal is to confirm 
that the transformation process is reversible and preserves the statistical 
data. The previously computed covariance matrix in the reference frame:
\begin{equation*}
P = 
\begin{bmatrix}
3252.05	 & -4947.23	&5574.43 &4.54908	 & -4.09574	 &7.49346 \\
-4947.23 &	7779.14	& -8494.81 & -7.09806	 &6.01336	 & -11.4802 \\
5574.43	 & -8494.81	&9895.48 &7.70011	 & -7.24646	 &13.1542 \\
4.54908	 & -7.09806	&7.70011 &0.00662308 & -0.00552239 & 	0.010546\\
-4.09574 & 	6.01336	& -7.24646 & -0.00552239 & 	0.00557713 & -0.00959254\\
7.49346	 & -11.4802	&13.1542 &0.010546	 & -0.00959254 &0.017722
\end{bmatrix}
\end{equation*}

is transformed back into the inertial frame using the inverse of the 
original transformation matrix:
\begin{equation*}
T = 
\begin{bmatrix}
-0.3528	& -0.2221 &	-0.9089 \\
 0.2857 &0.8994 &	-0.3307 \\
 0.891	& -0.3764 &	-0.2539
\end{bmatrix}
\end{equation*}

The expected result of this reverse transformation is the recovery of the 
original covariance matrix in the inertial frame:
\begin{equation*}
P = 
\begin{bmatrix}
8333.33	 & 8500	& -5166.67 &	11.9444	 & 8	& -5.72222 \\
8500	 & 9098.89	& -5572.22 &	12.6122	 & 7.93333	& -6.11444 \\
-5166.67 & -5572.22	& 3494.44	 & -7.73889	 & -4.83333	& 3.83889 \\
11.9444	 & 12.6122	& -7.73889 &	0.0177656	 & 0.0112778	& -0.00858556 \\
8	 & 7.93333	& -4.83333 &	0.0112778	 & 0.00786667	& -0.00542222 \\
-5.72222 & -6.11444	& 3.83889	 & -0.00858556 &	-0.00542222	&0.00429
\end{bmatrix}
\end{equation*}
}
\end{enumerate}

The results for all these cases are as expected.

\subsection{Correlation Matrix}
The following tests in \textit{RUN\_correlation\_coefficients}
evalute the response of the \\ \texttt{extract\_correlation\_coefficients}
method with different covariance matrices, $P$. Since the method assumes the 
covariance matrix is populated at least on the lower diagonal, the upper
triangle is filled with larger values to verify they are ignored in the
calculations. For easier understanding, the covariance matrices can be 
visualized as symmetric, where only the lower diagonal contains meaningful
data and the upper triangle has garbage data. The expected correlation matrices, $\rho$ are
described below.

\begin{enumerate}
\item{
An ordinary covariance matrix with all positive variances (diagonals) and
covariances (off-diagonals) simply applies
(\ref{eqn:corr_coeff}):
\begin{equation*}
P = 
\begin{bmatrix}
4 & 100 & 100 \\
2 & 3 & 100 \\
0.6 & 0.9 & 2
\end{bmatrix} \quad \Longrightarrow \quad 
\rho = 
\begin{bmatrix}
1 & 0.577 & 0.212 \\
0.577 & 1 & 0.367 \\
0.212 & 0.367 & 1
\end{bmatrix}
\end{equation*}
}
\item{
A covariance matrix with all positive variances (positive diagonal elements) can also
contain zero or negative covariances (off-diagonal elements). In such cases, 
(\ref{eqn:corr_coeff}) still applies:
\begin{equation*}
P = 
\begin{bmatrix}
4 & 100 & 100 \\
-1.2 & 3 & 100 \\
0.6 & -0.9 & 2.5
\end{bmatrix} \quad \Longrightarrow \quad 
\rho = 
\begin{bmatrix}
1 & -0.346 & 0.190 \\
-0.346 & 1 & -0.329 \\
0.190 & -0.329 & 1
\end{bmatrix}
\end{equation*}
}
\item{
A covariance matrix with any negative variances (i.e. negative diagonal elements) 
like below isn't defined as variables cannot be
negatively correlated with themselves:
\begin{equation*}
P = 
\begin{bmatrix}
-1 & 100 & 100 \\
0 & -2 & 100 \\
0 & 0 & -3
\end{bmatrix} 
\end{equation*}
The method is expected to broadcast an appropriate error.
}
\item{
A covariance matrix with any only positive variances (i.e. positive diagonal elements
and off-diagonal elements of zero) like below can simply apply 
(\ref{eqn:corr_coeff}) to produce an identity matrix of that size, because the variables
will only be correlated with themselves (i.e. the diagonal is filled with ones and
the rest is zero):
\begin{equation*}
P = 
\begin{bmatrix}
1 & 100 & 100 \\
0 & 2 & 100 \\
0 & 0 & 3
\end{bmatrix} \quad \Longrightarrow \quad 
\rho = 
\begin{bmatrix}
1 & 0 & 0 \\
0 & 1 & 0 \\
0 & 0 & 1
\end{bmatrix} 
\end{equation*}
}
\item{
A covariance matrix with any zero variances and non-zero respective covariances also isn't 
defined (i.e. diagonal elements of zero with non-zero elements in that column and row), 
because a variable without variance cannot have a covariance with other variables. This test 
case includes a zero variance in the middle variable and a non-zero covariance to its left.
It is specifically designed to trigger the error-checking logic that inspects the columns
of the lower triangle in a specific row for non-zero elements 
(the first \texttt{i\_lower\_col} loop).
\begin{equation*}
P = 
\begin{bmatrix}
1 & 100 & 100 \\
1 & 0 & 100 \\
0 & 0 & 1
\end{bmatrix} 
\end{equation*}
The method is expected to broadcast an appropriate error.
}
\item{
This test case is similar to the last test, with a zero variance in the middle variable 
again, but with a non-zero covariance below it. It is specifically designed to trigger 
the error-checking logic that inspects the rows of the lower triangle in a specific column
for non-zero elements (the \texttt{i\_lower\_row} loop).
\begin{equation*}
P = 
\begin{bmatrix}
1 & 100 & 100 \\
0 & 0 & 100 \\
0 & 1 & 1
\end{bmatrix} 
\end{equation*}
The method is expected to broadcast an appropriate error.
}
\item{
A covariance matrix with any zero variances and zero respective covariances 
(i.e. diagonal elements of zero with elements of zero in that column and row) would be
defined as degenerate random variables, where the variables would have 100\% certainty.
These cases aren't regularly found, and according to (\ref{eqn:corr_coeff}) the 
correlation coefficients for that column and row would be undefined. However, for
computing purposes, the correlation coefficient for the diagonal would be set to one, and the
coefficients for the column and row would result in zero:
\begin{equation*}
P = 
\begin{bmatrix}
0 & 100 & 100 \\
0 & 1 & 100 \\
0 & 1 & 1
\end{bmatrix} \quad \Longrightarrow \quad
\rho = 
\begin{bmatrix}
1 & 0 & 0 \\
0 & 1 & 1 \\
0 & 1 & 1
\end{bmatrix}
\end{equation*}
The method is also expected to log an appropriate inform message about the possible
ambiguity of the undefined correlation coefficients being calculated as if the 
first diagonal element was one.
}
\item{
A covariance matrix with all zero variance and covariances (i.e. all elements are zero)
would result in an identity matrix of that size, because the variables will only be 
correlated with themselves (i.e. the diagonal is filled with ones and the rest is zero):
\begin{equation*}
P = 
\begin{bmatrix}
0 & 100 & 100 \\
0 & 0 & 100 \\
0 & 0 & 0
\end{bmatrix} \quad \Longrightarrow \quad
\rho = 
\begin{bmatrix}
1 & 0 & 0 \\
0 & 1 & 0 \\
0 & 0 & 1
\end{bmatrix}
\end{equation*}
The method is also expected to log three appropriate inform messages about the possible
ambiguity of the undefined correlation coefficients being calculated as if the 
all the diagonal elements were one.
}
\end{enumerate}

The resulting correlation matrix for all these cases are as expected, and the matrices
are further confirmed to symmetric and have diagonals of only ones.

\subsection{Vector and Matrix Algebra}
The verification of the basic vector operations is conducted in 
\textit{SIM\_std\_array\_ops/RUN\_test\_arr\_ops}, where the following
three outputs are reported after every operation to display how they change
as the sequence of operations are applied:
\begin{itemize}
    \item scalar: scalar value 
    \item out\_3: \texttt{std::array} of size 3
    \item out\_4: \texttt{std::array} of size 4
\end{itemize}

In the following description, the subscript on each vector indicates its size.
If a superscript $C$ is present (e.g $in^C$), the vector is a C-style array, otherwise it is
an \texttt{std::array}. The test run \textit{RUN\_test\_arr\_ops} performs the following 
sequence of operations 
with inputs $in^C_3 = [3.1, 3.2, 3.3]$ and $in^C_4 = [4.1, 4.2, 4.3, 4.4]$:
\begin{enumerate}
    \item \texttt{std\_copy} is applied to $in^C_3$ and $in^C_4$ to generate: \\
        $out_3 = [3.1, 3.2, 3.3]$  \\ $out_4 = [4.1, 4.2, 4.3, 4.4]$.
    \item \texttt{std::array operator=} is applied to generate copies of the inputs for 
        later operations: \\
        $a_3 = [3.1, 3.2, 3.3]$ \\
        $a_4 = [4.1, 4.2, 4.3, 4.4]$.
    \item \texttt{std::array operator+=} is applied with lvalue arrays to generate: \\
        $out_3 = out_3 + a_3 = [6.2, 6.4, 6.6]$ \\ 
        $out_4 = out_4 + a_3 = [8.2, 8.4, 8.6, 8.8]$.
    \item \texttt{std::array operator+=} is applied with rvalue arrays to generate: \\
        $out_3 = out_3 + [1.0, 2.0, 3.0] = [7.2, 8.4, 9.6]$  \\ 
        $out_4 = out_4 + [1.0, 2.0, 3.0, 4.0] = [9.2, 10.4, 11.6, 12.8]$.
    \item \texttt{std::array operator=} is applied to create copies of $out_3$ and $out_4$ for 
        later operations: \\
        $b_3 = [7.2, 8.4, 9.6]$  \\ $b_4 = [9.2, 10.4, 11.6, 12.8]$.
    \item \texttt{std::array operator-=} is applied with lvalue and rvalue arrays to generate: \\
        $out_3 = out_3 - a_3 = [4.1, 5.2, 6.3]$  \\ 
        $out_4 = out_4 - [1.5, 2.5, 3.5, 4.5] = [7.7, 7.9, 8.1, 8.3]$.
    \item \texttt{std::array operator-} is applied to generate: \\
        $out_3 = b_3 - a_3 = [4.1, 5.2, 6.3]$  \\ 
        $out_4 = b_4 - a_4 = [5.1, 6.2, 7.3, 8.4]$.
    \item \texttt{std::array operator*=} is appled with rvalues of type \texttt{double} 
        to generate: \\
        $out_3 = out_3 * 3 = [12.3, 15.6, 18.9]$  \\ 
        $out_4 = out_4 * 3 = [15.3, 18.6, 21.9, 25.2]$.
    \item \texttt{std::array operator*=} and \texttt{std::array operator*} are applied with 
        lvalues arrays and rvalues of type \texttt{double}, respectively
        to generate: \\
        $out_3 = out_3 * (a_3 * 2) = [76.26, 99.84, 124.74]$  \\ 
        $out_4 = out_4 * (a_4 * 2) = [125.46, 156.24, 188.34, 221.76]$.
    \item \texttt{std::array operator*=} is applied again with lvalue arrays to generate: \\
        $out_3 = out_3 * a_3 = [236.406, 319.488, 411.642]$  \\ 
        $out_4 = out_4 * a_4 = [514.386, 656.208, 809.862, 975.744]$.
    \item \texttt{std::array operator*} is applied to generate: \\
        $out_3 = b_3 * a_3 = [22.32, 26.88, 31.68]$  \\ 
        $out_4 = b_4 * a_4 = [37.72, 43.68, 49.88, 56.32]$.
    \item \texttt{std::array operator/} is applied to generate: \\
        $out_3 = b_3 / a_3 = [2.32258, 2.625, 2.90909]$  \\ 
        $out_4 = b_4 / a_4 = [2.2439, 2.47619, 2.69767, 2.90909]$.
    \item \texttt{vector\_scalar\_product} is applied twice to generate:
        \begin{itemize}
            \item $scalar = b_3 \cdot a_3 = 80.88$ 
            \item $scalar = b_4 \cdot a_4 = 187.6$
        \end{itemize}
    \item \texttt{vector\_cross\_product} is applied four times to generate:
        \begin{itemize}
            \item $out_3 = b_3 \times a_3 = [-3, 6, -3]$ 
            \item $out_3 = b_3 \times in^C_3 = [-3, 6, -3]$ 
            \item $out_3 = in^C_3 \times b_3 = [3, -6, 3]$ 
            \item $out_3 = in^C_3 \times in^C_3 = [0, 0, 0]$ 
        \end{itemize}
    \item \texttt{vec\_mag\_sq} is applied to generate: \\
        $ scalar = ||a_4||^2 = 72.3$.
    \item \texttt{vec\_mag} is applied to generate: \\
        $ scalar = ||a_4|| = 8.50294$.
    \item \texttt{vec\_mag\_xy} is applied to generate: \\
        $ scalar = ||a_3||_{xy} = 4.45533$.
    \item \texttt{vec\_mag\_xy} is applied to generate: \\
        $ scalar = ||in^C_3||_{xy} = 4.45533$.
    \item \texttt{vec\_mag\_xz} is applied to generate: \\
        $ scalar = ||a_3||_{xz} = 4.52769$.
    \item \texttt{vec\_mag\_xz} is applied to generate: \\
        $ scalar = ||in^C_3||_{xz} = 4.52769$.
    \item \texttt{vec\_mag\_yz} is applied to generate: \\
        $ scalar = ||a_3||_{yz} = 4.59674$.
    \item \texttt{vec\_mag\_yz} is applied to generate: \\
        $ scalar = ||in^C_3||_{yz} = 4.59674$.
    \item \texttt{unit\_vector} is applied twice to generate:
        \begin{itemize}        
            \item $ out_3 = \frac{a_3}{||a_3||} = [0.559126, 0.577162, 0.595199]$ \\
                  $ out_4 = \frac{a_4}{||a_4||} = [0.482186, 0.493947, 0.505707, 0.517468]$
            \item $ out_3 = \frac{\vec{0}}{||\vec{0}||} = [0, 0, 0]$
        \end{itemize}
    \item \texttt{zero\_vector} is applied to generate: \\
        $ out_3 = \vec{0} = [0, 0, 0]$ \\
        $ out_4 = \vec{0} = [0, 0, 0, 0]$.
    \item \texttt{vector\_elementwise\_sqrt} is applied to generate: \\
        $ out_3 = \sqrt{a_3} = [1.76068, 1.78885, 1.81659]$ \\
        $ out_4 = \sqrt{a_4} = [2.02485, 2.04939, 2.07364, 2.09762]$.
    \item \texttt{diff} is applied to generate: \\
        $ out_3 = in^C_3 - in^C_3 = [0, 0, 0]$.
    \item \texttt{sum} is applied to generate: \\
        $ out_3 = in^C_3 + in^C_3 = [6.2, 6.4, 6.6]$.
\end{enumerate}

The verification of the basic matrix operations is conducted in
\textit{SIM\_unit\_math\_utils/RUN\_matrix},
with verification provided in the Octave script
\textit{SIM\_unit\_math\_utils/matrix\_analysis.m}. The verification script
reads in the logged data and applies the same matrix operations to the same
inputs,
then compares the results of those operations with the logged output data.

In the following description, the subscripts on each vector or matrix indicates the
size of the structure. The run \textit{RUN\_matrix} executes the same sequence
of operations for a set of four initial conditions:
\begin{enumerate}
    \item The matrices $[A]_{42}$ and $[B]_{33}$ are populated from file
        \textit{Unit\_test/math\_utils\_matrix\_data.txt}.
    \item{Matrix $[C]_{43}$ and vector $D_8$ are constructed manually from the values in
        $[A]_{42}$, with
            \begin{itemize}
                \item $C_{ij} =
                        \begin{cases}
                            A_{ij} & \qquad \text{for } 0 \leqslant i < 4, \quad 0 \leqslant j < 2 \\
                            1 & \qquad \text{for } 0 \leqslant i < 4, \quad j = 2 \\
                        \end{cases}$
                \item $D_{k} = A_{ij} \qquad \text{where } k=2i+j, \qquad 0 \leqslant i < 4, \quad 0 \leqslant j < 2$
            \end{itemize}}
    \item \texttt{matrix\_trans} is applied to $[A]_{42}$, generating
        $[E]_{24}$.
    \item \texttt{matrix\_mult\_left\_trans} is applied to generate
        $[F]_{32} = ([C]_{43})^T [A]_{42}$.
    \item \texttt{matrix\_mult\_right\_trans} is applied to generate
        $[G]_{34} = [F]_{32} ([A]_{42})^T$.
    \item \texttt{matrix\_mult\_trans\_trans} is applied to generate
        $[H]_{23} = ([A]_{42})^T ([G]_{34})^T$.
    \item \texttt{matrix\_mult} is applied to generate
        $[J]_{43} = [A]_{42} [H]_{23}$.
    \item \texttt{col\_maj\_vec\_to\_matrix} is applied to generate
        $[K]_{42}$ by reshaping the vector $D_8$ into a $4 \times 2$ matrix
        using column-major order.
    \item \texttt{row\_maj\_vec\_to\_matrix} is applied to generate
        $[L]_{42}$ by reshaping the vector $D_8$ into a $4 \times 2$ matrix
        using row-major order.
    \item \texttt{matrix\_inv\_using\_cholesky} is applied to generate
        $[M]_{33} = [B]_{33}^{-1}$
    \item \texttt{matrix\_copy} is applied to generate
        $[N]_{33} = [M]_{33}$
    \item \texttt{matrix\_incr} is applied to $[N]_{33}$:
        $[N]_{33} = [N]_{33} + [B]_{33}$
    \item \texttt{matrix\_copy} is applied to generate
        $[O]_{33} = [N]_{33}$
    \item \texttt{matrix\_decr} is applied to $[O]_{33}$:
        $[O]_{33} = [O]_{33} - [B]_{33} = [M]_{33}$
    \item \texttt{matrix\_transformation} is applied to generate
        $[P]_{33} = [B]_{33} [M]_{33} [B]_{33}^T = [B]_{33}^T$.
        This identity holds because $[M]_{33}=[B]^{-1}_{33}$, so the product
        simplifies as:
        \[
        [P]_{33} = [B]_{33}[B]^{-1}_{33}[B]^T_{33} = [I]_{33}[B]^T_{33} = [B]^T_{33}
        \]

        If $[B]_{33}$ is symmetric, then $[B]^T_{33} = [B]_{33}$, and thus $[P]_{33}=[B]_{33}$.
        This symmetry is required for the Cholesky-based inversion used in step 8, and therefore
        is a necessary condition for the test sequence.

        Furthermore, if $[B]_{33}$ is a transformation matrix (i.e. orthonormal), 
        then $[B]^T_{33} = [B]^{-1}_{33} = [M]_{33}$,
        and the transformation simplifies to:
        \[
        [P]_{33} = [B]_{33} [M]_{33} [B]^T_{33} = [B]_{33} [B]^{-1}_{33} [B]^T_{33} = [B]^T_{33} = [M]_{33}
        \]
        The only matrix in the test cases that satisfies both symmetry and orthonormality
        is the identity matrix from test case 1.

    \item \texttt{matrix\_inverse\_transformation} is applied to generate
        $[Q]_{33} = [M]^T_{33} [M]_{33} [M]_{33}$.

        If $[M]_{33}$ is a transformation matrix (i.e. orthonormal), then it satisfies
        $[M]^T_{33} = [M]^{-1}_{33}$, and the expression simplifies as:
        \[
        [Q]_{33} = [M]^T_{33} [M]_{33} [M]_{33} = [M]^{-1}_{33} [M]_{33} [M]_{33} = [I]_{33} [M]_{33} = [M]_{33}
        \]

        This simplification occurs only in the first test case, where $[B]_{33}$ 
        is the identity matrix, because it is the only matrix that is  
        both symmetric and a transformation matrix.

    \item \texttt{zero\_matrix} is applied to generate $[R]_{33} = [0]_{33}$.
    \item \texttt{matrix\_copy} is applied to generate $[S]_{33} = [M]_{33}$.
    \item \texttt{matrix\_scale} is applied to $[S]_{33}$: $[S]_{33} = 2.5 * [S]_{33}$.
    \item Matrices $[T]_{33}$ and $[U]_{22}$ are manually populated with elements 
        of a larger number like $100$.
    \item \texttt{matrix\_copy} is applied four times to generate:
            \begin{itemize}
                \item $[V]_{22} = [U]_{22}$
                \item $[W]_{22} = [U]_{22}$
                \item $[X]_{33} = [T]_{33}$
                \item $[Y]_{33} = [T]_{33}$
            \end{itemize}
        \item \texttt{matrix\_copy\_submatrix\_out} is applied three times with different starting indices to extract 
        $2 \times 2$ submatrices from $[A]_{42}$:
        \[
        \begin{aligned}
        [U]_{22} &= ([A]_{42})_{ij} \qquad \text{for } 2 \leqslant i < 4,\quad 0 \leqslant j < 2 \quad &\text{(starting at } (2,0)) \\
        [V]_{22} &= ([A]_{42})_{ij} \qquad \text{for } 2 \leqslant i < 4,\quad 1 \leqslant j < 2 \quad &\text{(starting at } (2,1)) \\
        [W]_{22} &= ([A]_{42})_{ij} \qquad \text{for } 3 \leqslant i < 4,\quad 0 \leqslant j < 2 \quad &\text{(starting at } (3,0)) \\
        \end{aligned}
        \]
    \item \texttt{matrix\_copy\_submatrix\_in} is applied three times with different starting indices to insert 
        $[U]_{22}$ as submatrices into $3 \times 3$ matrices:
        \[
        [T]_{33} = 
        \begin{bmatrix}
            100 & 100 & 100 \\
            100 & ([U]_{22})_{11} & ([U]_{22})_{12} \\
            100 & ([U]_{22})_{21} & ([U]_{22})_{22}
        \end{bmatrix}
        \]
        \[
        [X]_{33} = 
        \begin{bmatrix}
            100 & 100 & 100 \\
            100 & 100 & ([U]_{22})_{11} \\
            100 & 100 & ([U]_{22})_{21}
        \end{bmatrix}
        \]
        \[
        [Y]_{33} = 
        \begin{bmatrix}
            100 & 100 & 100 \\
            100 & 100 & 100 \\
            100 & ([U]_{22})_{11} & ([U]_{22})_{12}
        \end{bmatrix}
        \]
\end{enumerate}

\section{Calculus Methods}
\subsection{Backward Difference}
This method is verified in
\textit{SIM\_unit\_math\_utils/RUN\_backward\_difference}.
For this test, the function $y = 4x^3 - 3x^2 +2x +1$ is evaluated at $x
\in \{5,4,3,2,1\}$ to yield corresponding values \newline
$y_{true} \in \{436, 217, 88, 25, 4\}$.

For each test, the \textit{history} list is populated with some subset of
$y_{true}$, and the backward-difference method called using this history. In
each case, the value \textit{dy} can be compared to the value of the derivative
function $y' = 12x^2-6x +2$ with $y'(5) = 272$

\begin{itemize}
\item \textbf{History is empty} This case produces an error message and
assignment of $dy=0$.
\item \textbf{History = \{436\}} This case produces a trivial result: with
one history point, the resulting function is approximated as constant.
$dy = 0$.
This value matches the simulation output.
\item \textbf{History = \{436,217\}} With two data points, the resulting
function is approximated as linear. $dy = 436-217 = 219$.
This value matches the simulation output.
\item \textbf{History = \{436,217,88\}} With three data points, the
resulting function is approximated as quadratic
\begin{align*}
f(x) & =ax^2 + bx +c \\
f'(x) & = 2ax+b
\end{align*}
Differencing adjacent values yields
\begin{equation*}
f(x-1) = f(x) - 2ax + a -b \quad \Rightarrow \quad f(x) - f(x-1) = 2ax - a + b
\end{equation*}
With the three values specified, we have:
\begin{align*}
f(x) &= 435 \\
f(x-1) &= 217 \\
f(x-2) &= 88
\end{align*}
The first difference, $f(x) - f(x-1)$ generates $219 = 2ax - a + b$.\newline
The second difference, $f(x-1) - f(x-2)$ generates $129 = 2ax - 3a + b$.\newline
Subtracting the second difference from the first, yields $a=45$. Subsequently,
for $x=5$, $b=-186$. \newline
Thus $dy = 2ax+b = 264$. \newline
This value matches the simulation output.
\item \textbf{History = \{436,217,88,25\}}. With four data points, the
representative function is treated as a cubic
\begin{align*}
f(x) &= ax^3 + bx^2 + cx +d \\
f'(x) &= 3ax^2 + 2bx + c
\end{align*}
We can follow the same process applied to the previous test:
\begin{itemize}
\item the difference between adjacent points is
$f(x-1) = f(x) - 3ax^2+3ax-a-2bx+b-c$.
\item the first difference yields
$(3ax^2 + 2bx + c) = 219 +(3ax-a+b)$.
\item the second difference yields $(3ax-a+b) = 45+2a$.
\item the third difference yields $a=4$.
\end{itemize}
Thus
$f'(x) = 3ax^2 + 2bx + c = 219 + (45 + 2(4)) = 272$.\newline
The model generates a value of 271.9999999999999.
\item \textbf{History = \{436,217,88,25,4\}} With five data points, the
representative function is treated as a quartic and should match
perfectly with the underlying quartic function.
We should expect the numerical result to match with the analytic
result, $dy = 272$.
The model generates the expected value of $272$.
\item \textbf{History = \{436,217,88,25,4,999999\}} For this test, we
populate the 6th member of the history list with an unrelated value.
If this number is used in the algorithm, it will bias the end
result. The algorithm should stop at the 5th value, so if the result
is still $dy = 272$, we know that this 6th value was successfully
ignored.
The model generates the expected value of $272$.
\end{itemize}



\subsection{Unit Vector Derivative}
This method is verified in
\textit{SIM\_unit\_math\_utils/RUN\_unit\_vec\_deriv}

The verification tests a sequence of vector and vector-derivative
specifications. In each of the following tests, the vector is presented
followed by the vector-derivative and the resulting unit-vector-derivative.
\begin{enumerate}
\item $[0,0,0]~~[0,0,0] \rightarrow [0,0,0]$
Because the vector (the first set of values) is
zero, an error message is posted and the unit-vector-derivative is
populated from the normalization of the vector-derivative (the
second set of values) which, in this
case, creates the zero vector for the unit-vector-derivative (the
third set of values).
\item $[0,0,0]~~[3,4,0] \rightarrow [0.6,0.8,0]$
As with the previous case, the vector is zero so the
unit-vector-derivative is
populated from the normalization of the vector-derivative.
\item $[0,0,1]~~[0,0,1] \rightarrow [0,0,0]$
The two vectors are aligned so the derivative only serves to extend
the original vector, not change its direction. The unit-vector
remains unchanged and the unit-vector-derivative is zero.
\item $[0,0,1]~~[0,1,0] \rightarrow [0,1,0]$
The vector and vector-derivative are perpendicular;
the unit-vector-derivative is the vector-derivative scaled by
the reciprocal of the magnitude of the vector.
\item $[0,0,1]~~[0,5,0] \rightarrow [0,5,0]$
For this test, we simply scale up the vector-derivative;
the unit-vector-derivative scales up by the same factor.
\item $[0,0,5]~~[0,5,0] \rightarrow [0,1,0]$
For this test, we simply scale up the vector by the same factor used
in the previous test (\#5);
the unit-vector-derivative scales down by the same factor, returning
to the value seen in test \#4.
\item $[1,1,0]~~[0,0,2] \rightarrow [0,0,\sqrt2]$
As with test \#4, the vector and vector-derivative are
perpendicular, although now not unit vectors.
\item $[1,1,0]~~[0,2,2] \rightarrow \frac{1}{\sqrt2}~[-1,1,2]$
Adding a y-component to the vector derivative unsurprisingly adds a
y-component to the unit-vector-derivative; the
z-component of the unit-vector-derivative is unaffected because it
remains perpendicular to the position vector. Perhaps suprisingly,
this also introduces an x-component to the unit-vector-derivative
even though there is no x-component in the vector-derivative.
Because the unit-vector includes x- and y-components, and the
vector-derivative includes a y-component, the resulting rotation of
the vector will immediately affect the x-component of the
unit-vector.
\item $[1,2,3]~~[3,2,1] \rightarrow \frac{2\sqrt2}{7\sqrt7}[4,1,-2]$
($\approx [0.611, 0.153, -0.305]$).
This test is used as a basis for the next two tests; there is no
trivially intuitive interpretation of these values.
\item $[1,2,3]~~[6,4,2] \rightarrow \frac{4\sqrt2}{7\sqrt7}[4,1,-2]$
($\approx [1.222, 0.305, -0.611]$).
Like test \#5, we scale up the vector-derivative from the previous
test and notice that the unit-vector-derivative scales up
accordingly even when the vector and vector-derivative are not
perpendicular.
\item $[4,8,12]~~[6,4,2] \rightarrow \frac{\sqrt2}{7\sqrt7}[4,1,-2]$
($\approx [0.305, 0.076, -0.153]$).
Like test \#6, we scale up the vector from the previous
test and notice that the unit-vector-derivative scales down
accordingly even when the vector and vector-derivative are not
perpendicular.
\end{enumerate}
Outputs from this run all match with expectations.

\section{Geometry Methods}
\subsection{Ellipsoid Intersection}
The ellipsoid intersection capability is verified in 
\textit{SIM\_ellipsoid/RUN\_test}, where lines defined from source 
points to the target points, as listed in the following table, are evaluated for 
intersection with the ellipsoid:
\begin{equation*}
\left( \frac{x}{2} \right)^2 +
\left( \frac{y}{3} \right)^2 +
\left( \frac{z}{4} \right)^2 = 1
\end{equation*}

\begin{center}
\begin{tabular}{|l|l|l||l|l|}
\hline
 \textbf{ } & \textbf{Source Point} & \textbf{Target Point} & \textbf{Intersection} &
 \textbf{Expected Intersection Points} \\
\hline \hline
Line 1 & $(5,0,0)$  & $(-5,0,0)$ & true & \{$(2,0,0),(-2,0,0)$\} \\
Line 2 & $(0,5,0)$  & $(0,-5,0)$ & true & \{$(0,3,0),(0,-3,0)$\} \\
Line 3 & $(0,0,5)$  & $(0,0,-5)$ & true & \{$(0,0,4),(0,0,-4)$\} \\
Line 4 & $(-1,0,4)$ & $(1,0,4)$  & true & $(0,0,4)$ \\
Line 5 & $(1,1,1)$  & $(0,0,1)$  & true & \{$(1.611, 1.611, 1),(-1.611,-1.611,1)$\} \\
Line 6 & $(1.5,2.5,0)$  & $(1.5,2.5,1)$ & false & -- \\
\hline \hline
\end{tabular}
\end{center}

The expected results shown in the table can be visually confirmed using Figure \ref{fig:ellip_verif_unique}, which illustrates 
the ellipsoid and the corresponding lines, allowing verification of whether and where each 
intersection occurs. These results are consistent with that of the run.

\begin{figure}[h!]
\begin{center}
    \includegraphics[width=\linewidth]{Figs/ellipsoid_multiview.png}
    \caption{Multiview Plots of Ellipsoid and Intersecting Lines}
    \label{fig:ellip_verif_unique}
\end{center}
\end{figure}

\end{document}
